\chapter{Datatypes}
\mpitermtitleindex{datatypes}
\label{chap:datatypes}

% Define a macro to format words in math mode. Otherwise, each letter is
% treated as a separate variable name (as is the convention in math)
%ALLOWLATEX:mword%
\def\mword#1{\textrm{\itshape{#1}}} %ALLOWLATEX%
%ALLOWLATEX:msword%
\def\msword#1{\textrm{\itshape{\scriptsize #1}}} %ALLOWLATEX%
%ALLOWLATEX:mscode%
\def\mscode#1{\mpicode{\scriptsize #1}} %%ALLOWLATEX%
Basic datatypes are introduced
in \sectionref{subsec:pt2pt-messagedata}
and
in \sectionref{sec:pt2pt-typematch}.
In this chapter, this model is
extended
to describe any data layout.
We consider general
datatypes that allow one to transfer efficiently
heterogeneous and noncontiguous data.
We conclude with the description of calls for explicit packing and unpacking of
messages.

\section{Derived Datatypes}
\mpitermtitleindex{derived datatype}
\label{sec:pt2pt-datatype}

Point-to-point communications on buffers containing a sequence of identical basic datatypes
is constraining.  One often wants to pass
messages that contain values with different datatypes (e.g., an integer count,
followed by a sequence of real numbers); and one often wants to send
noncontiguous data (e.g., a sub-block of a matrix).  One solution is to
pack noncontiguous data into a contiguous buffer
at the sender site and unpack it
at the receiver site.   This has
the disadvantage of requiring additional memory-to-memory copy operations
at both sites, even when the communication
subsystem has scatter-gather capabilities.   Instead, \MPI/ provides
mechanisms to specify more general, mixed, and noncontiguous
communication buffers. It is up to the implementation to decide
whether data should be first packed in a contiguous buffer before being
transmitted, or whether it can be collected directly from where it
resides.

The general mechanisms provided here allow one to transfer directly,
without copying, objects of various shapes and sizes.  It is not assumed
that the \MPI/ library is cognizant of the objects declared in the host
language. Thus, if one wants to transfer a structure, or an array
section, it will be necessary to provide in \MPI/ a definition of a
communication buffer that mimics the definition of the structure or
array section in question.  These facilities can be used by library
designers to define communication functions that can transfer objects
defined in the host language---by decoding their definitions as
available in a symbol table or a dope vector.  Such higher-level
communication functions are not part of \MPI/.

More general communication buffers are specified by replacing the
basic datatypes that have been used so far with derived datatypes that
are constructed from basic datatypes using the constructors described
in this section.  These methods of constructing derived datatypes can
be applied recursively.

A \mpitermdef{general datatype} is an opaque object that specifies two
things:
\begin{itemize}
\item
A sequence of basic datatypes.
\item
A sequence of integer (byte) displacements.
\end{itemize}

The displacements are not required to be positive, distinct, or
in increasing order. Therefore, the order of items need not
coincide with their order in store, and an item may appear more than
once.
We call such a pair of sequences (or sequence of pairs) a \mpitermdef{type map}.
The sequence of basic datatypes (displacements ignored) is the \mpitermdef{type
signature} of the datatype.

Let
\[
\mword{Typemap} = \{ (\mword{type}_0,\mword{disp}_0), \ldots , (\mword{type}_{n-1}, \mword{disp}_{n-1}) \} ,
\]
be such a type map, where $\mword{type}_i$ are basic types, and
$\mword{disp}_i$ are  displacements.
Let
\[
\mword{Typesig} = \{ \mword{type}_0 , \ldots , \mword{type}_{n-1} \}
\]
be the associated type signature.
This type map, together with a base address \mpicode{buf},
specifies a communication buffer: the communication buffer that consists of $n$
entries, where the $i$-th entry is at address $\mpicode{buf} +
\mword{disp}_i$ and has type $\mword{type}_i$.
A message assembled from such a
communication buffer will consist of $n$ values, of the types defined
by $\mword{Typesig}$.

Most datatype constructors have replication count or block length arguments.
Allowed values are nonnegative integers. If the value is zero, no elements are
generated in the type map and there is no effect on datatype bounds or
extent.

We can use a handle to a general datatype as an argument in a send or
receive operation, instead of a basic datatype argument.  The
operation
\mpifunc{MPI\_SEND}\mpicode{(buf, 1, datatype,$\ldots$)} will use the send buffer
defined by the base address \mpiarg{buf} and the general datatype
associated with \mpiarg{datatype}; it will generate a message with the type
signature determined by the \mpiarg{datatype} argument.
\mpifunc{MPI\_RECV}\mpicode{(buf, 1, datatype,$\ldots$)} will use the receive buffer
defined by the base address \mpiarg{buf} and the general datatype
associated with \mpiarg{datatype}.

General datatypes can be used in all send and receive
operations.  We discuss, in \sectionref{subsec:pt2pt-datatypeuse}, the
case where the second argument \mpiarg{count} has value $> 1$.

The basic datatypes presented in
Section~\ref{subsec:pt2pt-messagedata}
are particular cases of a general datatype, and are predefined.
Thus, \type{MPI\_INT} is a predefined handle to a datatype with type
map $\{ (\ctype{int}, 0) \}$, with one entry of type \ctype{int} and
displacement zero.  The other basic datatypes are similar.

The \mpitermdefni{extent}\mpitermdefindex{extent of datatypes} of a datatype is defined to
be the span from the first byte to the last byte occupied by entries in this
datatype, rounded up to satisfy alignment requirements.
That is, if
\[
\mword{Typemap} = \{ (\mword{type}_0,\mword{disp}_0), \ldots , (\mword{type}_{n-1}, \mword{disp}_{n-1}) \} ,
\]
then
\begin{eqnarray}
lb(\mword{Typemap}) & = & \min_j \mword{disp}_j , \nonumber \\
ub(\mword{Typemap}) & = & \max_j (\mword{disp}_j + \mpicode{sizeof}(\mword{type}_j)) + \epsilon , \mbox{ and}
\nonumber \\ extent(\mword{Typemap}) & = & ub(\mword{Typemap}) - lb(\mword{Typemap}).
\label{soft-lb-ub-definition}
\end{eqnarray}
If $\mword{type}_j$ requires alignment to a byte address that
is
a multiple
of $k_j$,
then $\epsilon$ is the least nonnegative increment needed to round
$extent(\mword{Typemap})$ to the next multiple of $\max_j k_j$.
In Fortran, it is implementation dependent whether the \MPI/ implementation
computes the alignments $k_j$ according to the alignments
used by the compiler in common blocks, \ftype{SEQUENCE} derived types,
\ftype{BIND(C)} derived types,
or derived types that are neither \ftype{SEQUENCE} nor \ftype{BIND(C)}.
The complete definition of \mpitermdefni{extent}\mpitermdefindex{extent of datatypes} is given
by Equation~\ref{soft-lb-ub-definition}.


\begin{example}
\label{pt2pt-exR}
\exindex{Typemap}%
\Exindex{Neutral}{Typemap}{}%
Assume that
$Type = \{ (\ctype{double},0), (\ctype{char}, 8) \}$
(a \ctype{double} at
displacement zero, followed by a \ctype{char} at displacement eight).
Assume, furthermore, that
doubles have to be strictly aligned at addresses that are multiples of eight.
Then, the extent of this datatype is 16 (9 rounded to the next multiple of 8).
A datatype that consists of a character immediately followed by a double will
also have an extent of 16.
\end{example}

\begin{rationale}
The definition of extent is motivated by the assumption that
the amount of padding added at the end of each structure in an array of
structures is the least needed to fulfill alignment constraints.
More explicit control of the extent is provided in
Section~\ref{subsec:pt2pt-markers}.  Such explicit control is needed
in cases where the assumption does not hold, for example, where union types
are used.
In Fortran, structures can be expressed with several language features, e.g.,
common blocks, \ftype{SEQUENCE} derived types, or \ftype{BIND(C)} derived types. The compiler may
use different alignments, and therefore, it is recommended to use \mpifunc{MPI\_TYPE\_CREATE\_RESIZED}
for arrays of structures if an alignment may cause an alignment-gap at the end of
a structure as described
in \sectionref{subsec:pt2pt-markers} and
in \sectionref{sec:f90-problems:derived-types}.
\end{rationale}




\subsection{Type Constructors with Explicit Addresses}
\label{sec:misc-extent}
In Fortran, the procedures
\mpifunc{MPI\_TYPE\_CREATE\_HVECTOR},\flushline  % force break to keep in margin
\mpifunc{MPI\_TYPE\_CREATE\_HINDEXED},
\mpifunc{MPI\_TYPE\_CREATE\_HINDEXED\_BLOCK},
\mpifunc{MPI\_TYPE\_CREATE\_STRUCT}, and \mpifunc{MPI\_GET\_ADDRESS}
accept arguments of type
\mpiftypekind{ADDRESS}, wherever arguments of type
\cdeclmainindex{MPI\_Aint}%
\mpictype{MPI\_Aint}
are used in C.
For Fortran compilers that do not support the Fortran~90 \ftype{KIND}
notation, and where addresses are 64~bits whereas default \ftype{INTEGER}s
are 32~bits, these arguments will be of type \ftype{INTEGER*8} (assuming
the Fortran compiler accepts the common extension of \ftype{INTEGER*8}
for eight-byte integers).

For the large count versions of three datatype constructors with explicit addresses,
\mpifunc{MPI\_TYPE\_CREATE\_HINDEXED}, \mpifunc{MPI\_TYPE\_CREATE\_HINDEXED\_BLOCK}, and
\mpifunc{MPI\_TYPE\_CREATE\_STRUCT}, absolute addresses shall not be used to
specify byte displacements since the parameter is of type \mpidtype{MPI\_COUNT}
instead of type \mpidtype{MPI\_AINT}.

\subsection{Datatype Constructors}
\label{subsec:pt2pt-datatypeconst}

\paraheading{Contiguous.}
The simplest datatype constructor is
\mpifunc{MPI\_TYPE\_CONTIGUOUS}, which
allows replication of a datatype into contiguous locations.

\cdeclmainindex{MPI\_Datatype}\cdeclmainindex{TYPE(MPI\_Datatype)}%
\begin{funcdef}{MPI\_TYPE\_CONTIGUOUS}{(\mbox{count}, \mbox{oldtype}, \mbox{newtype})}
\funcarg{\IN}{count}{replication count (nonnegative integer)}
\funcarg{\IN}{oldtype}{old datatype (handle)}
\funcarg{\OUT}{newtype}{new datatype (handle)}
\end{funcdef}
\mpibindmain{MPI\_Type\_contiguous}{(\gb{}int~count, MPI\_Datatype~oldtype, MPI\_Datatype~*newtype)}
\MPIbindindex{MPI\_Type\_contiguous\_c}
\mpibindmain{MPI\_Type\_contiguous\_c}{(\gb{}MPI\_Count~count, MPI\_Datatype~oldtype, MPI\_Datatype~*newtype)}
\mpifnewbindmain{MPI\_Type\_contiguous}{(count, oldtype, newtype, ierror) \fargs INTEGER, INTENT(IN)\ ::\ \mbox{count}\fnarg{}TYPE(MPI\_Datatype), INTENT(IN)\ ::\ \mbox{oldtype}\fnarg{}TYPE(MPI\_Datatype), INTENT(OUT)\ ::\ \mbox{newtype}\fnfinalarg{}INTEGER, OPTIONAL, INTENT(OUT)\ ::\ \mbox{ierror}}
\mpifnewbindmain{MPI\_Type\_contiguous}{(count, oldtype, newtype, ierror) !(\_c) \fargs INTEGER(KIND=MPI\_COUNT\_KIND), INTENT(IN)\ ::\ \mbox{count}\fnarg{}TYPE(MPI\_Datatype), INTENT(IN)\ ::\ \mbox{oldtype}\fnarg{}TYPE(MPI\_Datatype), INTENT(OUT)\ ::\ \mbox{newtype}\fnfinalarg{}INTEGER, OPTIONAL, INTENT(OUT)\ ::\ \mbox{ierror}}
\mpifbindmain{MPI\_TYPE\_CONTIGUOUS}{(COUNT, OLDTYPE, NEWTYPE, IERROR) \fargs \fnfinalarg{}INTEGER \mbox{COUNT}, \mbox{OLDTYPE}, \mbox{NEWTYPE}, \mbox{IERROR}}


\noindent
\mpiarg{newtype} is the datatype obtained by concatenating
\mpiarg{count} copies of
\mpiargshort{oldtype}.
Concatenation is defined using \emph{extent} as the size of
the concatenated copies.


\begin{example}
\label{pt2pt-exS}
\exindex{Typemap}%
\Exindex{Neutral}{Typemap for contiguous}{MPI\_TYPE\_CONTIGUOUS}%
Let \mpiarg{oldtype} have type map
\(
\{ (\ctype{double}, 0), (\ctype{char}, 8) \} ,
\)
with extent 16,
and let $\mpiarg{count} = 3$.  The type map of
the datatype returned by \mpiarg{newtype} is
\begin{displaymath}
\{ (\ctype{double}, 0), (\ctype{char}, 8), (\ctype{double}, 16),
   (\ctype{char}, 24), (\ctype{double}, 32), (\ctype{char}, 40) \};
\end{displaymath}
i.e., alternating \ctype{double} and \ctype{char} elements, with displacements
$0, 8, 16, 24, 32, 40$.
\end{example}

In general,
assume that the type map of \mpiarg{oldtype} is
\begin{displaymath}
\{ (\mword{type}_0,\mword{disp}_0), \ldots , (\mword{type}_{n-1}, \mword{disp}_{n-1}) \} ,
\end{displaymath}
with extent $ex$.
Then \mpiarg{newtype} has a type map with $\mpicode{count} \cdot \mpicode{n}$
entries defined by:
\begin{displaymath}
% The \hskip is to keep the equation within the text margins
\hskip-1em\{ (\mword{type}_0, \mword{disp}_0), \ldots , (\mword{type}_{n-1}, \mword{disp}_{n-1}), %ALLOWLATEX%
(\mword{type}_0, \mword{disp}_0 +ex), \ldots ,(\mword{type}_{n-1}, \mword{disp}_{n-1} + ex),\\
\end{displaymath}
% This is *not* the way to display a multiline equation - FIXME
\begin{displaymath}
\hskip-1em\ldots,(\mword{type}_0, \mword{disp}_0 +ex \cdot(\mpicode{count}-1) ), \ldots ,%ALLOWLATEX%
(\mword{type}_{n-1} , \mword{disp}_{n-1} + ex \cdot (\mpicode{count}-1)) \} .
\end{displaymath}

\paraheading{Vector.}
The procedure
\mpifunc{MPI\_TYPE\_VECTOR} is a more general constructor that
allows replication of a datatype
into locations that consist of equally spaced blocks.  Each block
is obtained by concatenating the same number of copies of the old datatype.
The spacing between blocks is a multiple of the extent of the old datatype.

\begin{funcdef}{MPI\_TYPE\_VECTOR}{(\mbox{count}, \mbox{blocklength}, \mbox{stride}, \mbox{oldtype}, \mbox{newtype})}
\funcarg{\IN}{count}{number of blocks (nonnegative integer)}
\funcarg{\IN}{blocklength}{number of elements in each block (nonnegative integer)}
\funcarg{\IN}{stride}{number of elements between start of each block (integer)}
\funcarg{\IN}{oldtype}{old datatype (handle)}
\funcarg{\OUT}{newtype}{new datatype (handle)}
\end{funcdef}
\mpibindmain{MPI\_Type\_vector}{(\gb{}int~count, int~blocklength, int~stride, MPI\_Datatype~oldtype, MPI\_Datatype~*newtype)}
\MPIbindindex{MPI\_Type\_vector\_c}
\mpibindmain{MPI\_Type\_vector\_c}{(\gb{}MPI\_Count~count, MPI\_Count~blocklength, MPI\_Count~stride, MPI\_Datatype~oldtype, MPI\_Datatype~*newtype)}
\mpifnewbindmain{MPI\_Type\_vector}{(count, blocklength, stride, oldtype, newtype, ierror) \fargs INTEGER, INTENT(IN)\ ::\ \mbox{count}, \mbox{blocklength}, \mbox{stride}\fnarg{}TYPE(MPI\_Datatype), INTENT(IN)\ ::\ \mbox{oldtype}\fnarg{}TYPE(MPI\_Datatype), INTENT(OUT)\ ::\ \mbox{newtype}\fnfinalarg{}INTEGER, OPTIONAL, INTENT(OUT)\ ::\ \mbox{ierror}}
\mpifnewbindmain{MPI\_Type\_vector}{(count, blocklength, stride, oldtype, newtype, ierror) !(\_c) \fargs INTEGER(KIND=MPI\_COUNT\_KIND), INTENT(IN)\ ::\ \mbox{count}, \mbox{blocklength}, \mbox{stride}\fnarg{}TYPE(MPI\_Datatype), INTENT(IN)\ ::\ \mbox{oldtype}\fnarg{}TYPE(MPI\_Datatype), INTENT(OUT)\ ::\ \mbox{newtype}\fnfinalarg{}INTEGER, OPTIONAL, INTENT(OUT)\ ::\ \mbox{ierror}}
\mpifbindmain{MPI\_TYPE\_VECTOR}{(COUNT, BLOCKLENGTH, STRIDE, OLDTYPE, NEWTYPE, IERROR) \fargs \fnfinalarg{}INTEGER \mbox{COUNT}, \mbox{BLOCKLENGTH}, \mbox{STRIDE}, \mbox{OLDTYPE}, \mbox{NEWTYPE}, \mbox{IERROR}}

\begin{example}
\label{pt2pt-exT}
\exindex{Typemap}%
\Exindex{Neutral}{Typemap for vector}{MPI\_TYPE\_VECTOR}%
Assume, again, that \mpiarg{oldtype} has type map
\(
\{ (\ctype{double}, 0), (\ctype{char}, 8) \} ,
\)
with extent 16.
A call to \mpifunc{MPI\_TYPE\_VECTOR}\mpicode{(2, 3, 4, oldtype, newtype)} will
create the datatype with type map,
\begin{displaymath}
\{
(\ctype{double}, 0), (\ctype{char}, 8), (\ctype{double}, 16), (\ctype{char},
24), (\ctype{double}, 32), (\ctype{char}, 40),
\end{displaymath}
\begin{displaymath}
(\ctype{double}, 64), (\ctype{char}, 72), (\ctype{double}, 80), (\ctype{char},
88), (\ctype{double}, 96), (\ctype{char}, 104)
\} .
\end{displaymath}
That is, two blocks with three copies each of the old
type, with a stride of 4 elements ($4 \cdot 16$ bytes) between the start of each block.
\end{example}

\begin{example}
\label{pt2pt-exTTT}
\exindex{Typemap}%
\Exindex{Neutral}{Typemap for vector}{MPI\_TYPE\_VECTOR}%
A call to \mpifunc{MPI\_TYPE\_VECTOR}\mpicode{(3, 1, -2, oldtype, newtype)} will create
the datatype,
\begin{displaymath}
\{
(\ctype{double}, 0), (\ctype{char}, 8), (\ctype{double}, -32), (\ctype{char},
-24), (\ctype{double}, -64), (\ctype{char}, -56)
\} .
\end{displaymath}
\end{example}

In general, assume that \mpiarg{oldtype} has type map,
\begin{displaymath}
\{ (\mword{type}_0,\mword{disp}_0), \ldots, (\mword{type}_{n-1}, \mword{disp}_{n-1}) \} ,
\end{displaymath}
with extent $ex$.  Let \mpicode{bl} be the \mpicode{blocklength}.
The newly created  datatype has a type map with
$\mpicode{count} \cdot \mpicode{bl} \cdot \mpicode{n}$
entries:
\begin{displaymath}
\{
(\mword{type}_0, \mword{disp}_0), \ldots , (\mword{type}_{n-1} , \mword{disp}_{n-1}),
\end{displaymath}
\begin{displaymath}
(\mword{type}_0 ,\mword{disp}_0 + ex) , \ldots ,
(\mword{type}_{n-1} , \mword{disp}_{n-1} + ex ), \ldots,
\end{displaymath}
\begin{displaymath}
(\mword{type}_0 , \mword{disp}_0 + (\mpicode{bl} -1) \cdot ex
) , \ldots , (\mword{type}_{n-1} , \mword{disp}_{n-1} + (\mpicode{bl} -1) \cdot ex ) ,
\end{displaymath}
\begin{displaymath}
(\mword{type}_0 ,\mword{disp}_0 + \mpicode{stride} \cdot ex ) , \ldots ,
(\mword{type}_{n-1} , \mword{disp}_{n-1} + \mpicode{stride} \cdot ex ), \ldots ,
\end{displaymath}
\begin{displaymath}
(\mword{type}_0 , \mword{disp}_0 + (\mpicode{stride} + \mpicode{bl} -1) \cdot ex ) , \ldots ,
(\mword{type}_{n-1}, \mword{disp}_{n-1} + (\mpicode{stride} + \mpicode{bl} -1) \cdot
ex ) , \ldots ,
\end{displaymath}
\begin{displaymath}
(\mword{type}_0 ,\mword{disp}_0 + \mpicode{stride} \cdot (\mpicode{count}-1) \cdot ex ) , \ldots ,
\end{displaymath}
\begin{displaymath}
(\mword{type}_{n-1} , \mword{disp}_{n-1} + \mpicode{stride} \cdot (\mpicode{count} -1) \cdot
ex )
, \ldots ,
\end{displaymath}
\begin{displaymath}
(\mword{type}_0 , \mword{disp}_0 + (\mpicode{stride} \cdot (\mpicode{count} -1)
+ \mpicode{bl} -1) \cdot ex ) , \ldots ,
\end{displaymath}
\begin{displaymath}
(\mword{type}_{n-1}, \mword{disp}_{n-1} + (\mpicode{stride} \cdot (\mpicode{count} -1)
+ \mpicode{bl} -1) \cdot ex )
\} .
\end{displaymath}

A call to \mpifunc{MPI\_TYPE\_CONTIGUOUS}\mpicode{(count, oldtype, newtype)} is
equivalent to a call to
\mpifunc{MPI\_TYPE\_VECTOR}\mpicode{(count, 1, 1, oldtype, newtype)}, or to a call to
\mpifunc{MPI\_TYPE\_VECTOR}\mpicode{(1, count, n, oldtype, newtype)}, where \mpicode{n}
is an arbitrary integer value.

\paraheading{Hvector.}
The procedure
\mpifunc{MPI\_TYPE\_CREATE\_HVECTOR}
is identical to
\mpifunc{MPI\_TYPE\_VECTOR}, except that \mpiarg{stride} is given in bytes,
rather than in elements.  The use for both types of vector
constructors is illustrated in Section~\ref{subsec:pt2pt-examples}.
(\mpicode{H} stands for ``heterogeneous'').


\cdeclindex{MPI\_Aint}%
\begin{funcdef}{MPI\_TYPE\_CREATE\_HVECTOR}{(\mbox{count}, \mbox{blocklength}, \mbox{stride}, \mbox{oldtype}, \mbox{newtype})}
\funcarg{\IN}{count}{number of blocks (nonnegative integer)}
\funcarg{\IN}{blocklength}{number of elements in each block (nonnegative integer)}
\funcarg{\IN}{stride}{number of bytes between start of each block (integer)}
\funcarg{\IN}{oldtype}{old datatype (handle)}
\funcarg{\OUT}{newtype}{new datatype (handle)}
\end{funcdef}
\mpibindmain{MPI\_Type\_create\_hvector}{(\gb{}int~count, int~blocklength, MPI\_Aint~stride, MPI\_Datatype~oldtype, MPI\_Datatype~*newtype)}
\MPIbindindex{MPI\_Type\_create\_hvector\_c}
\mpibindmain{MPI\_Type\_create\_hvector\_c}{(\gb{}MPI\_Count~count, MPI\_Count~blocklength, MPI\_Count~stride, MPI\_Datatype~oldtype, MPI\_Datatype~*newtype)}
\mpifnewbindmain{MPI\_Type\_create\_hvector}{(count, blocklength, stride, oldtype, newtype, ierror) \fargs INTEGER, INTENT(IN)\ ::\ \mbox{count}, \mbox{blocklength}\fnarg{}INTEGER(KIND=MPI\_ADDRESS\_KIND), INTENT(IN)\ ::\ \mbox{stride}\fnarg{}TYPE(MPI\_Datatype), INTENT(IN)\ ::\ \mbox{oldtype}\fnarg{}TYPE(MPI\_Datatype), INTENT(OUT)\ ::\ \mbox{newtype}\fnfinalarg{}INTEGER, OPTIONAL, INTENT(OUT)\ ::\ \mbox{ierror}}
\mpifnewbindmain{MPI\_Type\_create\_hvector}{(count, blocklength, stride, oldtype, newtype, ierror) !(\_c) \fargs INTEGER(KIND=MPI\_COUNT\_KIND), INTENT(IN)\ ::\ \mbox{count}, \mbox{blocklength}, \mbox{stride}\fnarg{}TYPE(MPI\_Datatype), INTENT(IN)\ ::\ \mbox{oldtype}\fnarg{}TYPE(MPI\_Datatype), INTENT(OUT)\ ::\ \mbox{newtype}\fnfinalarg{}INTEGER, OPTIONAL, INTENT(OUT)\ ::\ \mbox{ierror}}
\mpifbindmain{MPI\_TYPE\_CREATE\_HVECTOR}{(COUNT, BLOCKLENGTH, STRIDE, OLDTYPE, NEWTYPE, IERROR) \fargs INTEGER \mbox{COUNT}, \mbox{BLOCKLENGTH}, \mbox{OLDTYPE}, \mbox{NEWTYPE}, \mbox{IERROR}\fnfinalarg{}INTEGER(KIND=MPI\_ADDRESS\_KIND) \mbox{STRIDE}}

Assume that \mpiarg{oldtype} has type map,
\begin{displaymath}
\{ (\mword{type}_0,\mword{disp}_0), \ldots, (\mword{type}_{n-1}, \mword{disp}_{n-1}) \} ,
\end{displaymath}
with extent $ex$.  Let \mpicode{bl} be the \mpicode{blocklength}.
The newly created  datatype has a type map with
$\mpicode{count} \cdot \mpicode{bl} \cdot n$
entries:
\begin{displaymath}
\{
(\mword{type}_0, \mword{disp}_0), \ldots , (\mword{type}_{n-1} , \mword{disp}_{n-1}),
\end{displaymath}
\begin{displaymath}
(\mword{type}_0 ,\mword{disp}_0 + ex) , \ldots ,
(\mword{type}_{n-1} , \mword{disp}_{n-1} + ex ), \ldots,
\end{displaymath}
\begin{displaymath}
(\mword{type}_0 , \mword{disp}_0 + (\mpicode{bl} -1) \cdot ex
) , \ldots , (\mword{type}_{n-1} , \mword{disp}_{n-1} + (\mpicode{bl} -1) \cdot ex ) ,
\end{displaymath}
\begin{displaymath}
(\mword{type}_0 ,\mword{disp}_0 + \mpicode{stride}  ) , \ldots ,
(\mword{type}_{n-1} , \mword{disp}_{n-1} + \mpicode{stride} ) , \ldots ,
\end{displaymath}
\begin{displaymath}
(\mword{type}_0 , \mword{disp}_0 + \mpicode{stride} + ( \mpicode{bl} -1) \cdot ex
) , \ldots ,
\end{displaymath}
\begin{displaymath}
(\mword{type}_{n-1}, \mword{disp}_{n-1} + \mpicode{stride} + (\mpicode{bl} -1) \cdot
ex ) ,
\ldots ,
\end{displaymath}
\begin{displaymath}
(\mword{type}_0 ,\mword{disp}_0 + \mpicode{stride} \cdot (\mpicode{count}-1) ) , \ldots ,
(\mword{type}_{n-1} , \mword{disp}_{n-1} + \mpicode{stride} \cdot (\mpicode{count} -1)  )
, \ldots ,
\end{displaymath}
\begin{displaymath}
(\mword{type}_0 , \mword{disp}_0 + \mpicode{stride} \cdot (\mpicode{count} -1)
+ (\mpicode{bl} -1) \cdot ex ) , \ldots ,
\end{displaymath}
\begin{displaymath}
(\mword{type}_{n-1}, \mword{disp}_{n-1} + \mpicode{stride} \cdot (\mpicode{count} -1)
+ (\mpicode{bl} -1) \cdot ex )
\} .
\end{displaymath}

\paraheading{Indexed.}
The procedure
\mpifunc{MPI\_TYPE\_INDEXED} allows
replication of an old datatype into a sequence of blocks (each block is
a concatenation of the old datatype), where
each block can contain a different number of copies and have a different
displacement.  All block displacements are multiples of the old type
extent.

\begin{funcdef}{MPI\_TYPE\_INDEXED}{(\mbox{count}, \mbox{array\_of\_blocklengths}, \mbox{array\_of\_displacements}, \mbox{oldtype}, \mbox{newtype})}
\funcarg{\IN}{count}{number of blocks---also number of entries in \mpiarg{array\_of\_displacements} and \mpiarg{array\_of\_blocklengths} (nonnegative integer)}
\funcarg{\IN}{array\_of\_blocklengths}{number of elements per block (array of nonnegative integers)}
\funcarg{\IN}{array\_of\_displacements}{displacement for each block, in multiples of \mpiarg{oldtype} (array of integers)}
\funcarg{\IN}{oldtype}{old datatype (handle)}
\funcarg{\OUT}{newtype}{new datatype (handle)}
\end{funcdef}
\mpibindmain{MPI\_Type\_indexed}{(\gb{}int~count, const~int~array\_of\_blocklengths[], const~int~array\_of\_displacements[], MPI\_Datatype~oldtype, MPI\_Datatype~*newtype)}
\MPIbindindex{MPI\_Type\_indexed\_c}
\mpibindmain{MPI\_Type\_indexed\_c}{(\gb{}MPI\_Count~count, const~MPI\_Count~array\_of\_blocklengths[], const~MPI\_Count~array\_of\_displacements[], MPI\_Datatype~oldtype, MPI\_Datatype~*newtype)}
\mpifnewbindmain{MPI\_Type\_indexed}{(count, array\_of\_blocklengths, array\_of\_displacements, oldtype, newtype, ierror) \fargs INTEGER, INTENT(IN)\ ::\ \mbox{count}, \mbox{array\_of\_blocklengths(count)}, \mbox{array\_of\_displacements(count)}\fnarg{}TYPE(MPI\_Datatype), INTENT(IN)\ ::\ \mbox{oldtype}\fnarg{}TYPE(MPI\_Datatype), INTENT(OUT)\ ::\ \mbox{newtype}\fnfinalarg{}INTEGER, OPTIONAL, INTENT(OUT)\ ::\ \mbox{ierror}}
\mpifnewbindmain{MPI\_Type\_indexed}{(count, array\_of\_blocklengths, array\_of\_displacements, oldtype, newtype, ierror) !(\_c) \fargs INTEGER(KIND=MPI\_COUNT\_KIND), INTENT(IN)\ ::\ \mbox{count}, \mbox{array\_of\_blocklengths(count)}, \mbox{array\_of\_displacements(count)}\fnarg{}TYPE(MPI\_Datatype), INTENT(IN)\ ::\ \mbox{oldtype}\fnarg{}TYPE(MPI\_Datatype), INTENT(OUT)\ ::\ \mbox{newtype}\fnfinalarg{}INTEGER, OPTIONAL, INTENT(OUT)\ ::\ \mbox{ierror}}
\mpifbindmain{MPI\_TYPE\_INDEXED}{(COUNT, ARRAY\_OF\_BLOCKLENGTHS, ARRAY\_OF\_DISPLACEMENTS, OLDTYPE, NEWTYPE, IERROR) \fargs \fnfinalarg{}INTEGER \mbox{COUNT}, \mbox{ARRAY\_OF\_BLOCKLENGTHS(*)}, \mbox{ARRAY\_OF\_DISPLACEMENTS(*)}, \mbox{OLDTYPE}, \mbox{NEWTYPE}, \mbox{IERROR}}


\begin{example}
\label{pt2pt-exU}
\exindex{Typemap}%
\Exindex{Neutral}{Typemap for indexed}{MPI\_TYPE\_INDEXED}%
Let \mpiarg{oldtype} have type map
\(
\{ (\ctype{double}, 0), (\ctype{char}, 8) \} ,
\)
with extent 16.
Let  \mpicode{B = (3, 1)} and let \mpicode{D = (4, 0)}.  A call to
\mpifunc{MPI\_TYPE\_INDEXED}\mpicode{(2, B, D, oldtype, newtype)} returns a datatype with
type map,
\begin{displaymath}
\{
(\ctype{double}, 64), (\ctype{char}, 72), (\ctype{double}, 80), (\ctype{char},
88), (\ctype{double}, 96), (\ctype{char}, 104),
\end{displaymath}
\begin{displaymath}
(\ctype{double}, 0), (\ctype{char}, 8)
\} .
\end{displaymath}
That is, three copies of the old type starting at displacement
64, and one copy starting at displacement 0.
\end{example}

In general,
assume that \mpiarg{oldtype} has type map,
\begin{displaymath}
\{ (\mword{type}_0,\mword{disp}_0), \ldots, (\mword{type}_{n-1}, \mword{disp}_{n-1}) \} ,
\end{displaymath}
with extent \emph{ex}.
Let \mpiarg{B} be the \mpiarg{array\_of\_blocklengths} argument and
\mpishortarg{D} be the\flushline
\mpiarg{array\_of\_displacements} argument. The newly created datatype
has $n \cdot \sum_{i=0}^{\mscode{count}-1}
\mpiarg{B[i]}$ entries:
\begin{displaymath}
\{
(\mword{type}_0, \mword{disp}_0 + \mpiarg{D[0]} \cdot ex ) , \ldots ,
(\mword{type}_{n-1} , \mword{disp}_{n-1} + \mpiarg{D[0]} \cdot ex ) , \ldots ,
\end{displaymath}
\begin{displaymath}
(\mword{type}_0 , \mword{disp}_0 + (\mpiarg{D[0]} + \mpiarg{B[0]} -1) \cdot ex) ,\ldots,
\end{displaymath}
\begin{displaymath}
(\mword{type}_{n-1} , \mword{disp}_{n-1} + (\mpiarg{D[0]} +\mpiarg{B[0]} -1) \cdot ex ) ,
\ldots,
\end{displaymath}
\begin{displaymath}
(\mword{type}_0, \mword{disp}_0 + \mpiarg{D[count-1]} \cdot ex ) , \ldots ,
(\mword{type}_{n-1} , \mword{disp}_{n-1} + \mpiarg{D[count-1]} \cdot ex ) , \ldots ,
\end{displaymath}
\begin{displaymath}
(\mword{type}_0 , \mword{disp}_0 + (\mpiarg{D[count-1]} + \mpiarg{B[count-1]} -1) \cdot ex)
,\ldots,
\end{displaymath}
\begin{displaymath}
(\mword{type}_{n-1} , \mword{disp}_{n-1} + (\mpiarg{D[count-1]} +\mpiarg{B[count-1]} -1)
\cdot ex )
\} .
\end{displaymath}

A call to \mpifunc{MPI\_TYPE\_VECTOR}\mpicode{(count, blocklength, stride, oldtype,
newtype)} is equivalent to a call to
\mpifunc{MPI\_TYPE\_INDEXED}\mpicode{(count, B, D, oldtype, newtype)} where
\begin{displaymath}
\mpiarg{D[j]} = j \cdot \mpiarg{stride},\ j=0 ,\ldots, \mpicode{count} -1 ,
\end{displaymath}
and
\begin{displaymath}
\mpiarg{B[j]} = \mpiarg{blocklength},\ j=0 ,\ldots, \mpicode{count} -1 .
\end{displaymath}

\paraheading{Hindexed.}
The procedure
\mpifunc{MPI\_TYPE\_CREATE\_HINDEXED}
is identical to
\mpifunc{MPI\_TYPE\_INDEXED}, except that block displacements in
\mpiarg{array\_of\_displacements} are specified in
bytes, rather than in multiples of the \mpiarg{oldtype} extent.

\cdeclindex{MPI\_Aint}%
\begin{funcdef}{MPI\_TYPE\_CREATE\_HINDEXED}{(\mbox{count}, \mbox{array\_of\_blocklengths}, \mbox{array\_of\_displacements}, \mbox{oldtype}, \mbox{newtype})}
\funcarg{\IN}{count}{number of blocks---also number of entries in \mpiarg{array\_of\_displacements} and \mpiarg{array\_of\_blocklengths} (nonnegative integer)}
\funcarg{\IN}{array\_of\_blocklengths}{number of elements in each block (array of nonnegative integers)}
\funcarg{\IN}{array\_of\_displacements}{byte displacement of each block (array of integers)}
\funcarg{\IN}{oldtype}{old datatype (handle)}
\funcarg{\OUT}{newtype}{new datatype (handle)}
\end{funcdef}
\mpibindmain{MPI\_Type\_create\_hindexed}{(\gb{}int~count, const~int~array\_of\_blocklengths[], const~MPI\_Aint~array\_of\_displacements[], MPI\_Datatype~oldtype, MPI\_Datatype~*newtype)}
\MPIbindindex{MPI\_Type\_create\_hindexed\_c}
\mpibindmain{MPI\_Type\_create\_hindexed\_c}{(\gb{}MPI\_Count~count, const~MPI\_Count~array\_of\_blocklengths[], const~MPI\_Count~array\_of\_displacements[], MPI\_Datatype~oldtype, MPI\_Datatype~*newtype)}
\mpifnewbindmain{MPI\_Type\_create\_hindexed}{(count, array\_of\_blocklengths, array\_of\_displacements, oldtype, newtype, ierror) \fargs INTEGER, INTENT(IN)\ ::\ \mbox{count}, \mbox{array\_of\_blocklengths(count)}\fnarg{}INTEGER(KIND=MPI\_ADDRESS\_KIND), INTENT(IN)\ ::\ \mbox{array\_of\_displacements(count)}\fnarg{}TYPE(MPI\_Datatype), INTENT(IN)\ ::\ \mbox{oldtype}\fnarg{}TYPE(MPI\_Datatype), INTENT(OUT)\ ::\ \mbox{newtype}\fnfinalarg{}INTEGER, OPTIONAL, INTENT(OUT)\ ::\ \mbox{ierror}}
\mpifnewbindmain{MPI\_Type\_create\_hindexed}{(count, array\_of\_blocklengths, array\_of\_displacements, oldtype, newtype, ierror) !(\_c) \fargs INTEGER(KIND=MPI\_COUNT\_KIND), INTENT(IN)\ ::\ \mbox{count}, \mbox{array\_of\_blocklengths(count)}, \mbox{array\_of\_displacements(count)}\fnarg{}TYPE(MPI\_Datatype), INTENT(IN)\ ::\ \mbox{oldtype}\fnarg{}TYPE(MPI\_Datatype), INTENT(OUT)\ ::\ \mbox{newtype}\fnfinalarg{}INTEGER, OPTIONAL, INTENT(OUT)\ ::\ \mbox{ierror}}
\mpifbindmain{MPI\_TYPE\_CREATE\_HINDEXED}{(COUNT, ARRAY\_OF\_BLOCKLENGTHS, ARRAY\_OF\_DISPLACEMENTS, OLDTYPE, NEWTYPE, IERROR) \fargs INTEGER \mbox{COUNT}, \mbox{ARRAY\_OF\_BLOCKLENGTHS(*)}, \mbox{OLDTYPE}, \mbox{NEWTYPE}, \mbox{IERROR}\fnfinalarg{}INTEGER(KIND=MPI\_ADDRESS\_KIND) \mbox{ARRAY\_OF\_DISPLACEMENTS(*)}}

Assume that \mpiarg{oldtype} has type map,
\begin{displaymath}
\{ (\mword{type}_0,\mword{disp}_0), \ldots, (\mword{type}_{n-1}, \mword{disp}_{n-1}) \} ,
\end{displaymath}
with extent $ex$.
Let \mpiarg{B} be the \mpiarg{array\_of\_blocklengths} argument and
\mpiarg{D} be the
\mpiarg{array\_of\_displacements} argument. The newly created datatype
has a type map with
$n \cdot \sum_{i=0}^{\mscode{count}-1}
\mpiarg{B[i]}$ entries:
\begin{displaymath}
\{
(\mword{type}_0, \mword{disp}_0 + \mpiarg{D[0]}  ) , \ldots ,
(\mword{type}_{n-1} , \mword{disp}_{n-1} + \mpiarg{D[0]} ) , \ldots ,
\end{displaymath}
\begin{displaymath}
(\mword{type}_0 , \mword{disp}_0 + \mpiarg{D[0]} +(\mpiarg{B[0]} -1) \cdot ex) ,\ldots,
\end{displaymath}
\begin{displaymath}
(\mword{type}_{n-1} , \mword{disp}_{n-1} + \mpiarg{D[0]} +(\mpiarg{B[0]} -1) \cdot ex ) ,
\ldots,
\end{displaymath}
\begin{displaymath}
(\mword{type}_0, \mword{disp}_0 + \mpiarg{D[count-1]} ) , \ldots ,
(\mword{type}_{n-1} , \mword{disp}_{n-1} + \mpiarg{D[count-1]} ) , \ldots ,
\end{displaymath}
\begin{displaymath}
(\mword{type}_0 , \mword{disp}_0 + \mpiarg{D[count-1]} +(\mpiarg{B[count-1]} -1) \cdot ex)
,\ldots,
\end{displaymath}
\begin{displaymath}
(\mword{type}_{n-1} , \mword{disp}_{n-1} + \mpiarg{D[count-1]} +(\mpiarg{B[count-1]} -1)
\cdot ex )
\} .
\end{displaymath}

\paraheading{Indexed\_block.}
This procedure is the same as \mpifunc{MPI\_TYPE\_INDEXED} except that the
blocklength is the same for all blocks.
There are many codes using indirect addressing arising from
unstructured grids where the blocksize is always 1 (gather/scatter).  The
following convenience procedure allows for constant blocksize and arbitrary
displacements.


\begin{funcdef}{MPI\_TYPE\_CREATE\_INDEXED\_BLOCK}{(\mbox{count}, \mbox{blocklength}, \mbox{array\_of\_displacements}, \mbox{oldtype}, \mbox{newtype})}
\funcarg{\IN}{count}{number of blocks---also number of entries in \mpiarg{array\_of\_displacements} (nonnegative integer)}
\funcarg{\IN}{blocklength}{number of elements in each block (nonnegative integer)}
\funcarg{\IN}{array\_of\_displacements}{array of displacements, in multiples of \mpiarg{oldtype} (array of integers)}
\funcarg{\IN}{oldtype}{old datatype (handle)}
\funcarg{\OUT}{newtype}{new datatype (handle)}
\end{funcdef}
\mpibindmain{MPI\_Type\_create\_indexed\_block}{(\gb{}int~count, int~blocklength, const~int~array\_of\_displacements[], MPI\_Datatype~oldtype, MPI\_Datatype~*newtype)}
\MPIbindindex{MPI\_Type\_create\_indexed\_block\_c}
\mpibindmain{MPI\_Type\_create\_indexed\_block\_c}{(\gb{}MPI\_Count~count, MPI\_Count~blocklength, const~MPI\_Count~array\_of\_displacements[], MPI\_Datatype~oldtype, MPI\_Datatype~*newtype)}
\mpifnewbindmain{MPI\_Type\_create\_indexed\_block}{(count, blocklength, array\_of\_displacements, oldtype, newtype, ierror) \fargs INTEGER, INTENT(IN)\ ::\ \mbox{count}, \mbox{blocklength}, \mbox{array\_of\_displacements(count)}\fnarg{}TYPE(MPI\_Datatype), INTENT(IN)\ ::\ \mbox{oldtype}\fnarg{}TYPE(MPI\_Datatype), INTENT(OUT)\ ::\ \mbox{newtype}\fnfinalarg{}INTEGER, OPTIONAL, INTENT(OUT)\ ::\ \mbox{ierror}}
\mpifnewbindmain{MPI\_Type\_create\_indexed\_block}{(count, blocklength, array\_of\_displacements, oldtype, newtype, ierror) !(\_c) \fargs INTEGER(KIND=MPI\_COUNT\_KIND), INTENT(IN)\ ::\ \mbox{count}, \mbox{blocklength}, \mbox{array\_of\_displacements(count)}\fnarg{}TYPE(MPI\_Datatype), INTENT(IN)\ ::\ \mbox{oldtype}\fnarg{}TYPE(MPI\_Datatype), INTENT(OUT)\ ::\ \mbox{newtype}\fnfinalarg{}INTEGER, OPTIONAL, INTENT(OUT)\ ::\ \mbox{ierror}}
\mpifbindmain{MPI\_TYPE\_CREATE\_INDEXED\_BLOCK}{(COUNT, BLOCKLENGTH, ARRAY\_OF\_DISPLACEMENTS, OLDTYPE, NEWTYPE, IERROR) \fargs \fnfinalarg{}INTEGER \mbox{COUNT}, \mbox{BLOCKLENGTH}, \mbox{ARRAY\_OF\_DISPLACEMENTS(*)}, \mbox{OLDTYPE}, \mbox{NEWTYPE}, \mbox{IERROR}}

\paraheading{Hindexed\_block.}
The procedure \mpifunc{MPI\_TYPE\_CREATE\_HINDEXED\_BLOCK} is identical
to \mpifunc{MPI\_TYPE\_CREATE\_INDEXED\_BLOCK}, except that block
displacements in \mpiarg{array\_of\_displacements} are specified in
bytes, rather than in multiples of the \mpiarg{oldtype} extent.

\begin{funcdef}{MPI\_TYPE\_CREATE\_HINDEXED\_BLOCK}{(\mbox{count}, \mbox{blocklength}, \mbox{array\_of\_displacements}, \mbox{oldtype}, \mbox{newtype})}
\funcarg{\IN}{count}{number of blocks---also number of entries in \mpiarg{array\_of\_displacements} (nonnegative integer)}
\funcarg{\IN}{blocklength}{number of elements in each block (nonnegative integer)}
\funcarg{\IN}{array\_of\_displacements}{byte displacement of each block (array of integers)}
\funcarg{\IN}{oldtype}{old datatype (handle)}
\funcarg{\OUT}{newtype}{new datatype (handle)}
\end{funcdef}
\mpibindmain{MPI\_Type\_create\_hindexed\_block}{(\gb{}int~count, int~blocklength, const~MPI\_Aint~array\_of\_displacements[], MPI\_Datatype~oldtype, MPI\_Datatype~*newtype)}
\MPIbindindex{MPI\_Type\_create\_hindexed\_block\_c}
\mpibindmain{MPI\_Type\_create\_hindexed\_block\_c}{(\gb{}MPI\_Count~count, MPI\_Count~blocklength, const~MPI\_Count~array\_of\_displacements[], MPI\_Datatype~oldtype, MPI\_Datatype~*newtype)}
\mpifnewbindmain{MPI\_Type\_create\_hindexed\_block}{(count, blocklength, array\_of\_displacements, oldtype, newtype, ierror) \fargs INTEGER, INTENT(IN)\ ::\ \mbox{count}, \mbox{blocklength}\fnarg{}INTEGER(KIND=MPI\_ADDRESS\_KIND), INTENT(IN)\ ::\ \mbox{array\_of\_displacements(count)}\fnarg{}TYPE(MPI\_Datatype), INTENT(IN)\ ::\ \mbox{oldtype}\fnarg{}TYPE(MPI\_Datatype), INTENT(OUT)\ ::\ \mbox{newtype}\fnfinalarg{}INTEGER, OPTIONAL, INTENT(OUT)\ ::\ \mbox{ierror}}
\mpifnewbindmain{MPI\_Type\_create\_hindexed\_block}{(count, blocklength, array\_of\_displacements, oldtype, newtype, ierror) !(\_c) \fargs INTEGER(KIND=MPI\_COUNT\_KIND), INTENT(IN)\ ::\ \mbox{count}, \mbox{blocklength}, \mbox{array\_of\_displacements(count)}\fnarg{}TYPE(MPI\_Datatype), INTENT(IN)\ ::\ \mbox{oldtype}\fnarg{}TYPE(MPI\_Datatype), INTENT(OUT)\ ::\ \mbox{newtype}\fnfinalarg{}INTEGER, OPTIONAL, INTENT(OUT)\ ::\ \mbox{ierror}}
\mpifbindmain{MPI\_TYPE\_CREATE\_HINDEXED\_BLOCK}{(COUNT, BLOCKLENGTH, ARRAY\_OF\_DISPLACEMENTS, OLDTYPE, NEWTYPE, IERROR) \fargs INTEGER \mbox{COUNT}, \mbox{BLOCKLENGTH}, \mbox{OLDTYPE}, \mbox{NEWTYPE}, \mbox{IERROR}\fnfinalarg{}INTEGER(KIND=MPI\_ADDRESS\_KIND) \mbox{ARRAY\_OF\_DISPLACEMENTS(*)}}

\paraheading{Struct.}
\mpifunc{MPI\_TYPE\_CREATE\_STRUCT} is the most general type constructor.
It further generalizes
\mpifunc{MPI\_TYPE\_CREATE\_HINDEXED}
in that it allows each block to consist of replications of
different datatypes.

\cdeclindex{MPI\_Aint}%
\begin{funcdef}{MPI\_TYPE\_CREATE\_STRUCT}{(\mbox{count}, \mbox{array\_of\_blocklengths}, \mbox{array\_of\_displacements}, \mbox{array\_of\_types}, \mbox{newtype})}
\funcarg{\IN}{count}{number of blocks---also number of entries in arrays \mpiarg{array\_of\_types}, \mpiarg{array\_of\_displacements}, and \mpiarg{array\_of\_blocklengths} (nonnegative integer)}
\funcarg{\IN}{array\_of\_blocklengths}{number of elements in each block (array of nonnegative integers)}
\funcarg{\IN}{array\_of\_displacements}{byte displacement of each block (array of integers)}
\funcarg{\IN}{array\_of\_types}{type of elements in each block (array of handles)}
\funcarg{\OUT}{newtype}{new datatype (handle)}
\end{funcdef}
\mpibindmain{MPI\_Type\_create\_struct}{(\gb{}int~count, const~int~array\_of\_blocklengths[], const~MPI\_Aint~array\_of\_displacements[], const~MPI\_Datatype~array\_of\_types[], MPI\_Datatype~*newtype)}
\MPIbindindex{MPI\_Type\_create\_struct\_c}
\mpibindmain{MPI\_Type\_create\_struct\_c}{(\gb{}MPI\_Count~count, const~MPI\_Count~array\_of\_blocklengths[], const~MPI\_Count~array\_of\_displacements[], const~MPI\_Datatype~array\_of\_types[], MPI\_Datatype~*newtype)}
\mpifnewbindmain{MPI\_Type\_create\_struct}{(count, array\_of\_blocklengths, array\_of\_displacements, array\_of\_types, newtype, ierror) \fargs INTEGER, INTENT(IN)\ ::\ \mbox{count}, \mbox{array\_of\_blocklengths(count)}\fnarg{}INTEGER(KIND=MPI\_ADDRESS\_KIND), INTENT(IN)\ ::\ \mbox{array\_of\_displacements(count)}\fnarg{}TYPE(MPI\_Datatype), INTENT(IN)\ ::\ \mbox{array\_of\_types(count)}\fnarg{}TYPE(MPI\_Datatype), INTENT(OUT)\ ::\ \mbox{newtype}\fnfinalarg{}INTEGER, OPTIONAL, INTENT(OUT)\ ::\ \mbox{ierror}}
\mpifnewbindmain{MPI\_Type\_create\_struct}{(count, array\_of\_blocklengths, array\_of\_displacements, array\_of\_types, newtype, ierror) !(\_c) \fargs INTEGER(KIND=MPI\_COUNT\_KIND), INTENT(IN)\ ::\ \mbox{count}, \mbox{array\_of\_blocklengths(count)}, \mbox{array\_of\_displacements(count)}\fnarg{}TYPE(MPI\_Datatype), INTENT(IN)\ ::\ \mbox{array\_of\_types(count)}\fnarg{}TYPE(MPI\_Datatype), INTENT(OUT)\ ::\ \mbox{newtype}\fnfinalarg{}INTEGER, OPTIONAL, INTENT(OUT)\ ::\ \mbox{ierror}}
\mpifbindmain{MPI\_TYPE\_CREATE\_STRUCT}{(COUNT, ARRAY\_OF\_BLOCKLENGTHS, ARRAY\_OF\_DISPLACEMENTS, ARRAY\_OF\_TYPES, NEWTYPE, IERROR) \fargs INTEGER \mbox{COUNT}, \mbox{ARRAY\_OF\_BLOCKLENGTHS(*)}, \mbox{ARRAY\_OF\_TYPES(*)}, \mbox{NEWTYPE}, \mbox{IERROR}\fnfinalarg{}INTEGER(KIND=MPI\_ADDRESS\_KIND) \mbox{ARRAY\_OF\_DISPLACEMENTS(*)}}

\begin{example}
\label{pt2pt-exV}
\exindex{Typemap}%
\Exindex{Neutral}{Typemap for create struct}{MPI\_TYPE\_CREATE\_STRUCT}%
Let \mpiarg{type1} have type map,
\[
\{ (\ctype{double}, 0), (\ctype{char}, 8) \} ,
\]
with extent 16.
Let \mpicode{B = (2, 1, 3)}, \mpicode{D = (0, 16, 26)},
and \mpicode{T = (MPI\_FLOAT, type1, MPI\_CHAR)}.  Then a call to
\mpiskipfunc{MPI\_TYPE\_CREATE\_STRUCT(3, B, D, T, newtype)}\mpifuncindex{MPI\_TYPE\_CREATE\_STRUCT} returns
a datatype with type map,
\begin{displaymath}
\{
(\ctype{float}, 0), (\ctype{float}, 4), (\ctype{double}, 16), (\ctype{char},
24), (\ctype{char}, 26), (\ctype{char}, 27), (\ctype{char}, 28)
\} .
\end{displaymath}
That is, two copies of \type{MPI\_FLOAT} starting at 0, followed by
one copy of \mpiarg{type1} starting at 16, followed by three copies of
\type{MPI\_CHAR}, starting at 26.
In this example, we assume that a float occupies four bytes.
\end{example}

In general,
let \mpiarg{T} be the \mpiarg{array\_of\_types} argument, where \mpiarg{T[i]}
is a handle to,
\begin{displaymath}
\mword{typemap}_i = \{ (\mword{type}_0^i , \mword{disp}_0^i ) , \ldots , (\mword{type}_{n_{i}-1}^i ,
\mword{disp}_{n_{i}-1}^i ) \} ,
\end{displaymath}
with extent $ex_i$.
Let
\mpiarg{B} be the \mpiarg{array\_of\_blocklength} argument and \mpiarg{D} be
the \mpiarg{array\_of\_displacements} argument.  %mansplit
Let \mpicode{c} be the
\mpicode{count} argument.
Then the newly created datatype has a type map with
$\sum_{i=0}^{\mscode{c}-1}\mpiarg{B[i]} \cdot n_i$
entries:
\begin{displaymath}
\{
(\mword{type}_0^0 , \mword{disp}_0^0 +\mpiarg{D[0]}) , \ldots , (\mword{type}_{n_0}^0 , \mword{disp}_{n_0}^0 +
\mpiarg{D[0]} ) ,
\ldots ,
\end{displaymath}
\begin{displaymath}
(\mword{type}_0^0 , \mword{disp}_0^0 + \mpiarg{D[0]} + (\mpiarg{B[0]}-1) \cdot ex_0 ) , \ldots
,
(\mword{type}_{n_0}^0 , \mword{disp}_{n_0}^0 + \mpiarg{D[0]} + (\mpiarg{B[0]-1)} \cdot
ex_0 )
, \ldots ,
\end{displaymath}
\begin{displaymath}
(\mword{type}_0^{\mscode{c}-1} , \mword{disp}_0^{\mscode{c}-1} +\mpiarg{D[c-1]}) , \ldots ,
(\mword{type}_{{n_{\mscode{c}-1}}-1}^{\mscode{c}-1} ,
\mword{disp}_{{n_{\mscode{c}-1}}-1}^{\mscode{c}-1} + \mpiarg{D[c-1]} ) ,
\ldots ,
\end{displaymath}
\begin{displaymath}
(\mword{type}_0^{\mscode{c}-1} , \mword{disp}_0^{\mscode{c}-1} +
\mpiarg{D[c-1]} + (\mpiarg{B[c-1]}-1) \cdot ex_{\mscode{c}-1} ) ,
\ldots ,
\end{displaymath}
\begin{displaymath}
(\mword{type}_{{n_{\mscode{c}-1}}-1}^{\mscode{c}-1} ,
\mword{disp}_{{n_{\mscode{c}-1}}-1}^{\mscode{c}-1} + \mpiarg{D[c-1]} +
(\mpiarg{B[c-1]-1)} \cdot ex_{\mscode{c}-1} )
\} .
\end{displaymath}

A call to
\mpifunc{MPI\_TYPE\_CREATE\_HINDEXED}\mpicode{(count, B, D, oldtype, newtype)}
is equivalent to a call to
\mpifunc{MPI\_TYPE\_CREATE\_STRUCT}\mpicode{(count, B, D, T, newtype)},
where each entry of \mpiarg{T} is equal to \mpiarg{oldtype}.

\subsection{Subarray Datatype Constructor}
%----------------------------------------------------------
\label{sec:io-const-array}

\begin{funcdef}{MPI\_TYPE\_CREATE\_SUBARRAY}{(\mbox{ndims}, \mbox{array\_of\_sizes}, \mbox{array\_of\_subsizes}, \mbox{array\_of\_starts}, \mbox{order}, \mbox{oldtype}, \mbox{newtype})}
\funcarg{\IN}{ndims}{number of array dimensions (positive integer)}
\funcarg{\IN}{array\_of\_sizes}{number of elements of type \mpiarg{oldtype} in each dimension of the full array (array of positive integers)}
\funcarg{\IN}{array\_of\_subsizes}{number of elements of type \mpiarg{oldtype} in each dimension of the subarray (array of positive integers)}
\funcarg{\IN}{array\_of\_starts}{starting coordinates of the subarray in each dimension (array of nonnegative integers)}
\funcarg{\IN}{order}{array storage order flag (state)}
\funcarg{\IN}{oldtype}{old datatype (handle)}
\funcarg{\OUT}{newtype}{new datatype (handle)}
\end{funcdef}
\mpibindmain{MPI\_Type\_create\_subarray}{(\gb{}int~ndims, const~int~array\_of\_sizes[], const~int~array\_of\_subsizes[], const~int~array\_of\_starts[], int~order, MPI\_Datatype~oldtype, MPI\_Datatype~*newtype)}
\MPIbindindex{MPI\_Type\_create\_subarray\_c}
\mpibindmain{MPI\_Type\_create\_subarray\_c}{(\gb{}int~ndims, const~MPI\_Count~array\_of\_sizes[], const~MPI\_Count~array\_of\_subsizes[], const~MPI\_Count~array\_of\_starts[], int~order, MPI\_Datatype~oldtype, MPI\_Datatype~*newtype)}
\mpifnewbindmain{MPI\_Type\_create\_subarray}{(ndims, array\_of\_sizes, array\_of\_subsizes, array\_of\_starts, order, oldtype, newtype, ierror) \fargs INTEGER, INTENT(IN)\ ::\ \mbox{ndims}, \mbox{array\_of\_sizes(ndims)}, \mbox{array\_of\_subsizes(ndims)}, \mbox{array\_of\_starts(ndims)}, \mbox{order}\fnarg{}TYPE(MPI\_Datatype), INTENT(IN)\ ::\ \mbox{oldtype}\fnarg{}TYPE(MPI\_Datatype), INTENT(OUT)\ ::\ \mbox{newtype}\fnfinalarg{}INTEGER, OPTIONAL, INTENT(OUT)\ ::\ \mbox{ierror}}
\mpifnewbindmain{MPI\_Type\_create\_subarray}{(ndims, array\_of\_sizes, array\_of\_subsizes, array\_of\_starts, order, oldtype, newtype, ierror) !(\_c) \fargs INTEGER, INTENT(IN)\ ::\ \mbox{ndims}, \mbox{order}\fnarg{}INTEGER(KIND=MPI\_COUNT\_KIND), INTENT(IN)\ ::\ \mbox{array\_of\_sizes(ndims)}, \mbox{array\_of\_subsizes(ndims)}, \mbox{array\_of\_starts(ndims)}\fnarg{}TYPE(MPI\_Datatype), INTENT(IN)\ ::\ \mbox{oldtype}\fnarg{}TYPE(MPI\_Datatype), INTENT(OUT)\ ::\ \mbox{newtype}\fnfinalarg{}INTEGER, OPTIONAL, INTENT(OUT)\ ::\ \mbox{ierror}}
\mpifbindmain{MPI\_TYPE\_CREATE\_SUBARRAY}{(NDIMS, ARRAY\_OF\_SIZES, ARRAY\_OF\_SUBSIZES, ARRAY\_OF\_STARTS, ORDER, OLDTYPE, NEWTYPE, IERROR) \fargs \fnfinalarg{}INTEGER \mbox{NDIMS}, \mbox{ARRAY\_OF\_SIZES(*)}, \mbox{ARRAY\_OF\_SUBSIZES(*)}, \mbox{ARRAY\_OF\_STARTS(*)}, \mbox{ORDER}, \mbox{OLDTYPE}, \mbox{NEWTYPE}, \mbox{IERROR}}

The subarray type constructor creates an \MPI/ datatype
describing an
$n$-dimensional
subarray of an
$n$-dimensional
array.
The subarray may be situated anywhere within the full array,
and may be of any nonzero size up to the size of the larger array
as long as it is confined within this array.
This type constructor facilitates creating filetypes to access
arrays distributed in blocks among processes to a
single file that contains the global array,
see \MPI/ I/O, especially \sectionref{subsec:io-2:definitions}.

This type constructor can handle arrays with an arbitrary number of
dimensions and works for both C and Fortran ordered matrices
(i.e., row-major or column-major).  Note that a C program may use Fortran
order and a Fortran program may use C order.


The \mpiarg{ndims} parameter specifies the number of dimensions in the full
data array and gives the number of elements in
\mpiarg{array\_of\_sizes}, \mpiarg{array\_of\_subsizes},
and \mpiarg{array\_of\_starts}.

The number of elements of type \mpiarg{oldtype} in
each dimension of the
$n$-dimensional
array and the
requested subarray are specified by
\mpiarg{array\_of\_sizes} and \mpiarg{array\_of\_subsizes},
respectively.
For any dimension
\mpiarg{i},
it is erroneous to specify
\mpiarg{array\_of\_subsizes[i]} $<$ 1 or
\mpiarg{array\_of\_subsizes[i]} $>$ \mpiarg{array\_of\_sizes[i]}.

The \mpiarg{array\_of\_starts} contains the
starting coordinates of each
dimension of the subarray.
Arrays are assumed to be indexed starting from zero.
For any dimension $i$, it is erroneous to specify
\mpiarg{array\_of\_starts[i]} $<$ 0 or
\mpiarg{array\_of\_starts[i]}
$>$ (\mpiarg{array\_of\_sizes[i]} $-$ \mpiarg{array\_of\_subsizes[i]}).

\begin{users}
In a Fortran program with arrays indexed starting from 1,
if the starting coordinate of a particular dimension
of the subarray is \mpiarg{n},
then the entry in \mpiarg{array\_of\_starts}
for that dimension is \mpiarg{n-1}.
\end{users}

The \mpiarg{order} argument specifies the storage order for the
subarray as well as the full array.
It must be set to one
of the following:
\begin{constlist}
\constmainitem{MPI\_ORDER\_C}{The ordering used by C arrays, (i.e., row-major order).}
\constmainitem{MPI\_ORDER\_FORTRAN}{The ordering used by Fortran arrays,
(i.e., column-major order).}
\end{constlist}

A \mpiarg{ndims}-dimensional subarray (\mpiarg{newtype}) with no
extra padding can be defined by the function Subarray() as
follows:
\begin{eqnarray*}
\mpiarg{newtype} & = &  \mbox{Subarray}( \mword{ndims},
                        \{\mword{size}_0, \mword{size}_1,\ldots,\mword{size}_{\msword{ndims}-1}\},        \\
                 &   &  \{\mword{subsize}_0, \mword{subsize}_1,\ldots,\mword{subsize}_{\msword{ndims}-1}\}, \\
                 &   &  \{\mword{start}_0, \mword{start}_1,\ldots,\mword{start}_{\msword{ndims}-1}\},
                        {\mpiarg{oldtype}} )
\end{eqnarray*}

Let the typemap of \mpiarg{oldtype} have the form:
\begin{displaymath}
\{(\mword{type}_0,\mword{disp}_0),(\mword{type}_1,\mword{disp}_1),\ldots,(\mword{type}_{n-1},\mword{disp}_{n-1})\}
\end{displaymath}
where $\mword{type}_i$ is a predefined \MPI/ datatype, and let $ex$ be the
extent of \mpiargshort{oldtype}.  Then we define the Subarray() function
recursively using the following three equations.
Equation~\ref{eq:subarray-base} defines the base step.
Equation~\ref{eq:subarray-fortran} defines the recursion step when
\mpiarg{order} = \mpiconst{MPI\_ORDER\_FORTRAN}, and
Equation~\ref{eq:subarray-c} defines the recursion step when
\mpiarg{order} = \mpiconst{MPI\_ORDER\_C}.
These equations use the conceptual datatypes \mpiublb{lb\_marker}
and \mpiublb{ub\_marker}; see \sectionref{subsec:pt2pt-markers} for details.

\begin{eqnarray}
\lefteqn{\mbox{Subarray}(1,\{\mword{size}_0\},\{\mword{subsize}_0\},\{\mword{start}_0\},}
\label{eq:subarray-base} \\
& & \quad \{(\mword{type}_0,\mword{disp}_0),(\mword{type}_1,\mword{disp}_1),\ldots,(\mword{type}_{n-1},\mword{disp}_{n-1})\}) \nonumber \\
& = & \{(\mpiublb{lb\_marker},0), \nonumber \\
& & (\mword{type}_0,\mword{disp}_0+\mword{start}_0 \times ex),\ldots,(\mword{type}_{n-1},
        \mword{disp}_{n-1} + \mword{start}_0 \times ex), \nonumber \\
& & (\mword{type}_0,\mword{disp}_0+(\mword{start}_0 + 1)\times ex),\ldots,(\mword{type}_{n-1},
\nonumber\\
& & \hspace{.5in}\mword{disp}_{n-1} + (\mword{start}_0+1) \times ex), \ldots \nonumber \\
& & (\mword{type}_0,\mword{disp}_0+(\mword{start}_0 + \mword{subsize}_0 - 1)\times ex),\ldots,
\nonumber \\
& & \hspace{.5in}(\mword{type}_{n-1},\mword{disp}_{n-1} + (\mword{start}_0+\mword{subsize}_0 - 1) \times ex),
\nonumber \\
& & (\mpiublb{ub\_marker}, \mword{size}_0 \times ex) \} \nonumber \\
& & \nonumber \\
\lefteqn{\mbox{Subarray}( \mword{ndims},
\{\mword{size}_0, \mword{size}_1,\ldots,\mword{size}_{\msword{ndims}-1}\},}       \label{eq:subarray-fortran} \\
& & \quad\{\mword{subsize}_0, \mword{subsize}_1,\ldots,\mword{subsize}_{\msword{ndims}-1}\}, \nonumber \\
& & \quad \{\mword{start}_0, \mword{start}_1,\ldots,\mword{start}_{\msword{ndims}-1}\},\mpiarg{oldtype} )
\nonumber \\
& = & \mbox{Subarray}(ndims-1,\{\mword{size}_1, \mword{size}_2,\ldots,\mword{size}_{\msword{ndims}-1}\},
\nonumber \\
& &\quad\{\mword{subsize}_1, \mword{subsize}_2,\ldots,\mword{subsize}_{\msword{ndims}-1}\}, \nonumber\\
& &\quad\{\mword{start}_1, \mword{start}_2,\ldots,\mword{start}_{\msword{ndims}-1}\}, \nonumber \\
& & \hspace{.5in}\mbox{Subarray}(1,\{\mword{size}_0\},\{\mword{subsize}_0\},\{\mword{start}_0\},
        \mpiarg{oldtype}))\nonumber \\
& & \nonumber \\
\lefteqn{\mbox{Subarray}( ndims,
\{\mword{size}_0, \mword{size}_1,\ldots,\mword{size}_{\msword{ndims}-1}\},} \label{eq:subarray-c}         \\
& & \quad\{\mword{subsize}_0, \mword{subsize}_1,\ldots,\mword{subsize}_{\msword{ndims}-1}\}, \nonumber \\
& & \quad\{\mword{start}_0, \mword{start}_1,\ldots,\mword{start}_{\msword{ndims}-1}\},\mpiarg{oldtype} )
\nonumber  \\
& = & \mbox{Subarray}(\mword{ndims}-1,\{\mword{size}_0, \mword{size}_1,\ldots,\mword{size}_{\msword{ndims}-2}\},
\nonumber \\
& &\quad\{\mword{subsize}_0, \mword{subsize}_1,\ldots,\mword{subsize}_{\msword{ndims}-2}\}, \nonumber \\
& &\quad\{\mword{start}_0, \mword{start}_1,\ldots,\mword{start}_{\msword{ndims}-2}\}, \nonumber \\
& & \hspace{.5in}\mbox{Subarray}(1,\{\mword{size}_{\msword{ndims}-1}\},\{\mword{subsize}_{\msword{ndims}-1}\},
        \{\mword{start}_{\msword{ndims}-1}\},\mpiarg{oldtype})) \nonumber
\end{eqnarray}
For an example use of \mpifunc{MPI\_TYPE\_CREATE\_SUBARRAY} in the
context of I/O see Section~\ref{sec:io-ex-subarray}.

\subsection{Distributed Array Datatype Constructor}
%---------------------------------------------------
\label{sec:io-const-darray}

The distributed array type constructor supports
HPF-like \cite{hpf-handbook} data distributions.
However, unlike in HPF, the storage order may be specified for C arrays
as well as for Fortran arrays.

\begin{users}
One can create an HPF-like file view using this type constructor
as follows.
Complementary filetypes are created by having every process of a group
call this constructor with identical arguments
(with the exception of \mpiarg{rank} which should be set appropriately).
These filetypes (along with identical \mpiarg{disp} and \mpiarg{etype})
are then used to define the view (via \mpifunc{MPI\_FILE\_SET\_VIEW}),
see \MPI/ I/O, especially \sectionref{subsec:io-2:definitions}
and \sectionref{sec:io-view}.
Using this view,
a collective data access operation (with identical offsets)
will yield an HPF-like distribution pattern.
\end{users}

\begin{funcdef}{MPI\_TYPE\_CREATE\_DARRAY}{(\mbox{size}, \mbox{rank}, \mbox{ndims}, \mbox{array\_of\_gsizes}, \mbox{array\_of\_distribs}, \mbox{array\_of\_dargs}, \mbox{array\_of\_psizes}, \mbox{order}, \mbox{oldtype}, \mbox{newtype})}
\funcarg{\IN}{size}{size of process group (positive integer)}
\funcarg{\IN}{rank}{rank in process group (nonnegative integer)}
\funcarg{\IN}{ndims}{number of array dimensions as well as process grid dimensions (positive integer)}
\funcarg{\IN}{array\_of\_gsizes}{number of elements of type \mpiarg{oldtype} in each dimension of global array (array of positive integers)}
\funcarg{\IN}{array\_of\_distribs}{distribution of array in each dimension (array of states)}
\funcarg{\IN}{array\_of\_dargs}{distribution argument in each dimension (array of positive integers)}
\funcarg{\IN}{array\_of\_psizes}{size of process grid in each dimension (array of positive integers)}
\funcarg{\IN}{order}{array storage order flag (state)}
\funcarg{\IN}{oldtype}{old datatype (handle)}
\funcarg{\OUT}{newtype}{new datatype (handle)}
\end{funcdef}
\mpibindmain{MPI\_Type\_create\_darray}{(\gb{}int~size, int~rank, int~ndims, const~int~array\_of\_gsizes[], const~int~array\_of\_distribs[], const~int~array\_of\_dargs[], const~int~array\_of\_psizes[], int~order, MPI\_Datatype~oldtype, MPI\_Datatype~*newtype)}
\MPIbindindex{MPI\_Type\_create\_darray\_c}
\mpibindmain{MPI\_Type\_create\_darray\_c}{(\gb{}int~size, int~rank, int~ndims, const~MPI\_Count~array\_of\_gsizes[], const~int~array\_of\_distribs[], const~int~array\_of\_dargs[], const~int~array\_of\_psizes[], int~order, MPI\_Datatype~oldtype, MPI\_Datatype~*newtype)}
\mpifnewbindmain{MPI\_Type\_create\_darray}{(size, rank, ndims, array\_of\_gsizes, array\_of\_distribs, array\_of\_dargs, array\_of\_psizes, order, oldtype, newtype, ierror) \fargs INTEGER, INTENT(IN)\ ::\ \mbox{size}, \mbox{rank}, \mbox{ndims}, \mbox{array\_of\_gsizes(ndims)}, \mbox{array\_of\_distribs(ndims)}, \mbox{array\_of\_dargs(ndims)}, \mbox{array\_of\_psizes(ndims)}, \mbox{order}\fnarg{}TYPE(MPI\_Datatype), INTENT(IN)\ ::\ \mbox{oldtype}\fnarg{}TYPE(MPI\_Datatype), INTENT(OUT)\ ::\ \mbox{newtype}\fnfinalarg{}INTEGER, OPTIONAL, INTENT(OUT)\ ::\ \mbox{ierror}}
\mpifnewbindmain{MPI\_Type\_create\_darray}{(size, rank, ndims, array\_of\_gsizes, array\_of\_distribs, array\_of\_dargs, array\_of\_psizes, order, oldtype, newtype, ierror) !(\_c) \fargs INTEGER, INTENT(IN)\ ::\ \mbox{size}, \mbox{rank}, \mbox{ndims}, \mbox{array\_of\_distribs(ndims)}, \mbox{array\_of\_dargs(ndims)}, \mbox{array\_of\_psizes(ndims)}, \mbox{order}\fnarg{}INTEGER(KIND=MPI\_COUNT\_KIND), INTENT(IN)\ ::\ \mbox{array\_of\_gsizes(ndims)}\fnarg{}TYPE(MPI\_Datatype), INTENT(IN)\ ::\ \mbox{oldtype}\fnarg{}TYPE(MPI\_Datatype), INTENT(OUT)\ ::\ \mbox{newtype}\fnfinalarg{}INTEGER, OPTIONAL, INTENT(OUT)\ ::\ \mbox{ierror}}
\mpifbindmain{MPI\_TYPE\_CREATE\_DARRAY}{(SIZE, RANK, NDIMS, ARRAY\_OF\_GSIZES, ARRAY\_OF\_DISTRIBS, ARRAY\_OF\_DARGS, ARRAY\_OF\_PSIZES, ORDER, OLDTYPE, NEWTYPE, IERROR) \fargs \fnfinalarg{}INTEGER \mbox{SIZE}, \mbox{RANK}, \mbox{NDIMS}, \mbox{ARRAY\_OF\_GSIZES(*)}, \mbox{ARRAY\_OF\_DISTRIBS(*)}, \mbox{ARRAY\_OF\_DARGS(*)}, \mbox{ARRAY\_OF\_PSIZES(*)}, \mbox{ORDER}, \mbox{OLDTYPE}, \mbox{NEWTYPE}, \mbox{IERROR}}

\mpifunc{MPI\_TYPE\_CREATE\_DARRAY} can be used to generate
the datatypes corresponding to the distribution
of an \mpiarg{ndims}-dimensional array of \mpiarg{oldtype} elements
onto
an \mpiarg{ndims}-dimensional grid of logical processes.
Unused dimensions of \mpiarg{array\_of\_psizes} should be set to \mpicode{1}
(see \namedref{Example}{ex:io-hpf}).
For a call to \mpifunc{MPI\_TYPE\_CREATE\_DARRAY} to be correct,
the equation \( \prod_{i=0}^{ndims-1} array\_of\_psizes[i] = size \)
must be satisfied.
The ordering of processes in the process grid is assumed to be
row-major, as in the case of virtual Cartesian process topologies.
\begin{users}
For both Fortran and C arrays, the ordering of processes in the
process grid is assumed to be row-major. This is consistent with the
ordering used in virtual Cartesian process topologies in
\MPI/.
To create such virtual process topologies, or to find the coordinates
of a process in the process grid, etc., users may use the corresponding
process topology procedures,
see \namedref{Chapter}{chap:topol}.
\end{users}

Each dimension of the array
can be
distributed in one of three ways:
\begin{constlist}
\constmainitem{MPI\_DISTRIBUTE\_BLOCK}{Block distribution.}
\constmainitem{MPI\_DISTRIBUTE\_CYCLIC}{Cyclic distribution.}
\constmainitem{MPI\_DISTRIBUTE\_NONE}{Dimension not distributed.}
\end{constlist}

The constant \mpiconstmain{MPI\_DISTRIBUTE\_DFLT\_DARG} specifies
a default distribution argument.
The distribution argument for a dimension that is not distributed
is ignored.
For any dimension
\mpiarg{i}
in which the distribution
is \mpiconst{MPI\_DISTRIBUTE\_BLOCK},
it is erroneous
to specify
\mpiarg{array\_of\_dargs[i]} $*$ \mpiarg{array\_of\_psizes[i]}
$<$ \mpiarg{array\_of\_gsizes[i]}.

For example, the HPF layout \code{ARRAY(CYCLIC(15))}
corresponds to \mpiconst{MPI\_DISTRIBUTE\_CYCLIC}
with a distribution argument of 15, and the HPF layout \gb\code{ARRAY(BLOCK)}
corresponds to
\mpiconst{MPI\_DISTRIBUTE\_BLOCK} with a distribution argument of
\mpiconst{MPI\_DISTRIBUTE\_DFLT\_DARG}.

The \mpiarg{order} argument is used as in \mpifunc{MPI\_TYPE\_CREATE\_SUBARRAY} to
specify the storage order.
Therefore, arrays described by this type constructor may be
stored in Fortran (column-major) or C (row-major) order.
Valid values for \mpiarg{order} are
\mpiconst{MPI\_ORDER\_FORTRAN} and \mpiconst{MPI\_ORDER\_C}.

This routine creates a new \MPI/ datatype with a typemap defined in
terms of a function called ``\mpicode{cyclic()}'' (see below).

Without loss of generality, it suffices to define the typemap
for the \mpiconst{MPI\_DISTRIBUTE\_CYCLIC} case where
\mpiconst{MPI\_DISTRIBUTE\_DFLT\_DARG} is not used.

\mpiconst{MPI\_DISTRIBUTE\_BLOCK} and \mpiconst{MPI\_DISTRIBUTE\_NONE}
can be reduced to the \mpiconst{MPI\_DISTRIBUTE\_CYCLIC} case
for dimension
\mpiarg{i}
as follows.

\mpiconst{MPI\_DISTRIBUTE\_BLOCK} with
\mpiarg{array\_of\_dargs[i]} equal to \mpiconst{MPI\_DISTRIBUTE\_DFLT\_DARG}
is equivalent to
\mpiconst{MPI\_DISTRIBUTE\_CYCLIC}
with \mpiarg{array\_of\_dargs[i]} set to
\[
(\mpiarg{array\_of\_gsizes[i]} + \mpiarg{array\_of\_psizes[i]} - 1)
        / \mpiarg{array\_of\_psizes[i]}.
\]
If \mpiarg{array\_of\_dargs[i]} is not \mpiconst{MPI\_DISTRIBUTE\_DFLT\_DARG},
then \mpiconst{MPI\_DISTRIBUTE\_BLOCK} and \mpiconst{MPI\_DISTRIBUTE\_CYCLIC}
are equivalent.

\mpiconst{MPI\_DISTRIBUTE\_NONE} is equivalent to
\mpiconst{MPI\_DISTRIBUTE\_CYCLIC}
with \mpiarg{array\_of\_dargs[i]} set to \mpiarg{array\_of\_gsizes[i]}.

Finally,
\mpiconst{MPI\_DISTRIBUTE\_CYCLIC} with
\mpiarg{array\_of\_dargs[i]} equal to \mpiconst{MPI\_DISTRIBUTE\_DFLT\_DARG}
is equivalent to
\mpiconst{MPI\_DISTRIBUTE\_CYCLIC} with
\mpiarg{array\_of\_dargs[i]} set to 1.

For \mpiconst{MPI\_ORDER\_FORTRAN},
an \mpiarg{ndims}-dimensional distributed array (\mpiargshort{newtype})
is defined by the following code fragment:

%%HEADER
%%FRAGMENT
%%DECL: MPI_Datatype oldtypes[10], oldtype, newtype;
%%DECL: int i, ndims, r[10], array_of_psizes[10];
%%DECL: int array_of_dargs[10], array_of_gsizes[10];
%%DECL: extern MPI_Datatype cyclic(int,int,int,int,MPI_Datatype);
%%TAIL: if (!newtype) abort();
%%ENDHEADER
\begin{lstlisting}[language=C]
oldtypes[0] = oldtype;
for (i = 0; i < ndims; i++) {
    oldtypes[i+1] = cyclic(array_of_dargs[i],
                           array_of_gsizes[i],
                           r[i],
                           array_of_psizes[i],
                           oldtypes[i]);
}
newtype = oldtypes[ndims];
\end{lstlisting}

For \mpiconst{MPI\_ORDER\_C}, the code is:

%%HEADER
%%FRAGMENT
%%DECL: MPI_Datatype oldtypes[10], oldtype, newtype;
%%DECL: int i, ndims, r[10], array_of_psizes[10];
%%DECL: int array_of_dargs[10], array_of_gsizes[10];
%%DECL: extern MPI_Datatype cyclic(int,int,int,int,MPI_Datatype);
%%TAIL: if (!newtype) abort();
%%ENDHEADER
\begin{lstlisting}[language=C]
oldtypes[0] = oldtype;
for (i = 0; i < ndims; i++) {
    oldtypes[i+1] = cyclic(array_of_dargs[ndims - i - 1],
                           array_of_gsizes[ndims - i - 1],
                           r[ndims - i - 1],
                           array_of_psizes[ndims - i - 1],
                           oldtypes[i]);
}
newtype = oldtypes[ndims];
\end{lstlisting}
where \code{r[i]} is the position of the process (with rank \mpiarg{rank})
in the process grid at dimension \code{i}.
The values of \code{r[i]} are given by the following code fragment:

%%HEADER
%%FRAGMENT
%%DECL: int t_rank, rank, t_size, i, ndims, array_of_psizes[10], r[10];
%%ENDHEADER
\begin{lstlisting}[language=C]
t_rank = rank;
t_size = 1;
for (i = 0; i < ndims; i++)
    t_size *= array_of_psizes[i];
for (i = 0; i < ndims; i++) {
    t_size = t_size / array_of_psizes[i];
    r[i] = t_rank / t_size;
    t_rank = t_rank % t_size;
}
\end{lstlisting}

\noindent Let the typemap of \mpiarg{oldtype} have the form:
\begin{displaymath}
\{(\mword{type}_0,\mword{disp}_0),(\mword{type}_1,\mword{disp}_1),\ldots,(\mword{type}_{n-1},\mword{disp}_{n-1})\}
\end{displaymath}
where $\mword{type}_i$ is a predefined \MPI/ datatype, and let $ex$ be the
extent of \mpiargshort{oldtype}.
The following function uses the conceptual datatypes \mpiublb{lb\_marker}
and \mpiublb{ub\_marker}, see \sectionref{subsec:pt2pt-markers} for details.

Given the above, the function cyclic() is defined as follows:
\begin{eqnarray*}
\lefteqn{\mbox{cyclic}(\mword{darg}, gsize, r, \mword{psize}, \mpiarg{oldtype})} \\
&=& \{ (\mpiublb{lb\_marker}, 0), \\
& &  (\mword{type}_0, \mword{disp}_0 + r \times \mword{darg} \times ex), \ldots , \\
& & \hspace{.5in}  (\mword{type}_{n-1}, \mword{disp}_{n-1} + r \times \mword{darg} \times ex), \\
& &  (\mword{type}_0, \mword{disp}_0 + (r \times \mword{darg} + 1) \times ex), \ldots , \\
& & \hspace{.5in} (\mword{type}_{n-1}, \mword{disp}_{n-1} + (r \times \mword{darg} + 1)\times ex), \\
& & \ldots \\
& &  (\mword{type}_0, \mword{disp}_0 + ((r+1) \times \mword{darg} -1) \times ex), \ldots , \\
& & \hspace{.5in} (\mword{type}_{n-1}, \mword{disp}_{n-1} + ((r+1)
                        \times \mword{darg} - 1)\times ex), \\
& & \\
& &  (\mword{type}_0, \mword{disp}_0 + r \times \mword{darg} \times ex + \mword{psize} \times \mword{darg} \times ex),
        \ldots , \\
& & \hspace{.5in} (\mword{type}_{n-1}, \mword{disp}_{n-1} + r \times \mword{darg} \times ex
                + \mword{psize} \times \mword{darg} \times ex), \\
& &  (\mword{type}_0, \mword{disp}_0 + (r \times \mword{darg} + 1) \times ex
         + \mword{psize} \times \mword{darg} \times ex),
        \ldots , \\
& & \hspace{.5in} (\mword{type}_{n-1}, \mword{disp}_{n-1} + (r \times \mword{darg} + 1) \times ex
                + \mword{psize} \times \mword{darg} \times ex), \\
& & \ldots \\
& &  (\mword{type}_0, \mword{disp}_0 + ((r+1) \times \mword{darg} -1) \times ex
                 + \mword{psize} \times \mword{darg} \times ex), \ldots , \\
& & \hspace{.5in} (\mword{type}_{n-1}, \mword{disp}_{n-1} + ((r+1) \times \mword{darg} - 1)\times ex
                 + \mword{psize} \times \mword{darg} \times ex), \\
& & \hspace{.5in}\vdots \\
& &  (\mword{type}_0, \mword{disp}_0 + r \times \mword{darg} \times ex + \mword{psize} \times \mword{darg} \times ex
                \times (count - 1)),
        \ldots , \\
& & \hspace{.5in} (\mword{type}_{n-1}, \mword{disp}_{n-1} + r \times \mword{darg} \times ex
                + \mword{psize} \times \mword{darg} \times ex \times (count - 1)), \\
& &  (\mword{type}_0, \mword{disp}_0 + (r \times \mword{darg} + 1) \times ex
         + \mword{psize} \times \mword{darg} \times ex \times (count - 1)),
        \ldots , \\
& & \hspace{.5in} (\mword{type}_{n-1}, \mword{disp}_{n-1} + (r \times \mword{darg} + 1) \times ex \\
& & \hspace{1in} + \mword{psize} \times \mword{darg} \times ex \times (count - 1)), \\
& & \ldots \\
& &  (\mword{type}_0, \mword{disp}_0 + (r \times \mword{darg} + \mword{darg}_{\msword{last}}-1) \times ex \\
& & \hspace{1in} + \mword{psize} \times \mword{darg} \times ex \times (count-1)),
        \ldots , \\
& & \hspace{.5in} (\mword{type}_{n-1}, \mword{disp}_{n-1} + (r \times \mword{darg} + \mword{darg}_{\msword{last}} - 1)
                \times ex \\
& & \hspace{1in} + \mword{psize} \times \mword{darg} \times ex \times (count - 1)), \\
& &   (\mpiublb{ub\_marker}, gsize * ex) \}
\end{eqnarray*}
where $count$ is defined by this code fragment:
%%HEADER
%%FRAGMENT
%%DECL: int nblocks, gsize, darg, count, psize, left_over, r;
%%ENDHEADER
\begin{lstlisting}[language=C]
nblocks = (gsize + (darg - 1)) / darg;
count = nblocks / psize;
left_over = nblocks - count * psize;
if (r < left_over)
    count = count + 1;
\end{lstlisting}
Here, \code{nblocks} is the number of blocks that must be
distributed among the processors.
Finally, $\mword{darg}_{\msword{last}}$ is defined by this code fragment:
%%HEADER
%%FRAGMENT
%%DECL: int r, num_in_last_cyclic, gsize, psize, darg, darg_last;
%%ENDHEADER
\begin{lstlisting}[language=C]
if ((num_in_last_cyclic = gsize % (psize * darg)) == 0)
    darg_last = darg;
else {
    darg_last = num_in_last_cyclic - darg * r;
    if (darg_last > darg)
        darg_last = darg;
    if (darg_last <= 0)
        darg_last = darg;
}
\end{lstlisting}

\begin{example}
Consider generating the filetypes corresponding to the HPF distribution:
\label{ex:io-hpf}
\exindex{Typemap}%
\Exindex{Fortran}{Datatypes for distributed arrays}{MPI\_Type\_create\_darray}%
%%HEADER
%%SKIP
%%ENDHEADER
\begin{lstlisting}[language={[HPF]Fortran}]
      <oldtype> FILEARRAY(100, 200, 300)
!HPF$ PROCESSORS PROCESSES(2, 3)
!HPF$ DISTRIBUTE FILEARRAY(CYCLIC(10), *, BLOCK) ONTO PROCESSES
\end{lstlisting}
This can be achieved by the following Fortran code,
assuming there will be six processes attached to the run:
%%HEADER
%%LANG: FORTRAN
%%FRAGMENT
%%DECL: integer array_of_gsizes(3), array_of_distribs(3), array_of_dargs(3)
%%DECL: integer ierr, size, rank, oldtype, newtype, ndims
%%DECL: integer array_of_psizes(3)
%%ENDHEADER
\begin{lstlisting}[language={[MPI]Fortran}]
ndims = 3
array_of_gsizes(1) = 100
array_of_distribs(1) = MPI_DISTRIBUTE_CYCLIC
array_of_dargs(1) = 10
array_of_gsizes(2) = 200
array_of_distribs(2) = MPI_DISTRIBUTE_NONE
array_of_dargs(2) = 0
array_of_gsizes(3) = 300
array_of_distribs(3) = MPI_DISTRIBUTE_BLOCK
array_of_dargs(3) = MPI_DISTRIBUTE_DFLT_DARG
array_of_psizes(1) = 2
array_of_psizes(2) = 1
array_of_psizes(3) = 3
call MPI_COMM_SIZE(MPI_COMM_WORLD, size, ierr)
call MPI_COMM_RANK(MPI_COMM_WORLD, rank, ierr)
call MPI_TYPE_CREATE_DARRAY(size, rank, ndims, array_of_gsizes, &
     array_of_distribs, array_of_dargs, array_of_psizes,        &
     MPI_ORDER_FORTRAN, oldtype, newtype, ierr)
\end{lstlisting}
\end{example}


\subsection{Address and Size Procedures}
\label{subsec:pt2pt-addfunc}

The displacements in a general datatype are relative to some initial buffer
address.
\mpitermdefni{Absolute addresses}\mpitermdefindex{absolute addresses}\mpitermdefindex{addresses!absolute}
can be substituted for these
displacements: we treat them as displacements relative to ``address
zero,'' the start of the address space.  This initial address zero is
indicated by the constant \mpiconstmain{MPI\_BOTTOM}.  Thus, a datatype can
specify the absolute address of the entries in the communication
buffer, in which case the \mpiarg{buf} argument is passed the value
\mpiconst{MPI\_BOTTOM}.
Note that in Fortran \mpiconst{MPI\_BOTTOM} is not usable for initialization or
assignment, see Section~\ref{subsec:named-constants}.

The address of a location in memory can be found by invoking the
procedure\flushline
\mpifunc{MPI\_GET\_ADDRESS}.
The \mpitermdef{relative displacement}\mpitermdefindex{addresses!relative displacement}
between two absolute addresses
can be calculated with the procedure \mpifunc{MPI\_AINT\_DIFF}. A new absolute
address as sum of an absolute base address and a relative displacement can be
calculated with the procedure \mpifunc{MPI\_AINT\_ADD}. To ensure portability,
arithmetic on absolute addresses should not be performed with the intrinsic
operators ``-'' and ``+''. See also Sections~\ref{subsec:displacement}
and~\ref{subsec:pt2pt-segmented} on pages \pageref{subsec:displacement} and
\pageref{subsec:pt2pt-segmented}.

\begin{rationale}
Address sized integer values, i.e., \mpictype{MPI\_Aint} or
\vlgb\mpiftypekind{ADDRESS} values, are signed integers, while
absolute addresses are unsigned quantities. Direct arithmetic on addresses
stored in address sized signed variables can cause overflows, resulting in
undefined behavior.
\end{rationale}

\cdeclindex{MPI\_Aint}%
\begin{funcdef}{MPI\_GET\_ADDRESS}{(\mbox{location}, \mbox{address})}
\funcarg{\IN}{location}{location in caller memory (choice)}
\funcarg{\OUT}{address}{address of location (integer)}
\end{funcdef}
\mpibindmain{MPI\_Get\_address}{(\gb{}const~void~*location, MPI\_Aint~*address)}
\mpifnewbindmain{MPI\_Get\_address}{(location, address, ierror) \fargs TYPE(*), DIMENSION(..), ASYNCHRONOUS\ ::\ \mbox{location}\fnarg{}INTEGER(KIND=MPI\_ADDRESS\_KIND), INTENT(OUT)\ ::\ \mbox{address}\fnfinalarg{}INTEGER, OPTIONAL, INTENT(OUT)\ ::\ \mbox{ierror}}
\mpifbindmain{MPI\_GET\_ADDRESS}{(LOCATION, ADDRESS, IERROR) \fargs <type> \mbox{LOCATION(*)}\fnarg{}INTEGER(KIND=MPI\_ADDRESS\_KIND) \mbox{ADDRESS}\fnfinalarg{}INTEGER \mbox{IERROR}}

Returns the (byte) address of \mpiarg{location}.

\begin{rationale}
In the \code{mpi\_f08} module, the \mpiarg{location} argument is not defined
with \ftype{INTENT(IN)} because existing applications may use
\mpifunc{MPI\_GET\_ADDRESS} as a
substitute for \mpifunc{MPI\_F\_SYNC\_REG}, which was not defined before \MPIIIIDOTO/.
\end{rationale}

\begin{example}
\label{pt2pt-exW}
\Exindex{Fortran}{Get\_address}{MPI\_Get\_address}%
Using \mpifunc{MPI\_GET\_ADDRESS} for an array.

%%HEADER
%%FRAGMENT
%%LANG: FORTRAN
%%DECL: INTEGER IERROR
%%ENDHEADER
\begin{lstlisting}[language={[MPI]Fortran}]
REAL A(100,100)
INTEGER(KIND=MPI_ADDRESS_KIND) I1, I2, DIFF
CALL MPI_GET_ADDRESS(A(1,1), I1, IERROR)
CALL MPI_GET_ADDRESS(A(10,10), I2, IERROR)
DIFF = MPI_AINT_DIFF(I2, I1)
! The value of DIFF is 909*SIZEOF(REAL); the values of I1 and I2 are
! implementation dependent.
\end{lstlisting}
\end{example}

\begin{users}
C users may be tempted to avoid the usage of
\mpifunc{MPI\_GET\_ADDRESS}
and rely on
the availability of the address operator \code{\&}.  Note, however, that
\code{\&}~\emph{cast-expression} is a pointer, not an
address.
ISO C
does not require that the value of a pointer
(or the pointer cast to \ctype{int}) be
the absolute address of the object pointed at---although this is
commonly the case.
Furthermore, referencing may not have a unique
definition on machines with a segmented address space.
The use of
\mpifunc{MPI\_GET\_ADDRESS}
to ``reference'' C
variables guarantees portability to such machines as well.
\end{users}

%% See notes in pt2pt on why the commented out text is both wrong and
%% unnecessary
\begin{users}
  To prevent problems with the argument copying and register optimization done
  by Fortran compilers, please note the hints in
Sections~\ref{sec:misc-problems}--\ref{sec:f90-problems:comparison-with-C}.
%% In particular, refer to
%% %
%% Sections~\ref{sec:misc-sequence} and~\ref{sec:f90-problems:vector-subscripts}
%% on pages~\pageref{sec:misc-sequence}-\pageref{sec:f90-problems:vector-subscripts}
%% about ``{\sf Problems Due to Data Copying and Sequence Association with Subscript Triplets}''
%% and ``{\sf Vector Subscripts}'',
%% and Sections~\ref{sec:misc-register}-\ref{sec:f90-problems:perm-data-movements}
%% on pages~\pageref{sec:misc-register}-\pageref{sec:f90-problems:perm-data-movements}
%% about ``{\sf Optimization Problems}'', ``{\sf Code Movements and Register Optimization}'',
%% ``{\sf Temporary Data Movements}'' and ``{\sf Permanent Data Movements}''.
\end{users}

To ensure portability, arithmetic on \MPI/ addresses must be
performed using the \mpifunc{MPI\_AINT\_ADD}
and \mpifunc{MPI\_AINT\_DIFF} procedures.

\cdeclindex{MPI\_Aint}%
\begin{funcdef}{MPI\_AINT\_ADD}{(\mbox{base}, \mbox{disp})}
\funcarg{\IN}{base}{base address (integer)}
\funcarg{\IN}{disp}{displacement (integer)}
\end{funcdef}
\mpibindnotintmain{MPI\_Aint}{MPI\_Aint\_add}{(\gb{}MPI\_Aint~base, MPI\_Aint~disp)}
\mpifnewbindretmain{INTEGER(KIND=MPI\_ADDRESS\_KIND)}{MPI\_Aint\_add}{(base, disp) \fargs \fnfinalarg{}INTEGER(KIND=MPI\_ADDRESS\_KIND), INTENT(IN)\ ::\ \mbox{base}, \mbox{disp}}
\mpifbindretmain{INTEGER(KIND=MPI\_ADDRESS\_KIND)}{MPI\_AINT\_ADD}{(BASE, DISP) \fargs \fnfinalarg{}INTEGER(KIND=MPI\_ADDRESS\_KIND) \mbox{BASE}, \mbox{DISP}}

\mpifunc{MPI\_AINT\_ADD} produces a new \mpictype{MPI\_Aint} value that is
equivalent to the sum of the \mpiarg{base} and \mpiarg{disp} arguments, where
\mpiarg{base} represents a base address returned by a call to
\mpifunc{MPI\_GET\_ADDRESS} and \mpiarg{disp} represents a signed integer
displacement. The resulting address is valid only at the process that generated
\mpiarg{base}, and it must correspond to a location in the same object
referenced by \mpiarg{base}, as described in
Section~\ref{subsec:pt2pt-segmented}. The addition is performed in a manner
that results in the correct \mpictype{MPI\_Aint} representation of the output
address, as if the process that originally produced \mpiarg{base} had called:
%%HEADER
%%FRAGMENT
%%LANG: C
%%DECL: MPI_Aint disp, result; char *base;
%%ENDHEADER
\begin{lstlisting}[language={[MPI]C}]
MPI_Get_address((char *) base + disp, &result);
\end{lstlisting}

\cdeclindex{MPI\_Aint}%
\begin{funcdef}{MPI\_AINT\_DIFF}{(\mbox{addr1}, \mbox{addr2})}
\funcarg{\IN}{addr1}{minuend address (integer)}
\funcarg{\IN}{addr2}{subtrahend address (integer)}
\end{funcdef}
\mpibindnotintmain{MPI\_Aint}{MPI\_Aint\_diff}{(\gb{}MPI\_Aint~addr1, MPI\_Aint~addr2)}
\mpifnewbindretmain{INTEGER(KIND=MPI\_ADDRESS\_KIND)}{MPI\_Aint\_diff}{(addr1, addr2) \fargs \fnfinalarg{}INTEGER(KIND=MPI\_ADDRESS\_KIND), INTENT(IN)\ ::\ \mbox{addr1}, \mbox{addr2}}
\mpifbindretmain{INTEGER(KIND=MPI\_ADDRESS\_KIND)}{MPI\_AINT\_DIFF}{(ADDR1, ADDR2) \fargs \fnfinalarg{}INTEGER(KIND=MPI\_ADDRESS\_KIND) \mbox{ADDR1}, \mbox{ADDR2}}

\mpifunc{MPI\_AINT\_DIFF} produces a new \mpictype{MPI\_Aint} value that is
equivalent to the difference between \mpiarg{addr1} and \mpiarg{addr2}
arguments, where \mpiarg{addr1} and \mpiarg{addr2} represent addresses returned
by calls to \mpifunc{MPI\_GET\_ADDRESS}.  The resulting address is valid only
at the process that generated \mpiarg{addr1} and \mpiarg{addr2}, and
\mpiarg{addr1} and \mpiarg{addr2} must correspond to locations in the same
object in the same process, as described in
Section~\ref{subsec:pt2pt-segmented}. The difference is calculated in a manner
that results in
the signed difference from \mpiarg{addr1} to \mpiarg{addr2}, as if
the process that originally produced the addresses had called
\code{(char *) addr1 - (char *) addr2}
on the addresses initially passed to \mpifunc{MPI\_GET\_ADDRESS}.

The following auxiliary procedures provide useful information on
derived datatypes.

\begin{funcdef}{MPI\_TYPE\_SIZE}{(\mbox{datatype}, \mbox{size})}
\funcarg{\IN}{datatype}{datatype to get information on (handle)}
\funcarg{\OUT}{size}{datatype size (integer)}
\end{funcdef}
\mpibindmain{MPI\_Type\_size}{(\gb{}MPI\_Datatype~datatype, int~*size)}
\MPIbindindex{MPI\_Type\_size\_c}
\mpibindmain{MPI\_Type\_size\_c}{(\gb{}MPI\_Datatype~datatype, MPI\_Count~*size)}
\mpifnewbindmain{MPI\_Type\_size}{(datatype, size, ierror) \fargs TYPE(MPI\_Datatype), INTENT(IN)\ ::\ \mbox{datatype}\fnarg{}INTEGER, INTENT(OUT)\ ::\ \mbox{size}\fnfinalarg{}INTEGER, OPTIONAL, INTENT(OUT)\ ::\ \mbox{ierror}}
\mpifnewbindmain{MPI\_Type\_size}{(datatype, size, ierror) !(\_c) \fargs TYPE(MPI\_Datatype), INTENT(IN)\ ::\ \mbox{datatype}\fnarg{}INTEGER(KIND=MPI\_COUNT\_KIND), INTENT(OUT)\ ::\ \mbox{size}\fnfinalarg{}INTEGER, OPTIONAL, INTENT(OUT)\ ::\ \mbox{ierror}}
\mpifbindmain{MPI\_TYPE\_SIZE}{(DATATYPE, SIZE, IERROR) \fargs \fnfinalarg{}INTEGER \mbox{DATATYPE}, \mbox{SIZE}, \mbox{IERROR}}

\mpifunc{MPI\_TYPE\_SIZE}
set the value of \mpiarg{size} to the total size, in bytes, of the entries in
the type signature
associated with \mpiarg{datatype}; i.e., the total size of
the data in a message that would be created with this datatype.  Entries that
occur multiple times in the datatype are counted with their multiplicity.
%
For both procedures, if the \OUT\ parameter cannot express the value to
be returned (e.g., if the parameter is too small to hold the output
value), it is set to \mpiconst{MPI\_UNDEFINED}.

\subsection{Lower-Bound and Upper-Bound Markers}
\mpitermtitleindex{lower-bound markers}
\mpitermtitleindex{upper-bound markers}
\label{subsec:pt2pt-markers}

It is often convenient to define explicitly the lower bound and upper
bound of a type map, and override the definition given
on page~\pageref{eq:pt2pt-extent}.
This allows one to define a datatype that
has ``holes'' at its beginning or its end, or a datatype with
entries that extend above the upper bound or below the lower bound.
Examples of
such usage are provided in Section~\ref{subsec:pt2pt-examples}.
Also, the user may want to override the alignment rules that are
used to compute upper bounds and extents.  E.g., a C compiler may allow
the user to override default alignment rules for some of the
structures within
a program.  The user has to specify explicitly the bounds of the datatypes
that match these structures.

To achieve this, we add two additional conceptual datatypes,
\mpitermdef{lb\_marker} and\flushline
\mpitermdef{ub\_marker}, that represent the lower
bound and upper
bound of a datatype.  These conceptual datatypes occupy no space
($extent(\mpiublb{lb\_marker}) = extent(\mpiublb{ub\_marker}) =0$).
They do not
affect the size or count of a datatype, and do not affect the
 content of a message created with this datatype.  However, they do
affect the definition of the
extent of a datatype and, therefore, affect the outcome of a replication of
this datatype by a datatype constructor.

\begin{example}
\label{pt2pt-exX}
\Exindex{Neutral}{Typemap of nested datatypes}{MPI\_TYPE\_CREATE\_STRUCT,MPI\_TYPE\_CONTIGUOUS}%
A call to
\mpifunc{MPI\_TYPE\_CREATE\_RESIZED}\mpicode{(MPI\_INT, -3, 9, type1)}
creates a new datatype that has an
extent of 9 (from -3 to 5, 5 included), and contains an integer at
displacement 0.   This is the datatype defined by the typemap
\{(\mpiublb{lb\_marker}, -3), (int, 0), (\mpiublb{ub\_marker}, 6)\}.
If this type is replicated twice by a call to
\mpifunc{MPI\_TYPE\_CONTIGUOUS}\mpicode{(2, type1, type2)} then the newly created
type can
be described by the typemap
\{(\mpiublb{lb\_marker}, -3), (int, 0), (int,9), (\mpiublb{ub\_marker}, 15)\}.
(An entry of type
\mpiublb{ub\_marker}
can be deleted if there is another entry of type \mpiublb{ub\_marker} with a
higher
displacement; an entry of type \mpiublb{lb\_marker} can be deleted if there
is another
entry of type \mpiublb{lb\_marker} with a lower displacement.)
\end{example}

In general, if
\begin{displaymath}
Typemap = \{ (\mword{type}_0 , \mword{disp}_0 ) , \ldots , (\mword{type}_{n-1} , \mword{disp}_{n-1}) \} ,
\end{displaymath}
then the \mpitermdef{lower bound} of $Typemap$ is defined to be
\[
lb(Typemap) = \left\{ \begin{array}{ll}
\min_j \mword{disp}_j & \parbox{1.5in}{\raggedright if no entry has type
\mpiublb{lb\_marker}} \\
\min_j \{ \mword{disp}_j \ \mbox{such that}\ \mword{type}_j = \mpiublb{lb\_marker} \} & \mbox{otherwise}
\end{array}
\right. \]
Similarly,
the \mpitermdef{upper bound} of $Typemap$ is defined to be
\[
ub(Typemap) = \left\{ \begin{array}{ll}
\max_j(\mword{disp}_j + \mpicode{sizeof}(\mword{type}_j)) + \epsilon & \parbox{1.4in}{\raggedright if no entry has type
\mpiublb{ub\_marker}}
\\ \max_j \{ \mword{disp}_j \ \mbox{such that}\ \mword{type}_j = \mpiublb{ub\_marker} \} & \mbox{otherwise}
\end{array}
\right. \]
Then
\pagelabel{eq:pt2pt-extent}
\[
extent(Typemap) = ub(Typemap) - lb(Typemap)
\]
If $\mword{type}_i$ requires alignment to a byte address that is a multiple of $k_i$,
then $\epsilon$ is the least nonnegative increment needed to round
$extent(Typemap)$ to the next multiple of $\max_i k_i$.
In Fortran, it is implementation dependent whether the \MPI/ implementation computes the alignments $k_i$ according to the alignments
used by the compiler in common blocks, \ftype{SEQUENCE} derived types, \ftype{BIND(C)} derived types,
or derived types that are neither \ftype{SEQUENCE} nor \ftype{BIND(C)}.

The formal definitions given for the various datatype constructors
apply now, with the amended definition of \mpitermdefni{extent}\mpitermdefindex{extent of datatypes}.

\begin{rationale}
Before Fortran 2003, \mpifunc{MPI\_TYPE\_CREATE\_STRUCT} could be applied to Fortran common blocks and
\ftype{SEQUENCE} derived types. With Fortran 2003, this list was extended by \ftype{BIND(C)} derived types
and \MPI/ implementors have implemented the alignments $k_i$ differently, i.e., some based on
the alignments used in \ftype{SEQUENCE} derived types, and others according to \ftype{BIND(C)} derived types.
\end{rationale}

\begin{implementors}
In Fortran, it is generally recommended to use \ftype{BIND(C)} derived types instead of
common blocks or \ftype{SEQUENCE} derived types. Therefore it is recommended to calculate the
alignments $k_i$ based on \ftype{BIND(C)} derived types.
\end{implementors}

\begin{users}
Structures combining different basic datatypes should be defined so that there will
be no gaps based on alignment rules.  If such a datatype is used to create an array of structures, users
should also avoid an alignment-gap at the end of the structure.
In \MPI/ communication, the content of such gaps would not be communicated
into the receiver's buffer.
For example, such an alignment-gap may occur between an odd number of
\ctype{float}s or \ftype{REAL}s before a \ctype{double} or
\ftype{DOUBLE PRECISION} data.  Such gaps
may be added explicitly to both the structure and the \MPI/ derived
datatype handle because the communication of a contiguous derived
datatype may be significantly faster than the communication of one
that is noncontiguous because of such alignment-gaps.

As an example, instead of

%%HEADER
%%FRAGMENT
%%LANG: FORTRAN08
%%ENDHEADER
\begin{lstlisting}[language={[MPI08]Fortran}]
TYPE, BIND(C) :: my_data
  REAL, DIMENSION(3) :: x
  ! there may be a gap of the size of one REAL
  ! if the alignment of a DOUBLE PRECISION is
  ! two times the size of a REAL
  DOUBLE PRECISION :: p
END TYPE
\end{lstlisting}

one should define

%%HEADER
%%FRAGMENT
%%LANG: FORTRAN08
%%ENDHEADER
\begin{lstlisting}[language={[MPI08]Fortran}]
TYPE, BIND(C) :: my_data
  REAL, DIMENSION(3) :: x
  REAL :: gap1
  DOUBLE PRECISION :: p
END TYPE
\end{lstlisting}

and also include \mpiarg{gap1} in the matching \MPI/ derived datatype.
It is required that all processes in a communication add the same
gaps, i.e., defined with the same basic datatype.
Both the original and the modified structures are portable,
but may have different performance implications for the
communication and memory accesses during computation on
systems with different alignment values.

In principle, a compiler may define an additional alignment rule
for structures, e.g., to use at least 4 or 8 byte alignment, although
the content may have a $\max_i k_i$ alignment less than this structure
alignment. To maintain portability, users should
always resize structure derived datatype handles
if used in an array of structures, see the Example in
\sectionref{sec:f90-problems:derived-types}.
\end{users}

\subsection{Extent and Bounds of Datatypes}
\mpitermtitleindex{extent of datatypes}
\mpitermtitleindex{bounds of datatypes}
\label{subsec:pt2pt-extent}

\cdeclindex{MPI\_Aint}%
\begin{funcdef}{MPI\_TYPE\_GET\_EXTENT}{(\mbox{datatype}, \mbox{lb}, \mbox{extent})}
\funcarg{\IN}{datatype}{datatype to get information on (handle)}
\funcarg{\OUT}{lb}{lower bound of datatype (integer)}
\funcarg{\OUT}{extent}{extent of datatype (integer)}
\end{funcdef}
\mpibindmain{MPI\_Type\_get\_extent}{(\gb{}MPI\_Datatype~datatype, MPI\_Aint~*lb, MPI\_Aint~*extent)}
\MPIbindindex{MPI\_Type\_get\_extent\_c}
\mpibindmain{MPI\_Type\_get\_extent\_c}{(\gb{}MPI\_Datatype~datatype, MPI\_Count~*lb, MPI\_Count~*extent)}
\mpifnewbindmain{MPI\_Type\_get\_extent}{(datatype, lb, extent, ierror) \fargs TYPE(MPI\_Datatype), INTENT(IN)\ ::\ \mbox{datatype}\fnarg{}INTEGER(KIND=MPI\_ADDRESS\_KIND), INTENT(OUT)\ ::\ \mbox{lb}, \mbox{extent}\fnfinalarg{}INTEGER, OPTIONAL, INTENT(OUT)\ ::\ \mbox{ierror}}
\mpifnewbindmain{MPI\_Type\_get\_extent}{(datatype, lb, extent, ierror) !(\_c) \fargs TYPE(MPI\_Datatype), INTENT(IN)\ ::\ \mbox{datatype}\fnarg{}INTEGER(KIND=MPI\_COUNT\_KIND), INTENT(OUT)\ ::\ \mbox{lb}, \mbox{extent}\fnfinalarg{}INTEGER, OPTIONAL, INTENT(OUT)\ ::\ \mbox{ierror}}
\mpifbindmain{MPI\_TYPE\_GET\_EXTENT}{(DATATYPE, LB, EXTENT, IERROR) \fargs INTEGER \mbox{DATATYPE}, \mbox{IERROR}\fnfinalarg{}INTEGER(KIND=MPI\_ADDRESS\_KIND) \mbox{LB}, \mbox{EXTENT}}

Returns the lower bound and the extent of
\mpiarg{datatype}
(as defined in \namedref{Equation}{soft-lb-ub-definition}).

If either \OUT\ parameter cannot express the value
to be returned (e.g., if the parameter is too small to hold the output
value), it is set to \mpiconst{MPI\_UNDEFINED}.


\MPI/ allows one to change the extent of a datatype, using lower bound
and upper bound markers.  This provides control over the stride of successive
datatypes that are replicated by datatype constructors, or are
replicated by the \mpiarg{count} argument in a send or receive
call.

\cdeclindex{MPI\_Aint}%
\begin{funcdef}{MPI\_TYPE\_CREATE\_RESIZED}{(\mbox{oldtype}, \mbox{lb}, \mbox{extent}, \mbox{newtype})}
\funcarg{\IN}{oldtype}{input datatype (handle)}
\funcarg{\IN}{lb}{new lower bound of datatype (integer)}
\funcarg{\IN}{extent}{new extent of datatype (integer)}
\funcarg{\OUT}{newtype}{output datatype (handle)}
\end{funcdef}
\mpibindmain{MPI\_Type\_create\_resized}{(\gb{}MPI\_Datatype~oldtype, MPI\_Aint~lb, MPI\_Aint~extent, MPI\_Datatype~*newtype)}
\MPIbindindex{MPI\_Type\_create\_resized\_c}
\mpibindmain{MPI\_Type\_create\_resized\_c}{(\gb{}MPI\_Datatype~oldtype, MPI\_Count~lb, MPI\_Count~extent, MPI\_Datatype~*newtype)}
\mpifnewbindmain{MPI\_Type\_create\_resized}{(oldtype, lb, extent, newtype, ierror) \fargs TYPE(MPI\_Datatype), INTENT(IN)\ ::\ \mbox{oldtype}\fnarg{}INTEGER(KIND=MPI\_ADDRESS\_KIND), INTENT(IN)\ ::\ \mbox{lb}, \mbox{extent}\fnarg{}TYPE(MPI\_Datatype), INTENT(OUT)\ ::\ \mbox{newtype}\fnfinalarg{}INTEGER, OPTIONAL, INTENT(OUT)\ ::\ \mbox{ierror}}
\mpifnewbindmain{MPI\_Type\_create\_resized}{(oldtype, lb, extent, newtype, ierror) !(\_c) \fargs TYPE(MPI\_Datatype), INTENT(IN)\ ::\ \mbox{oldtype}\fnarg{}INTEGER(KIND=MPI\_COUNT\_KIND), INTENT(IN)\ ::\ \mbox{lb}, \mbox{extent}\fnarg{}TYPE(MPI\_Datatype), INTENT(OUT)\ ::\ \mbox{newtype}\fnfinalarg{}INTEGER, OPTIONAL, INTENT(OUT)\ ::\ \mbox{ierror}}
\mpifbindmain{MPI\_TYPE\_CREATE\_RESIZED}{(OLDTYPE, LB, EXTENT, NEWTYPE, IERROR) \fargs INTEGER \mbox{OLDTYPE}, \mbox{NEWTYPE}, \mbox{IERROR}\fnfinalarg{}INTEGER(KIND=MPI\_ADDRESS\_KIND) \mbox{LB}, \mbox{EXTENT}}

Returns in \mpiarg{newtype} a handle to a new datatype that is
identical to \mpiarg{oldtype}, except that the lower bound of this new
datatype is set to be \mpiarg{lb}, and its upper bound is set to be
\mpiarg{lb}\mpicode{ + }\mpiarg{extent}.
Any previous \mpitermdefni{lb}\mpitermdefindex{lb\_marker!erased}
and \mpitermdefni{ub}\mpitermdefindex{ub\_marker!erased} markers are erased,
and a new pair of lower bound and upper bound markers are put in the
positions indicated by the \mpiarg{lb} and \mpiarg{extent} arguments.
This affects the behavior of the datatype when used in communication
operations, with  \mpiarg{count} $>1$,  and when used in the
construction of new derived datatypes.

\subsection{True Extent of Datatypes}
\mpitermtitleindex{true extent of datatypes}
\mpitermtitleindex{extent of datatypes!true extent}
\label{subsec:pt2pt-true-extent}

Suppose we implement gather
(see also \sectionref{sec:coll-gather})
as a spanning tree implemented on top of
point-to-point routines.  Since the receive buffer is only valid on the
root process, one will need to allocate some temporary space for
receiving data on intermediate nodes.  However, the datatype extent
cannot be used as an estimate of the amount of space that needs to be
allocated, if the user has modified the extent, for example by
using \mpifunc{MPI\_TYPE\_CREATE\_RESIZED}.
The procedure \mpifunc{MPI\_TYPE\_GET\_TRUE\_EXTENT} returns
the true extent of the datatype.

\cdeclindex{MPI\_Aint}%
\begin{funcdef}{MPI\_TYPE\_GET\_TRUE\_EXTENT}{(\mbox{datatype}, \mbox{true\_lb}, \mbox{true\_extent})}
\funcarg{\IN}{datatype}{datatype to get information on (handle)}
\funcarg{\OUT}{true\_lb}{true lower bound of datatype (integer)}
\funcarg{\OUT}{true\_extent}{true extent of datatype (integer)}
\end{funcdef}
\mpibindmain{MPI\_Type\_get\_true\_extent}{(\gb{}MPI\_Datatype~datatype, MPI\_Aint~*true\_lb, MPI\_Aint~*true\_extent)}
\MPIbindindex{MPI\_Type\_get\_true\_extent\_c}
\mpibindmain{MPI\_Type\_get\_true\_extent\_c}{(\gb{}MPI\_Datatype~datatype, MPI\_Count~*true\_lb, MPI\_Count~*true\_extent)}
\mpifnewbindmain{MPI\_Type\_get\_true\_extent}{(datatype, true\_lb, true\_extent, ierror) \fargs TYPE(MPI\_Datatype), INTENT(IN)\ ::\ \mbox{datatype}\fnarg{}INTEGER(KIND=MPI\_ADDRESS\_KIND), INTENT(OUT)\ ::\ \mbox{true\_lb}, \mbox{true\_extent}\fnfinalarg{}INTEGER, OPTIONAL, INTENT(OUT)\ ::\ \mbox{ierror}}
\mpifnewbindmain{MPI\_Type\_get\_true\_extent}{(datatype, true\_lb, true\_extent, ierror) !(\_c) \fargs TYPE(MPI\_Datatype), INTENT(IN)\ ::\ \mbox{datatype}\fnarg{}INTEGER(KIND=MPI\_COUNT\_KIND), INTENT(OUT)\ ::\ \mbox{true\_lb}, \mbox{true\_extent}\fnfinalarg{}INTEGER, OPTIONAL, INTENT(OUT)\ ::\ \mbox{ierror}}
\mpifbindmain{MPI\_TYPE\_GET\_TRUE\_EXTENT}{(DATATYPE, TRUE\_LB, TRUE\_EXTENT, IERROR) \fargs INTEGER \mbox{DATATYPE}, \mbox{IERROR}\fnfinalarg{}INTEGER(KIND=MPI\_ADDRESS\_KIND) \mbox{TRUE\_LB}, \mbox{TRUE\_EXTENT}}

\mpiarg{true\_lb} returns the offset of the lowest unit of storage that
is addressed by the datatype, i.e., the lower bound of the
corresponding typemap, ignoring explicit lower bound markers.
\mpiarg{true\_extent} returns the
true size of
the datatype, i.e., the extent of the corresponding typemap, ignoring
explicit lower bound and upper bound markers, and performing no
rounding for alignment.  If the typemap associated with
\mpiarg{datatype} is
\[
\mword{Typemap} = \{ (\mword{type}_0, \mword{disp}_0), \ldots , (\mword{type}_{n-1}, \mword{disp}_{n-1})\}
\]
Then
\[
true\_lb(\mword{Typemap}) = \min_j  \{ \mword{disp}_j ~:~ \mword{type}_j \ne \mpiublb{lb\_marker}, \mpiublb{ub\_marker} \},
\]
\[
true\_ub (\mword{Typemap}) = \max_j \{\mword{disp}_j + \mpicode{sizeof}(\mword{type}_j) ~:~ \mword{type}_j \ne
\mpiublb{lb\_marker}, \mpiublb{ub\_marker}\} ,
\]
and
\[
true\_extent (\mword{Typemap}) = true\_ub(\mword{Typemap}) - true\_lb(\mword{Typemap}).
\]
(Readers should compare this with the definitions in
\sectionref{subsec:pt2pt-markers} and
\sectionref{subsec:pt2pt-extent},
which describe the procedure
\mpifunc{MPI\_TYPE\_GET\_EXTENT}.)

The \mpiarg{true\_extent} is the minimum number of bytes of
memory necessary to hold a datatype, uncompressed.

If either \OUT\ parameter cannot express the value to
be returned (e.g., if the parameter is too small to hold the output
value), it is set to \mpiconst{MPI\_UNDEFINED}.

\subsection{Commit and Free}
\label{subsec:pt2pt-comfree}

A datatype object has to be \mpitermdefni{committed}\mpitermdefindex{commit} before it can be used in a
communication.
As an argument in datatype constructors, uncommitted and also
committed datatypes can be used.
There is no need to commit basic datatypes. They are ``pre-committed.''

\begin{funcdef}{MPI\_TYPE\_COMMIT}{(\mbox{datatype})}
\funcarg{\INOUT}{datatype}{datatype that is committed (handle)}
\end{funcdef}
\mpibindmain{MPI\_Type\_commit}{(\gb{}MPI\_Datatype~*datatype)}
\mpifnewbindmain{MPI\_Type\_commit}{(datatype, ierror) \fargs TYPE(MPI\_Datatype), INTENT(INOUT)\ ::\ \mbox{datatype}\fnfinalarg{}INTEGER, OPTIONAL, INTENT(OUT)\ ::\ \mbox{ierror}}
\mpifbindmain{MPI\_TYPE\_COMMIT}{(DATATYPE, IERROR) \fargs \fnfinalarg{}INTEGER \mbox{DATATYPE}, \mbox{IERROR}}

The commit operation commits the datatype, that is, the formal description of
a communication buffer, not the content of that buffer.  Thus, after a datatype
has been committed, it can be repeatedly reused to communicate the changing
content
of a buffer or, indeed, the content of different buffers, with different
starting addresses.

\begin{implementors}
The system may ``compile'' at commit time an internal
representation for the datatype that facilitates
communication, e.g.,
change from a compacted representation to a flat representation of the
datatype, and select the most convenient transfer mechanism.

The optimizations chosen during \mpifunc{MPI\_TYPE\_COMMIT} may no longer be optimal
if a session (or the \wpm{}) is initialized or finalized.
\end{implementors}

\mpifunc{MPI\_TYPE\_COMMIT}
will accept a committed datatype;
in this case, it is equivalent to a no-op.


\begin{example}
\label{datatype-exY}
\Exindex{Fortran}{Type\_commit}{MPI\_Type\_commit}%
The following code fragment gives examples of using
\mpifunc{MPI\_TYPE\_COMMIT}.
%%HEADER
%%FRAGMENT
%%LANG: FORTRAN
%%DECL: INTEGER ierr
%%ENDHEADER
\begin{lstlisting}[language={[MPI]Fortran}]
INTEGER type1, type2
CALL MPI_TYPE_CONTIGUOUS(5, MPI_REAL, type1, ierr)
              ! new type object created
CALL MPI_TYPE_COMMIT(type1, ierr)
              ! now type1 can be used for communication
type2 = type1
              ! type2 can be used for communication
              ! (it is a handle to same object as type1)
CALL MPI_TYPE_VECTOR(3, 5, 4, MPI_REAL, type1, ierr)
              ! new uncommitted type object created
CALL MPI_TYPE_COMMIT(type1, ierr)
              ! now type1 can be used anew for communication
\end{lstlisting}
\end{example}

\begin{funcdef}{MPI\_TYPE\_FREE}{(\mbox{datatype})}
\funcarg{\INOUT}{datatype}{datatype that is freed (handle)}
\end{funcdef}
\mpibindmain{MPI\_Type\_free}{(\gb{}MPI\_Datatype~*datatype)}
\mpifnewbindmain{MPI\_Type\_free}{(datatype, ierror) \fargs TYPE(MPI\_Datatype), INTENT(INOUT)\ ::\ \mbox{datatype}\fnfinalarg{}INTEGER, OPTIONAL, INTENT(OUT)\ ::\ \mbox{ierror}}
\mpifbindmain{MPI\_TYPE\_FREE}{(DATATYPE, IERROR) \fargs \fnfinalarg{}INTEGER \mbox{DATATYPE}, \mbox{IERROR}}

Marks the datatype object associated with \mpiarg{datatype} for
deallocation and sets \mpiarg{datatype} to \mpiconstmain{MPI\_DATATYPE\_NULL}.
Any communication that is currently using this datatype will complete
normally.
%
Freeing a datatype does not affect any other datatype that was
built from the freed datatype.  The system behaves as if
input datatype arguments to derived datatype constructors are passed by value.

\begin{implementors}
The implementation may keep a reference count of active communications
that use the datatype, in order to decide when to free it.
Also, one may implement constructors of derived datatypes so that they keep
pointers to their datatype arguments, rather than copying them.  In this
case, one needs to keep track of active datatype definition references in order
to know when a datatype object can be freed.
\end{implementors}

%%%%%%%%%%%%%%%%%%%%%%%%%%%%%%%%%%%%%%%%%%%%%%%%%%%%%%%%%%%%%%
\subsection{Duplicating a Datatype}
\label{sec:ei-dupdata}

\begin{funcdef}{MPI\_TYPE\_DUP}{(\mbox{oldtype}, \mbox{newtype})}
\funcarg{\IN}{oldtype}{datatype (handle)}
\funcarg{\OUT}{newtype}{copy of \mpiarg{oldtype} (handle)}
\end{funcdef}
\mpibindmain{MPI\_Type\_dup}{(\gb{}MPI\_Datatype~oldtype, MPI\_Datatype~*newtype)}
\mpifnewbindmain{MPI\_Type\_dup}{(oldtype, newtype, ierror) \fargs TYPE(MPI\_Datatype), INTENT(IN)\ ::\ \mbox{oldtype}\fnarg{}TYPE(MPI\_Datatype), INTENT(OUT)\ ::\ \mbox{newtype}\fnfinalarg{}INTEGER, OPTIONAL, INTENT(OUT)\ ::\ \mbox{ierror}}
\mpifbindmain{MPI\_TYPE\_DUP}{(OLDTYPE, NEWTYPE, IERROR) \fargs \fnfinalarg{}INTEGER \mbox{OLDTYPE}, \mbox{NEWTYPE}, \mbox{IERROR}}

\mpifunc{MPI\_TYPE\_DUP} is
a type constructor
that duplicates the existing \mpishortarg{oldtype}
with associated key values.
For each key value, the respective copy
callback function determines the attribute value associated with this
key in the new datatype; one particular action that a copy
callback may take is to delete the attribute from the new datatype.
Returns in \mpiarg{newtype} a new datatype with exactly the same
properties as \mpiarg{oldtype} and any copied cached
information, see \sectionref{subsec:caching:datatypes}.
The new datatype has identical
upper bound and lower bound and yields the same net result when fully
decoded with the procedures in Section~\ref{subsec:ei-dataaccess}.
The
\mpiarg{newtype} has the same committed state as the old
\mpiarg{oldtype}.

\subsection{Use of General Datatypes in Communication}
\label{subsec:pt2pt-datatypeuse}

Handles to derived datatypes can be passed to a communication call wherever a
datatype argument is required.
A call of the form \mpifunc{MPI\_SEND}\mpicode{(buf, count, datatype, ...)}, where
$\mpicode{count} > 1$, is interpreted as if the call was passed a new datatype
that is the
concatenation of \mpiarg{count} copies of \mpiarg{datatype}.
Thus,
\mpifunc{MPI\_SEND}\mpicode{(buf, count, datatype, dest, tag, comm)} is equivalent to,
%%HEADER
%%SKIP
%%ENDHEADER
\begin{lstlisting}[language={[MPI]Fortran}]
MPI_TYPE_CONTIGUOUS(count, datatype, newtype)
MPI_TYPE_COMMIT(newtype)
MPI_SEND(buf, 1, newtype, dest, tag, comm)
MPI_TYPE_FREE(newtype).
\end{lstlisting}
Similar statements apply to all other communication procedures that have a
\mpiarg{count} and \mpiarg{datatype} argument.

Suppose that a send operation \mpifunc{MPI\_SEND}\mpicode{(buf, count,
datatype, dest, tag, comm)} is executed, where
\mpiarg{datatype} has type map,
\[
\{(\mword{type}_0, \mword{disp}_0),\ldots,(\mword{type}_{n-1}, \mword{disp}_{n-1})\},
\]
and extent $extent$.  (Explicit lower bound and upper bound markers are not listed in the type map, but
they affect the value of $extent$.)
The send
operation sends $n \cdot \mpicode{count}$ entries, where entry $i
\cdot n + j$ is at location
$addr_{i,j} = \mpicode{buf} + extent \cdot i + \mword{disp}_j$
and has type $\mword{type}_j$,
for $i = 0 ,\ldots, \mpicode{count}-1$ and $j = 0 ,\ldots, n-1$.
These entries need not be contiguous, nor distinct; their order can be
arbitrary.

The variable stored at address $addr_{i,j}$ in the calling program
should be of a type that matches $\mword{type}_j$, where
type matching is defined as in Section~\ref{subsec:pt2pt-typematch}.
The message sent contains $n \cdot \mpicode{count}$ entries, where entry
$i \cdot n +j$ has type $\mword{type}_j$.

Similarly, suppose that a receive operation
\mpifunc{MPI\_RECV}\mpicode{(buf, count, datatype, source, tag, comm, status)} is
executed, where \mpiarg{datatype} has type map,
\begin{displaymath}
\{(\mword{type}_0, \mword{disp}_0) ,\ldots,(\mword{type}_{n-1}, \mword{disp}_{n-1}) \},
\end{displaymath}
\noindent
with extent $extent$.  (Again,  explicit lower bound and upper bound markers are not listed in the type map, but
they affect the value of $extent$.)
This receive operation receives
$n \cdot \mpicode{count}$ entries, where entry $i \cdot n + j$ is at location
$\mpicode{buf} + extent \cdot i + \mword{disp}_j$
and has type $\mword{type}_j$.  If the incoming message consists of $k$
elements, then we must have $k \le n \cdot \mpicode{count}$; the $i \cdot n +
j$-th element of the message should have a type that matches $\mword{type}_j$.

\mpitermdefni{Type matching}\mpitermdefindex{type matching}\mpitermdefindex{matching!type}
is defined according to the type signature of
the corresponding datatypes, that is, the sequence of basic type
components.  Type matching does not depend on some aspects of the
datatype definition, such as the displacements (layout in memory) or the
intermediate types used.

\begin{example}
\label{pt2pt-exAA}
\Exindex{Fortran}{Matching type with datatypes}{MPI\_Type\_contiguous,MPI\_Send,MPI\_Recv}%
\exindex{Datatype!matching type}%
This example shows that type matching is defined in terms of
the basic types that a derived type consists of.
%%HEADER
%%SKIP
%%ENDHEADER
\begin{lstlisting}[language={[MPI]Fortran}]
...
CALL MPI_TYPE_CONTIGUOUS(2, MPI_REAL, type2, ...)
CALL MPI_TYPE_CONTIGUOUS(4, MPI_REAL, type4, ...)
CALL MPI_TYPE_CONTIGUOUS(2, type2, type22, ...)
...
CALL MPI_SEND(a, 4, MPI_REAL, ...)
CALL MPI_SEND(a, 2, type2, ...)
CALL MPI_SEND(a, 1, type22, ...)
CALL MPI_SEND(a, 1, type4, ...)
...
CALL MPI_RECV(a, 4, MPI_REAL, ...)
CALL MPI_RECV(a, 2, type2, ...)
CALL MPI_RECV(a, 1, type22, ...)
CALL MPI_RECV(a, 1, type4, ...)
\end{lstlisting}
Each of the sends matches any of the receives.
\end{example}

A datatype may specify overlapping entries. The use of such a
datatype in any communication in association with a buffer updated by
the operation is erroneous. (This is erroneous even
if the actual message received is short enough not to write any entry
more than once.)

Suppose that
\mpifunc{MPI\_RECV}\mpicode{(buf, count, datatype, dest, tag, comm, status)} is
executed, where \mpiarg{datatype} has type map,
\begin{displaymath}
\{(\mword{type}_0, \mword{disp}_0) ,\ldots,(\mword{type}_{n-1}, \mword{disp}_{n-1}) \}.
\end{displaymath}
The received message need not fill all the receive buffer, nor does it
need to fill a number of locations that is a multiple of $n$.
Any number, $k$, of basic
elements can be received, where $0 \le k \le \mpicode{count} \cdot n$.
The number of basic elements received can be retrieved from
\mpiarg{status} using the query procedure \mpifunc{MPI\_GET\_ELEMENTS}.

\cdeclindex{MPI\_Status}%
\begin{funcdef}{MPI\_GET\_ELEMENTS}{(\mbox{status}, \mbox{datatype}, \mbox{count})}
\funcarg{\IN}{status}{return status of receive operation (status)}
\funcarg{\IN}{datatype}{datatype used by receive operation (handle)}
\funcarg{\OUT}{count}{number of received basic elements (integer)}
\end{funcdef}
\mpibindmain{MPI\_Get\_elements}{(\gb{}const~MPI\_Status~*status, MPI\_Datatype~datatype, int~*count)}
\MPIbindindex{MPI\_Get\_elements\_c}
\mpibindmain{MPI\_Get\_elements\_c}{(\gb{}const~MPI\_Status~*status, MPI\_Datatype~datatype, MPI\_Count~*count)}
\mpifnewbindmain{MPI\_Get\_elements}{(status, datatype, count, ierror) \fargs TYPE(MPI\_Status), INTENT(IN)\ ::\ \mbox{status}\fnarg{}TYPE(MPI\_Datatype), INTENT(IN)\ ::\ \mbox{datatype}\fnarg{}INTEGER, INTENT(OUT)\ ::\ \mbox{count}\fnfinalarg{}INTEGER, OPTIONAL, INTENT(OUT)\ ::\ \mbox{ierror}}
\mpifnewbindmain{MPI\_Get\_elements}{(status, datatype, count, ierror) !(\_c) \fargs TYPE(MPI\_Status), INTENT(IN)\ ::\ \mbox{status}\fnarg{}TYPE(MPI\_Datatype), INTENT(IN)\ ::\ \mbox{datatype}\fnarg{}INTEGER(KIND=MPI\_COUNT\_KIND), INTENT(OUT)\ ::\ \mbox{count}\fnfinalarg{}INTEGER, OPTIONAL, INTENT(OUT)\ ::\ \mbox{ierror}}
\mpifbindmain{MPI\_GET\_ELEMENTS}{(STATUS, DATATYPE, COUNT, IERROR) \fargs \fnfinalarg{}INTEGER \mbox{STATUS(MPI\_STATUS\_SIZE)}, \mbox{DATATYPE}, \mbox{COUNT}, \mbox{IERROR}}

The \mpiarg{datatype} argument should match the argument provided by the
receive call that set the \mpiarg{status} variable.
For both procedures, if the \OUT\ parameter cannot express the value to
be returned (e.g., if the parameter is too small to hold the output
value), it is set to \mpiconst{MPI\_UNDEFINED}.

The previously defined procedure
\mpifunc{MPI\_GET\_COUNT}
(Section~\ref{subsec:pt2pt-status}),
has a different behavior.
It returns the number of ``top-level
entries'' received, i.e., the number of ``copies'' of type
\mpiarg{datatype}.
In the previous example, \mpifunc{MPI\_GET\_COUNT}
may return any integer value $k$, where $0 \le k \le \mpicode{count}$.
If \mpifunc{MPI\_GET\_COUNT} returns $k$, then the number of basic
elements received (and the value returned by
\mpifunc{MPI\_GET\_ELEMENTS})
is $n \cdot k$.   If the number of basic elements received is not a
multiple of $n$, that is, if the receive operation has not received an
integral number of \mpiarg{datatype} ``copies,'' then
\mpifunc{MPI\_GET\_COUNT} sets the value of \mpiarg{count} to \mpiconst{MPI\_UNDEFINED}.

\begin{example}
\label{pt2pt-exBB}
\Exindex{Fortran}{Using Get\_count and Get\_elements}{MPI\_Get\_count,MPI\_Get\_elements,MPI\_Type\_contiguous}%
Usage of \mpifunc{MPI\_GET\_COUNT} and
\mpifunc{MPI\_GET\_ELEMENTS}.
%%HEADER
%%LANG: FORTRAN
%%FRAGMENT
%%DECL: integer Type2, ierr, comm, stat(MPI_STATUS_SIZE), i, rank
%%DECL: real a(10)
%%SKIPELIPSIS
%%ENDHEADER
\begin{lstlisting}[language={[MPI]Fortran},basicstyle=\lstsmall]
...
CALL MPI_TYPE_CONTIGUOUS(2, MPI_REAL, Type2, ierr)
CALL MPI_TYPE_COMMIT(Type2, ierr)
...
CALL MPI_COMM_RANK(comm, rank, ierr)
IF (rank .EQ. 0) THEN
   CALL MPI_SEND(a, 2, MPI_REAL, 1, 0, comm, ierr)
   CALL MPI_SEND(a, 3, MPI_REAL, 1, 0, comm, ierr)
ELSE IF (rank .EQ. 1) THEN
   CALL MPI_RECV(a, 2, Type2, 0, 0, comm, stat, ierr)
   CALL MPI_GET_COUNT(stat, Type2, i, ierr)     ! returns i=1
   CALL MPI_GET_ELEMENTS(stat, Type2, i, ierr)  ! returns i=2
   CALL MPI_RECV(a, 2, Type2, 0, 0, comm, stat, ierr)
   CALL MPI_GET_COUNT(stat, Type2, i, ierr)     ! returns i=MPI_UNDEFINED
   CALL MPI_GET_ELEMENTS(stat, Type2, i, ierr)  ! returns i=3
END IF
\end{lstlisting}
\end{example}

The procedure \mpifunc{MPI\_GET\_ELEMENTS}
can also be used after a probe
to find the number of elements in the probed message.
Note that the
\mpifunc{MPI\_GET\_COUNT} and
\mpifunc{MPI\_GET\_ELEMENTS}
return the same values when they are used
with basic datatypes as long as the limits of their respective \mpiarg{count} arguments are not exceeded.


\begin{rationale}
The extension given to the definition of \mpifunc{MPI\_GET\_COUNT} seems
natural: one would expect this procedure to return the value of the
\mpiarg{count} argument, when the receive buffer is filled.
Sometimes \mpiarg{datatype} represents
a basic unit of data one wants to transfer,
for example, a record in an array of records (structures).
One should be able to find out how many components were received
without bothering to divide by the number of elements in each
component.   However, on other occasions, \mpiarg{datatype} is used to
define a complex layout of data in the receiver memory, and does not represent
a basic unit of data for transfers.  In such cases, one needs to use
the procedure \mpifunc{MPI\_GET\_ELEMENTS}.
\end{rationale}

\begin{implementors}
The definition implies that a receive cannot change the value of
storage outside the entries defined to compose the communication
buffer.
In particular, the definition implies that padding space in a structure
should not be modified when such a structure is copied from one process to
another. This would
prevent the obvious optimization of copying the structure, together
with the padding, as one contiguous block.
The implementation is free to do this optimization when it does not
impact the outcome of the computation.
The user can ``force'' this optimization by explicitly including
padding as part of the message.
\end{implementors}

\subsection{Correct Use of Addresses}
\mpitermtitleindex{addresses!correct use}
\label{subsec:pt2pt-segmented}

Successively declared variables in C or Fortran are not necessarily
stored at contiguous locations.  Thus, care must be exercised
that displacements do not cross from one variable
to another.  Also, in machines with a segmented address space,
addresses are not unique and address arithmetic has some peculiar
properties.   Thus, the use of \mpitermdef{addresses},
that is, displacements relative to the
start address \mpiconst{MPI\_BOTTOM}, has to be restricted.

Variables belong
to the same \mpitermdef{sequential storage} if they belong to the same
array,
to the same \code{COMMON} block in Fortran, or to the same structure in C.
Valid addresses are defined recursively as follows:

\begin{enumerate}
\item
The procedure
\mpifunc{MPI\_GET\_ADDRESS}
returns a valid address, when
passed as argument a variable of the calling program.
\item
The \mpiarg{buf} argument of a communication procedure evaluates to a
valid address, when passed as argument a variable of the calling program.
\item
If \mpiarg{v} is a valid address, and \mpiarg{i} is an
integer, then \mpiarg{v+i} is a valid address, provided \mpiarg{v} and
\mpiarg{v+i} are in the same sequential storage.
\end{enumerate}

A correct program uses only valid addresses to identify the
locations of entries in communication buffers.
Furthermore, if \mpiarg{u} and \mpiarg{v} are two valid addresses, then
the (integer)
difference $\mpiarg{u}-\mpiarg{v}$ can be computed only if both \mpiarg{u} and \mpiarg{v}
are in the same sequential storage.
No other arithmetic operations can be meaningfully executed on addresses.

The rules above impose no constraints on the use of derived
datatypes, as long as they are used to define a communication buffer
that is wholly contained within the same sequential storage.
However, the construction of a communication buffer that contains
variables that are not within the same sequential storage must obey
certain restrictions.  Basically, a communication buffer with
variables that are not within the same sequential storage can be used
only by specifying in the communication call \mpiarg{buf}\mpicode{ = }%
\mpiconst{MPI\_BOTTOM}, \mpiarg{count}\mpicode{ = 1}, and using a \mpiarg{datatype}
argument where all displacements are valid (absolute) addresses.

\begin{users}
It is not expected that \MPI/ implementations will be able to detect
erroneous, ``out of bound'' displacements---unless those overflow the
user address space---since the \MPI/ call may not know the extent of the
arrays and records in the host program.
\end{users}



\begin{implementors}
There is no need to distinguish (absolute) addresses and (relative)
displacements on a machine with contiguous address space:
\mpiconst{MPI\_BOTTOM} is
zero, and both addresses and displacements are integers.  On machines where the
distinction is required, addresses are recognized as expressions that involve
\mpiconst{MPI\_BOTTOM}.
\end{implementors}

% Note that in Fortran, Fortran INTEGERs may be too small to contain an
% address (e.g., 32 bit INTEGERs on a machine with 64bit pointers).
% Because of this, in Fortran, implementations may restrict the use of
% absolute addresses to only part of the process memory, and restrict the
% use of relative displacements to subranges of the process memory where they
% are constrained by the size of Fortran INTEGERs.

%%%%%%%%%%%%%%%%%%%%%%%%%%%%%%%%%%%%%%%%%%%%%%%%%%%%%%%%%%%%%%%%%%%%%%
\subsection{Decoding a Datatype}
\label{subsec:typeaccess}
\label{subsec:ei-dataaccess}

\MPI/
datatype objects
allow users to specify
an arbitrary layout of data in memory.
There are several cases
where accessing the layout information in
opaque datatype objects would be useful.
The opaque datatype
object has found a number of uses outside \mpi/.  Furthermore, a
number of tools wish to display internal information about a datatype.
To achieve this, datatype decoding procedures are provided.
The two procedures in this section are used together to decode datatypes
to recreate the calling sequence used in their initial definition.
These can be used to allow a user to determine the type map and type
signature of a datatype.

\begin{funcdef}{MPI\_TYPE\_GET\_ENVELOPE}{(\mbox{datatype}, \mbox{num\_integers}, \mbox{num\_addresses}, \mbox{num\_large\_counts}, \mbox{num\_datatypes}, \mbox{combiner})}
\funcarg{\IN}{datatype}{datatype to decode (handle)}
\funcarg{\OUT}{num\_integers}{number of input integers used in call constructing \mpiarg{combiner} (nonnegative integer)}
\funcarg{\OUT}{num\_addresses}{number of input addresses used in call constructing \mpiarg{combiner} (nonnegative integer)}
\funcarg{\OUT}{num\_large\_counts}{number of input large counts used in call constructing \mpiarg{combiner} (nonnegative integer, \textbf{only present for large count variants})}
\funcarg{\OUT}{num\_datatypes}{number of input datatypes used in call constructing \mpiarg{combiner} (nonnegative integer)}
\funcarg{\OUT}{combiner}{combiner (state)}
\end{funcdef}
\mpibindmain{MPI\_Type\_get\_envelope}{(\gb{}MPI\_Datatype~datatype, int~*num\_integers, int~*num\_addresses, int~*num\_datatypes, int~*combiner)}
\MPIbindindex{MPI\_Type\_get\_envelope\_c}
\mpibindmain{MPI\_Type\_get\_envelope\_c}{(\gb{}MPI\_Datatype~datatype, MPI\_Count~*num\_integers, MPI\_Count~*num\_addresses, MPI\_Count~*num\_large\_counts, MPI\_Count~*num\_datatypes, int~*combiner)}
\mpifnewbindmain{MPI\_Type\_get\_envelope}{(datatype, num\_integers, num\_addresses, num\_datatypes, combiner, ierror) \fargs TYPE(MPI\_Datatype), INTENT(IN)\ ::\ \mbox{datatype}\fnarg{}INTEGER, INTENT(OUT)\ ::\ \mbox{num\_integers}, \mbox{num\_addresses}, \mbox{num\_datatypes}, \mbox{combiner}\fnfinalarg{}INTEGER, OPTIONAL, INTENT(OUT)\ ::\ \mbox{ierror}}
\mpifnewbindmain{MPI\_Type\_get\_envelope}{(datatype, num\_integers, num\_addresses, num\_large\_counts, num\_datatypes, combiner, ierror) !(\_c) \fargs TYPE(MPI\_Datatype), INTENT(IN)\ ::\ \mbox{datatype}\fnarg{}INTEGER(KIND=MPI\_COUNT\_KIND), INTENT(OUT)\ ::\ \mbox{num\_integers}, \mbox{num\_addresses}, \mbox{num\_large\_counts}, \mbox{num\_datatypes}\fnarg{}INTEGER, INTENT(OUT)\ ::\ \mbox{combiner}\fnfinalarg{}INTEGER, OPTIONAL, INTENT(OUT)\ ::\ \mbox{ierror}}
\mpifbindmain{MPI\_TYPE\_GET\_ENVELOPE}{(DATATYPE, NUM\_INTEGERS, NUM\_ADDRESSES, NUM\_DATATYPES, COMBINER, IERROR) \fargs \fnfinalarg{}INTEGER \mbox{DATATYPE}, \mbox{NUM\_INTEGERS}, \mbox{NUM\_ADDRESSES}, \mbox{NUM\_DATATYPES}, \mbox{COMBINER}, \mbox{IERROR}}

For the given \mpiarg{datatype},
\mpifunc{MPI\_TYPE\_GET\_ENVELOPE} returns information on the
number and type of input arguments used in the call that created the
\mpiarg{datatype}.  The number-of-arguments values returned can be
used to provide sufficiently large arrays in the decoding routine
\mpifunc{MPI\_TYPE\_GET\_CONTENTS}.  This call and the meaning of the
returned values is described below.  The \mpiarg{combiner} reflects
the \mpi/ datatype constructor call that was used in creating
\mpiarg{datatype}.

\begin{rationale}
By requiring that the \mpiarg{combiner} reflect the constructor used
in the creation
of the \mpiarg{datatype}, the decoded information can be used
to effectively recreate the calling sequence used in the original
creation.
This is the most useful information and
was
felt to be reasonable even though it constrains implementations to
remember the original constructor sequence even if the internal
representation is different.

The decoded information keeps track of datatype duplications.  This is
important as one needs to distinguish between a predefined datatype
and a dup of a predefined datatype.  The former is a constant object
that cannot be freed, while the latter is a derived datatype that can
be freed.
\end{rationale}

The list of values that can be returned from \mpiskipfunc{MPI\_TYPE\_GET\_ENVELOPE} in
\mpiarg{combiner} (on the left) and the call associated with them (on the right) are as follows:
\begin{constlist}[180pt]
\constitem{MPI\_COMBINER\_NAMED}{a named predefined datatype}
\constitem{MPI\_COMBINER\_DUP}{\mpifunc{MPI\_TYPE\_DUP}}
\constitem{MPI\_COMBINER\_CONTIGUOUS}{\mpifunc{MPI\_TYPE\_CONTIGUOUS}}
\constitem{MPI\_COMBINER\_VECTOR}{\mpifunc{MPI\_TYPE\_VECTOR}}
\constitem{MPI\_COMBINER\_HVECTOR}{\mpifunc{MPI\_TYPE\_CREATE\_HVECTOR}}
\constitem{MPI\_COMBINER\_INDEXED}{\mpifunc{MPI\_TYPE\_INDEXED}}
\constitem{MPI\_COMBINER\_HINDEXED}{\mpifunc{MPI\_TYPE\_CREATE\_HINDEXED}}
\constitem{MPI\_COMBINER\_INDEXED\_BLOCK}{\mpifunc{MPI\_TYPE\_CREATE\_INDEXED\_BLOCK}}
\constitem{MPI\_COMBINER\_HINDEXED\_BLOCK}{\mpifunc{MPI\_TYPE\_CREATE\_HINDEXED\_BLOCK}}
\constitem{MPI\_COMBINER\_STRUCT}{\mpifunc{MPI\_TYPE\_CREATE\_STRUCT}}
\constitem{MPI\_COMBINER\_SUBARRAY}{\mpifunc{MPI\_TYPE\_CREATE\_SUBARRAY}}
\constitem{MPI\_COMBINER\_DARRAY}{\mpifunc{MPI\_TYPE\_CREATE\_DARRAY}}
\constitem{MPI\_COMBINER\_F90\_REAL}{\mpifunc{MPI\_TYPE\_CREATE\_F90\_REAL}}
\constitem{MPI\_COMBINER\_F90\_COMPLEX}{\mpifunc{MPI\_TYPE\_CREATE\_F90\_COMPLEX}}
\constitem{MPI\_COMBINER\_F90\_INTEGER}{\mpifunc{MPI\_TYPE\_CREATE\_F90\_INTEGER}}
\constitem{MPI\_COMBINER\_RESIZED}{\mpifunc{MPI\_TYPE\_CREATE\_RESIZED}}
\constitem{MPI\_COMBINER\_VALUE\_INDEX}{\mpifunc{MPI\_TYPE\_GET\_VALUE\_INDEX}}
\end{constlist}

If \mpiarg{combiner} is \mpiconst{MPI\_COMBINER\_NAMED} then \mpiarg{datatype} is a
named predefined datatype.

If the \mpifunc{MPI\_TYPE\_GET\_ENVELOPE} variant without
\mpiarg{num\_large\_counts} is invoked with a \mpiarg{datatype} that
requires an output value of \mpiarg{num\_large\_counts} $>0$,
then an error of class \errormain{MPI\_ERR\_TYPE} is raised.

\begin{rationale}
  The large count variant of this \MPI/ procedure was added in \MPIIV/.
  It contains a new \mpiarg{num\_large\_counts} parameter.  The
  other variant---the variant that existed before \MPIIV/---was
  not changed in order to preserve backwards compatibility.
\end{rationale}

The actual arguments used in the creation call for a \mpiarg{datatype}
can be obtained using \mpifunc{MPI\_TYPE\_GET\_CONTENTS}.

\mpifunc{MPI\_TYPE\_GET\_ENVELOPE} and \mpifunc{MPI\_TYPE\_GET\_CONTENTS} also support large \gb{}count types in
separate additional \MPI/ procedures in C (suffixed with the ``\code{\_c}'') and interface
polymorphism  in Fortran when using \code{USE mpi\_f08}.

\cdeclindex{MPI\_Aint}%
\begin{funcdef}{MPI\_TYPE\_GET\_CONTENTS}{(\mbox{datatype}, \mbox{max\_integers}, \mbox{max\_addresses}, \mbox{max\_large\_counts}, \mbox{max\_datatypes}, \mbox{array\_of\_integers}, \mbox{array\_of\_addresses}, \mbox{array\_of\_large\_counts}, \mbox{array\_of\_datatypes})}
\funcarg{\IN}{datatype}{datatype to decode (handle)}
\funcarg{\IN}{max\_integers}{number of elements in \mpiarg{array\_of\_integers} (nonnegative integer)}
\funcarg{\IN}{max\_addresses}{number of elements in \mpiarg{array\_of\_addresses} (nonnegative integer)}
\funcarg{\IN}{max\_large\_counts}{number of elements in \mpiarg{array\_of\_large\_counts} (nonnegative integer, \textbf{only present for large count variants})}
\funcarg{\IN}{max\_datatypes}{number of elements in \mpiarg{array\_of\_datatypes} (nonnegative integer)}
\funcarg{\OUT}{array\_of\_integers}{contains integer arguments used in constructing \mpiarg{datatype} (array of integers)}
\funcarg{\OUT}{array\_of\_addresses}{contains address arguments used in constructing \mpiarg{datatype} (array of integers)}
\funcarg{\OUT}{array\_of\_large\_counts}{contains large count arguments used in constructing \mpiarg{datatype} (array of integers, \textbf{only present for large count variants})}
\funcarg{\OUT}{array\_of\_datatypes}{contains datatype arguments used in constructing \mpiarg{datatype} (array of handles)}
\end{funcdef}
\mpibindmain{MPI\_Type\_get\_contents}{(\gb{}MPI\_Datatype~datatype, int~max\_integers, int~max\_addresses, int~max\_datatypes, int~array\_of\_integers[], MPI\_Aint~array\_of\_addresses[], MPI\_Datatype~array\_of\_datatypes[])}
\MPIbindindex{MPI\_Type\_get\_contents\_c}
\mpibindmain{MPI\_Type\_get\_contents\_c}{(\gb{}MPI\_Datatype~datatype, MPI\_Count~max\_integers, MPI\_Count~max\_addresses, MPI\_Count~max\_large\_counts, MPI\_Count~max\_datatypes, int~array\_of\_integers[], MPI\_Aint~array\_of\_addresses[], MPI\_Count~array\_of\_large\_counts[], MPI\_Datatype~array\_of\_datatypes[])}
\mpifnewbindmain{MPI\_Type\_get\_contents}{(datatype, max\_integers, max\_addresses, max\_datatypes, array\_of\_integers, array\_of\_addresses, array\_of\_datatypes, ierror) \fargs TYPE(MPI\_Datatype), INTENT(IN)\ ::\ \mbox{datatype}\fnarg{}INTEGER, INTENT(IN)\ ::\ \mbox{max\_integers}, \mbox{max\_addresses}, \mbox{max\_datatypes}\fnarg{}INTEGER, INTENT(OUT)\ ::\ \mbox{array\_of\_integers(max\_integers)}\fnarg{}INTEGER(KIND=MPI\_ADDRESS\_KIND), INTENT(OUT)\ ::\ \mbox{array\_of\_addresses(max\_addresses)}\fnarg{}TYPE(MPI\_Datatype), INTENT(OUT)\ ::\ \mbox{array\_of\_datatypes(max\_datatypes)}\fnfinalarg{}INTEGER, OPTIONAL, INTENT(OUT)\ ::\ \mbox{ierror}}
\mpifnewbindmain{MPI\_Type\_get\_contents}{(datatype, max\_integers, max\_addresses, max\_large\_counts, max\_datatypes, array\_of\_integers, array\_of\_addresses, array\_of\_large\_counts, array\_of\_datatypes, ierror) !(\_c) \fargs TYPE(MPI\_Datatype), INTENT(IN)\ ::\ \mbox{datatype}\fnarg{}INTEGER(KIND=MPI\_COUNT\_KIND), INTENT(IN)\ ::\ \mbox{max\_integers}, \mbox{max\_addresses}, \mbox{max\_large\_counts}, \mbox{max\_datatypes}\fnarg{}INTEGER, INTENT(OUT)\ ::\ \mbox{array\_of\_integers(max\_integers)}\fnarg{}INTEGER(KIND=MPI\_ADDRESS\_KIND), INTENT(OUT)\ ::\ \mbox{array\_of\_addresses(max\_addresses)}\fnarg{}INTEGER(KIND=MPI\_COUNT\_KIND), INTENT(OUT)\ ::\ \mbox{array\_of\_large\_counts(max\_large\_counts)}\fnarg{}TYPE(MPI\_Datatype), INTENT(OUT)\ ::\ \mbox{array\_of\_datatypes(max\_datatypes)}\fnfinalarg{}INTEGER, OPTIONAL, INTENT(OUT)\ ::\ \mbox{ierror}}
\mpifbindmain{MPI\_TYPE\_GET\_CONTENTS}{(DATATYPE, MAX\_INTEGERS, MAX\_ADDRESSES, MAX\_DATATYPES, ARRAY\_OF\_INTEGERS, ARRAY\_OF\_ADDRESSES, ARRAY\_OF\_DATATYPES, IERROR) \fargs INTEGER \mbox{DATATYPE}, \mbox{MAX\_INTEGERS}, \mbox{MAX\_ADDRESSES}, \mbox{MAX\_DATATYPES}, \mbox{ARRAY\_OF\_INTEGERS(*)}, \mbox{ARRAY\_OF\_DATATYPES(*)}, \mbox{IERROR}\fnfinalarg{}INTEGER(KIND=MPI\_ADDRESS\_KIND) \mbox{ARRAY\_OF\_ADDRESSES(*)}}

\mpiarg{datatype} must be a predefined unnamed or a derived datatype;
the call is erroneous if \mpiarg{datatype} is a predefined named
datatype.

The values given for \mpiarg{max\_integers}, \mpiarg{max\_addresses},
\mpiarg{max\_large\_counts},
and \mpiarg{max\_datatypes}
must be at least as large as the value
returned in \mpiarg{num\_integers}, \mpiarg{num\_addresses},
\mpiarg{num\_large\_counts},
and \mpiarg{num\_datatypes},
respectively, in the call \mpifunc{MPI\_TYPE\_GET\_ENVELOPE}
for the same \mpiarg{datatype} argument.

\begin{rationale}
The arguments \mpiarg{max\_integers}, \mpiarg{max\_addresses},
\mpiarg{max\_large\_counts},
and
\mpiarg{max\_datatypes} allow for error checking in the
call.
\end{rationale}

If the \mpifunc{MPI\_TYPE\_GET\_CONTENTS} variant without
\mpiarg{max\_large\_counts} is invoked with a \mpiarg{datatype} that
requires $>0$ values in \mpiarg{array\_of\_large\_counts},
then an error of class \error{MPI\_ERR\_TYPE} is raised.

\begin{rationale}
  The large count variant of this \MPI/ procedure was added in \MPIIV/.
  It contains new \mpiarg{max\_large\_counts} and
  \mpiarg{array\_of\_large\_counts} parameters.  The
  other variant---the variant that existed before \MPIIV/---was
  not changed in order to preserve backwards compatibility.
\end{rationale}

The datatypes returned in \mpiarg{array\_of\_datatypes} are handles to
datatype objects that are equivalent to the datatypes used in the
original construction call.  If these were derived datatypes, then
the returned datatypes are new datatype objects, and the
user is responsible for freeing these datatypes with
\mpifunc{MPI\_TYPE\_FREE}.
If these were predefined datatypes, then
the returned datatype is equal to that (constant) predefined datatype
and cannot be freed.

The committed state of returned derived datatypes is undefined,
i.e., the datatypes may or
may not be committed.  Furthermore, the content of attributes of
returned datatypes is undefined.

Note that \mpifunc{MPI\_TYPE\_GET\_CONTENTS} can be invoked with a
\mpishortarg{datatype} argument that was constructed using
\mpifunc{MPI\_TYPE\_CREATE\_F90\_REAL},\flushline
\mpifunc{MPI\_TYPE\_CREATE\_F90\_INTEGER}, or
\mpifunc{MPI\_TYPE\_CREATE\_F90\_COMPLEX}  (an unnamed predefined datatype).
In such a case, an empty \mpiarg{array\_of\_datatypes} is returned.

\begin{rationale}
The definition of datatype equivalence implies that equivalent
predefined datatypes are equal.
By requiring the same handle for named predefined datatypes, it is
possible to use the \code{==} or \code{.EQ.} comparison operator to determine the
datatype involved.
\end{rationale}

\begin{implementors}
The datatypes returned in \mpiarg{array\_of\_datatypes} must appear to the
user as if each is an equivalent copy of the datatype used in the type
constructor call.
Whether this is done by
creating a new datatype or via another mechanism such as a reference
count mechanism is up to the implementation as long as the semantics
are preserved.
\end{implementors}

\begin{rationale}
The committed state and attributes of the returned datatype is
deliberately left vague.  The datatype used in the original
construction may have been modified since its use in the constructor
call.  Attributes can be added, removed, or modified as well as having
the datatype committed.  The semantics given allow for a
reference count implementation without having to track these changes.
\end{rationale}

In the
deprecated
 datatype constructor calls, the address arguments in Fortran are of
type \ftype{INTEGER}.  In the
preferred
calls, the address arguments are of
type \mpiftypekind{ADDRESS}.  The call
\mpifunc{MPI\_TYPE\_GET\_CONTENTS} returns all addresses in an argument
of type \mpiftypekind{ADDRESS}.  This is true even if the
deprecated
 calls were used.  Thus, the location of values returned can
be thought of as being returned by the C bindings.  It can also be
determined by examining the
preferred
 calls for datatype constructors
for the
deprecated
calls that involve addresses.

\begin{rationale}
By having all address arguments returned in the
\mpiarg{array\_of\_addresses} argument, the result from a C and Fortran
decoding of a \mpiarg{datatype} gives the result in the same
argument.  It is assumed that an integer
of type \mpiftypekind{ADDRESS} will be at least as large as
the \ftype{INTEGER} argument used in datatype construction with the old \mpii/
calls so no loss of information will occur.
\end{rationale}

The following defines what values are placed in each entry of the
returned arrays depending on the datatype constructor used for
\mpiarg{datatype}.  It also specifies the size of the arrays needed,
which is the values returned by \mpifunc{MPI\_TYPE\_GET\_ENVELOPE}.
In Fortran, the following calls were made:

% Note: in Fortran, ``TYPE'' is a keyword. While it can be used as a
% variable name, for clarity, a unique name, dtype (for DataTYPE) is
% used instead.
%%HEADER
%%LANG: FORTRAN
%%FRAGMENT
%%DECL: INTEGER LARGE
%%ENDHEADER
\begin{lstlisting}[language={[MPI]Fortran}]
PARAMETER (LARGE = 1000)
INTEGER DTYPE, NI, NA, ND, COMBINER, I(LARGE), D(LARGE), IERROR
INTEGER(KIND=MPI_ADDRESS_KIND) A(LARGE)
! CONSTRUCT DATATYPE DTYPE (NOT SHOWN)
CALL MPI_TYPE_GET_ENVELOPE(DTYPE, NI, NA, ND, COMBINER, IERROR)
IF ((NI .GT. LARGE) .OR. (NA .GT. LARGE) .OR. (ND .GT. LARGE)) THEN
   WRITE (*, *) "NI, NA, OR ND = ", NI, NA, ND, &
   " RETURNED BY MPI_TYPE_GET_ENVELOPE IS LARGER THAN LARGE = ", LARGE
   CALL MPI_ABORT(MPI_COMM_WORLD, 99, IERROR)
ENDIF
CALL MPI_TYPE_GET_CONTENTS(DTYPE, NI, NA, ND, I, A, D, IERROR)
\end{lstlisting}

\noindent
or in C the analogous calls of:

%%HEADER
%%LANG: C
%%FRAGMENT
%%ENDHEADER
\begin{lstlisting}[language={[MPI]C}]
#define LARGE 1000
int ni, na, nd, combiner, i[LARGE];
MPI_Aint a[LARGE];
MPI_Datatype dtype, d[LARGE];
/* construct datatype dtype (not shown) */
MPI_Type_get_envelope(dtype, &ni, &na, &nd, &combiner);
if ((ni > LARGE) || (na > LARGE) || (nd > LARGE)) {
   fprintf(stderr, "ni, na, or nd = %d %d %d returned by ", ni, na, nd);
   fprintf(stderr, "MPI_Type_get_envelope is larger than LARGE = %d\n",
           LARGE);
   MPI_Abort(MPI_COMM_WORLD, 99);
}
MPI_Type_get_contents(dtype, ni, na, nd, i, a, d);
\end{lstlisting}

\noindent
The following describes the values of the arguments for each combiner.
The lower case name of arguments is used.
Also, the descriptions below refer to \MPI/ datatypes created by
procedures without large count arguments.
\begin{description}[style=unboxed,leftmargin=0cm]

\item[\mpiconstmain{MPI\_COMBINER\_NAMED}]
the \mpiarg{datatype} represents a predefined type and therefore
it is erroneous to call \mpifuncshort{MPI\_TYPE\_GET\_CONTENTS}.

\item[\mpiconstmain{MPI\_COMBINER\_DUP}]
\code{ni = 0}, \code{na = 0}, \code{nd = 1}, and

\begin{center}
\begin{tabular}{l c c}
\hline
Constructor argument &  C &    Fortran location                \\
\hline
oldtype &               d[0] &          D(1)                            \\
\hline
\end{tabular}
\end{center}

\item[\mpiconstmain{MPI\_COMBINER\_CONTIGUOUS}]
\code{ni = 1}, \code{na = 0}, \code{nd = 1}, and

\begin{center}
\begin{tabular}{l c c}
\hline
Constructor argument &  C &    Fortran location                \\
\hline
count &                 i[0] &          I(1)                            \\
oldtype &               d[0] &          D(1)                            \\
\hline
\end{tabular}
\end{center}

\filbreak %%ALLOWLATEX%% % Intended to mark this as a good place to insert
% a page break
\item[\mpiconstmain{MPI\_COMBINER\_VECTOR}]
\code{ni = 3}, \code{na = 0}, \code{nd = 1}, and

\begin{center}
\begin{tabular}{l c c}
\hline
Constructor argument &  C &    Fortran location                \\
\hline
count &                 i[0] &          I(1)                            \\
blocklength &           i[1] &          I(2)                            \\
stride &                i[2] &          I(3)                            \\
oldtype &               d[0] &          D(1)                            \\
\hline
\end{tabular}
\end{center}

\item[\mpiconstmain{MPI\_COMBINER\_HVECTOR}]
\code{ni = 2}, \code{na = 1}, \code{nd = 1}, and

\begin{center}
\begin{tabular}{l c c}
\hline
Constructor argument &  C &    Fortran location                \\
\hline
count &                 i[0] &          I(1)                            \\
blocklength &           i[1] &          I(2)                            \\
stride &                a[0] &          A(1)                            \\
oldtype &               d[0] &          D(1)                            \\
\hline
\end{tabular}
\end{center}

\item[\mpiconstmain{MPI\_COMBINER\_INDEXED}]
\code{ni = 2*count+1}, \code{na = 0}, \code{nd = 1}, and

\begin{center}
\begin{tabular}{l c c}
\hline
Constructor argument &          C &                    Fortran location                \\
\hline
count                   &       i[0] &                          I(1)                            \\
array\_of\_blocklengths &       i[1] to i[i[0]] &               I(2) to I(I(1)+1)               \\
array\_of\_displacements &      i[i[0]+1] to i[2*i[0]] &        I(I(1)+2) to I(2*I(1)+1)        \\
oldtype &                       d[0] &                          D(1)                            \\
\hline
\end{tabular}
\end{center}

\item[\mpiconstmain{MPI\_COMBINER\_HINDEXED}]
\code{ni = count+1}, \code{na = count}, \code{nd = 1}, and

\begin{center}
\begin{tabular}{l c c}
\hline
Constructor argument &          C &                    Fortran location                \\
\hline
count                   &       i[0] &                          I(1)                            \\
array\_of\_blocklengths &       i[1] to i[i[0]] &               I(2) to I(I(1)+1)               \\
array\_of\_displacements &      a[0] to a[i[0]-1] &             A(1) to A(I(1))                 \\
oldtype &                       d[0] &                          D(1)                            \\
\hline
\end{tabular}
\end{center}

\item[\mpiconstmain{MPI\_COMBINER\_INDEXED\_BLOCK}]
\code{ni = count+2}, \code{na = 0}, \code{nd = 1}, and

\begin{center}
\begin{tabular}{l c c}
\hline
Constructor argument &          C &                    Fortran location                \\
\hline
count                   &       i[0] &                          I(1)                            \\
blocklength             &       i[1] &                          I(2)                            \\
array\_of\_displacements &      i[2] to i[i[0]+1] &             I(3) to I(I(1)+2)    \\
oldtype &                       d[0] &                          D(1)                            \\
\hline
\end{tabular}
\end{center}

\item[\mpiconstmain{MPI\_COMBINER\_HINDEXED\_BLOCK}]
\code{ni = 2}, \code{na = count}, \code{nd = 1}, and

\begin{center}
\begin{tabular}{l c c}
\hline
Constructor argument &          C &                    Fortran location                \\
\hline
count                   &       i[0] &                          I(1)                            \\
blocklength             &       i[1] &                          I(2)                            \\
array\_of\_displacements &      a[0] to a[i[0]-1] &             A(1) to A(I(1))                 \\
oldtype &                       d[0] &                          D(1)                            \\
\hline
\end{tabular}
\end{center}

\item[\mpiconstmain{MPI\_COMBINER\_STRUCT}]
\code{ni = count+1}, \code{na = count}, \code{nd = count}, and

\begin{center}
\begin{tabular}{l c c}
\hline
Constructor argument &          C &                    Fortran location                \\
\hline
count &                         i[0] &                          I(1)                            \\
array\_of\_blocklengths &       i[1] to i[i[0]] &               I(2) to I(I(1)+1)               \\
array\_of\_displacements &      a[0] to a[i[0]-1] &             A(1) to A(I(1))                 \\
array\_of\_types &              d[0] to d[i[0]-1] &             D(1) to D(I(1))                 \\
\hline
\end{tabular}
\end{center}

\item[\mpiconstmain{MPI\_COMBINER\_SUBARRAY}]
\code{ni = 3*ndims+2}, \code{na = 0}, \code{nd = 1}, and

\begin{center}
\begin{tabular}{l c c}
\hline
Constructor argument &          C &                    Fortran location                \\
\hline
ndims &                         i[0] &                          I(1)                            \\
array\_of\_sizes &              i[1] to i[i[0]] &               I(2) to I(I(1)+1)               \\
array\_of\_subsizes &           i[i[0]+1] to i[2*i[0]] &        I(I(1)+2) to I(2*I(1)+1)        \\
array\_of\_starts &             i[2*i[0]+1] to i[3*i[0]] &      I(2*I(1)+2) to I(3*I(1)+1)      \\
order &                         i[3*i[0]+1] &                   I(3*I(1)+2]                     \\
oldtype &                       d[0] &                          D(1)                            \\
\hline
\end{tabular}
\end{center}

\filbreak %%ALLOWLATEX%% % Intended to mark this as a good place to insert
% a page break
\item[\mpiconstmain{MPI\_COMBINER\_DARRAY}]
\code{ni = 4*ndims+4}, \code{na = 0}, \code{nd = 1}, and

\begin{center}
\begin{tabular}{l c c}
\hline
Constructor argument &          C &                    Fortran location                \\
\hline
size &                          i[0] &                          I(1)                            \\
rank &                          i[1] &                          I(2)                            \\
ndims &                         i[2] &                          I(3)                            \\
array\_of\_gsizes &             i[3] to i[i[2]+2] &             I(4) to I(I(3)+3)               \\
array\_of\_distribs &           i[i[2]+3] to i[2*i[2]+2] &      I(I(3)+4) to I(2*I(3)+3)        \\
array\_of\_dargs &              i[2*i[2]+3] to i[3*i[2]+2] &    I(2*I(3)+4) to I(3*I(3)+3)      \\
array\_of\_psizes &             i[3*i[2]+3] to i[4*i[2]+2] &    I(3*I(3)+4) to I(4*I(3)+3)      \\
order &                         i[4*i[2]+3] &                   I(4*I(3)+4)                     \\
oldtype &                       d[0] &                          D(1)                            \\
\hline
\end{tabular}
\end{center}

\item[\mpiconstmain{MPI\_COMBINER\_F90\_REAL}]
\code{ni = 2}, \code{na = 0}, \code{nd = 0}, and

\begin{center}
\begin{tabular}{l c c}
\hline
Constructor argument &          C &                    Fortran location                \\
\hline
p &                             i[0] &                          I(1)                            \\
r &                             i[1] &                          I(2)                            \\
\hline
\end{tabular}
\end{center}

\item[\mpiconstmain{MPI\_COMBINER\_F90\_COMPLEX}]
\code{ni = 2}, \code{na = 0}, \code{nd = 0}, and

\begin{center}
\begin{tabular}{l c c}
\hline
Constructor argument &          C &                    Fortran location                \\
\hline
p &                             i[0] &                          I(1)                            \\
r &                             i[1] &                          I(2)                            \\
\hline
\end{tabular}
\end{center}

\item[\mpiconstmain{MPI\_COMBINER\_F90\_INTEGER}]
\code{ni = 1}, \code{na = 0}, \code{nd = 0}, and

\begin{center}
\begin{tabular}{l c c}
\hline
Constructor argument &          C &                    Fortran location                \\
\hline
r &                             i[0] &                          I(1)                            \\
\hline
\end{tabular}
\end{center}

\item[\mpiconstmain{MPI\_COMBINER\_RESIZED}]
\code{ni = 0}, \code{na = 2}, \code{nd = 1}, and

\begin{center}
\begin{tabular}{l c c}
\hline
Constructor argument &          C &                    Fortran location                \\
\hline
lb &                            a[0] &                          A(1)                            \\
extent &                        a[1] &                          A(2)                            \\
oldtype &                       d[0] &                          D(1)                            \\
\hline
\end{tabular}
\end{center}

\item[\mpiconstmain{MPI\_COMBINER\_VALUE\_INDEX}]
\code{ni = 0}, \code{na = 0}, \code{nd = 2}, and

\smallskip%ALLOWLATEX%
\begin{center}
\begin{tabular}{l c c}
\hline
Constructor argument &          C &                    Fortran location                \\
\hline
value\_type &                   d[0] &                          D(1)                            \\
index\_type &                   d[1] &                          D(2)                            \\
\hline
\end{tabular}
\end{center}

\end{description}

\subsection{Examples}
\label{subsec:pt2pt-examples}

The following examples illustrate the use of derived datatypes.

\begin{example}
\label{pt2pt-exCC}
\Exindex{Fortran}{Send/receive of a 3D array}{MPI\_Type\_get\_extent,MPI\_Type\_vector,MPI\_Type\_create\_hvector,MPI\_Sendrecv,MPI\_Type\_commit}%
\exindex{Datatype!3D array}%
Send and receive a section of a 3D array.

%%HEADER
%%LANG: FORTRAN
%%FRAGMENT
%%ENDHEADER
\begin{lstlisting}[language={[MPI]Fortran}]
REAL a(100,100,100), e(9,9,9)
INTEGER oneslice, twoslice, threeslice, myrank, ierr
INTEGER(KIND=MPI_ADDRESS_KIND) lb, sizeofreal
INTEGER status(MPI_STATUS_SIZE)

! extract the section a(1:17:2, 3:11, 2:10)
! and store it in e(:,:,:).

CALL MPI_COMM_RANK(MPI_COMM_WORLD, myrank, ierr)

CALL MPI_TYPE_GET_EXTENT(MPI_REAL, lb, sizeofreal, ierr)

! create datatype for a 1D section
CALL MPI_TYPE_VECTOR(9, 1, 2, MPI_REAL, oneslice, ierr)

! create datatype for a 2D section
CALL MPI_TYPE_CREATE_HVECTOR(9, 1, 100*sizeofreal, oneslice, &
                             twoslice, ierr)

! create datatype for the entire section
CALL MPI_TYPE_CREATE_HVECTOR(9, 1, 100*100*sizeofreal, twoslice, &
                             threeslice, ierr)

CALL MPI_TYPE_COMMIT(threeslice, ierr)
CALL MPI_SENDRECV(a(1,3,2), 1, threeslice, myrank, 0, e, 9*9*9, &
                  MPI_REAL, myrank, 0, MPI_COMM_WORLD, status, ierr)
\end{lstlisting}
\end{example}


\begin{example}
\label{pt2pt-exDD}
\Exindex{Fortran}{Using indexed datatype}{MPI\_Type\_indexed,MPI\_Sendrecv,MPI\_Type\_commit}%
Copy the (strictly) lower triangular part of a matrix.
%%HEADER
%%LANG: FORTRAN
%%FRAGMENT
%%DECL: integer i
%%ENDHEADER
\begin{lstlisting}[language={[MPI]Fortran}]
REAL a(100,100), b(100,100)
INTEGER  disp(100), blocklen(100), ltype, myrank, ierr
INTEGER status(MPI_STATUS_SIZE)

! copy lower triangular part of array a
! onto lower triangular part of array b

CALL MPI_COMM_RANK(MPI_COMM_WORLD, myrank, ierr)

! compute start and size of each column
DO i=1,100
   disp(i) = 100*(i-1) + i
   blocklen(i) = 100-i
END DO

! create datatype for lower triangular part
CALL MPI_TYPE_INDEXED(100, blocklen, disp, MPI_REAL, ltype, ierr)

CALL MPI_TYPE_COMMIT(ltype, ierr)
CALL MPI_SENDRECV(a, 1, ltype, myrank, 0, b, 1, &
                  ltype, myrank, 0, MPI_COMM_WORLD, status, ierr)
\end{lstlisting}
\end{example}

\begin{example}
\label{pt2pt-exEE}
\Exindex{Fortran}{Nested vector datatypes}{MPI\_Type\_get\_extent,MPI\_Type\_vector,MPI\_Type\_create\_hvector,MPI\_Sendrecv,MPI\_Type\_commit}%
\exindex{Datatype!matrix transpose}%
Transpose a matrix.

%%HEADER
%%LANG: FORTRAN
%%FRAGMENT
%%ENDHEADER
\begin{lstlisting}[language={[MPI]Fortran}]
REAL a(100,100), b(100,100)
INTEGER row, xpose, myrank, ierr
INTEGER(KIND=MPI_ADDRESS_KIND) lb, sizeofreal
INTEGER status(MPI_STATUS_SIZE)

! transpose matrix a onto b

CALL MPI_COMM_RANK(MPI_COMM_WORLD, myrank, ierr)

CALL MPI_TYPE_GET_EXTENT(MPI_REAL, lb, sizeofreal, ierr)

! create datatype for one row
CALL MPI_TYPE_VECTOR(100, 1, 100, MPI_REAL, row, ierr)

! create datatype for matrix in row-major order
CALL MPI_TYPE_CREATE_HVECTOR(100, 1, sizeofreal, row, xpose, ierr)

CALL MPI_TYPE_COMMIT(xpose, ierr)

! send matrix in row-major order and receive in column major order
CALL MPI_SENDRECV(a, 1, xpose, myrank, 0, b, 100*100, &
                  MPI_REAL, myrank, 0, MPI_COMM_WORLD, status, ierr)
\end{lstlisting}
\end{example}

\begin{example}
\label{pt2pt-exFF}
\Exindex{Fortran}{Transpose with datatypes}{MPI\_Type\_get\_extent,MPI\_Type\_vector,MPI\_Type\_create\_struct,MPI\_Sendrecv,MPI\_Type\_commit}%
\exindex{Datatype!matrix transpose}%
Another approach to the transpose problem:
%%HEADER
%%LANG: FORTRAN
%%FRAGMENT
%%ENDHEADER
\begin{lstlisting}[language={[MPI]Fortran}]
REAL a(100,100), b(100,100)
INTEGER row, row1
INTEGER(KIND=MPI_ADDRESS_KIND) lb, sizeofreal
INTEGER myrank, ierr
INTEGER status(MPI_STATUS_SIZE)

CALL MPI_COMM_RANK(MPI_COMM_WORLD, myrank, ierr)

! transpose matrix a onto b

CALL MPI_TYPE_GET_EXTENT(MPI_REAL, lb, sizeofreal, ierr)

! create datatype for one row
CALL MPI_TYPE_VECTOR(100, 1, 100, MPI_REAL, row, ierr)

! create datatype for one row, with the extent of one real number
lb = 0
CALL MPI_TYPE_CREATE_RESIZED(row, lb, sizeofreal, row1, ierr)

CALL MPI_TYPE_COMMIT(row1, ierr)

! send 100 rows and receive in column major order
CALL MPI_SENDRECV(a, 100, row1, myrank, 0, b, 100*100, &
                  MPI_REAL, myrank, 0, MPI_COMM_WORLD, status, ierr)
\end{lstlisting}
\end{example}


\begin{example}
\label{pt2pt-exGG}
\Exindex{C}{Using datatypes with array of structures}{d:MPI\_Aint,MPI\_Get\_address,MPI\_Type\_create\_struct,MPI\_Type\_indexed,MPI\_Send,MPI\_Type\_get\_extent,MPI\_Type\_create\_hvector,MPI\_Type\_commit}%
\exindex{Datatype!array of structures}%
Use of \MPI/ datatypes to manipulate an array of structures.

%%HEADER
%%LANG: C
%%FRAGMENT
%%ENDHEADER
\begin{lstlisting}[language={[MPI]C}]
struct Partstruct
{
   int    type;   /* particle type */
   double d[6];   /* particle coordinates */
   char   b[7];   /* some additional information */
};

struct Partstruct    particle[1000];

int          i, dest, tag;
MPI_Comm     comm;


/* build datatype describing structure */

MPI_Datatype Particlestruct, Particletype;
MPI_Datatype type[3] = {MPI_INT, MPI_DOUBLE, MPI_CHAR};
int          blocklen[3] = {1, 6, 7};
MPI_Aint     disp[3];
MPI_Aint     base, lb, sizeofentry;


/* compute displacements of structure components */

MPI_Get_address(particle, disp);
MPI_Get_address(particle[0].d, disp+1);
MPI_Get_address(particle[0].b, disp+2);
base = disp[0];
for (i=0; i < 3; i++) disp[i] = MPI_Aint_diff(disp[i], base);

MPI_Type_create_struct(3, blocklen, disp, type, &Particlestruct);

/* Since the compiler may pad the structure, it is best to explicitly
   set the extent of the MPI datatype for a structure element using
   MPI_Type_create_resized */

/* compute extent of the structure */
MPI_Get_address(particle+1, &sizeofentry);
sizeofentry = MPI_Aint_diff(sizeofentry, base);

/* build datatype describing structure */
MPI_Type_create_resized(Particlestruct, 0, sizeofentry, &Particletype);


/* 4.1: send the entire array */

MPI_Type_commit(&Particletype);
MPI_Send(particle, 1000, Particletype, dest, tag, comm);


/* 4.2: send only the entries of type zero particles,
        preceded by the number of such entries */

MPI_Datatype Zparticles;   /* datatype describing all particles
                              with type zero (needs to be recomputed
                              if types change) */
MPI_Datatype Ztype;

int          zdisp[1000];
int          zblock[1000], j, k;
int          zzblock[2] = {1,1};
MPI_Aint     zzdisp[2];
MPI_Datatype zztype[2];

/* compute displacements of type zero particles */
j = 0;
for (i=0; i < 1000; i++)
   if (particle[i].type == 0)
      {
        zdisp[j] = i;
        zblock[j] = 1;
        j++;
      }

/* create datatype for type zero particles  */
MPI_Type_indexed(j, zblock, zdisp, Particletype, &Zparticles);

/* prepend particle count */
MPI_Get_address(&j, zzdisp);
MPI_Get_address(particle, zzdisp+1);
zztype[0] = MPI_INT;
zztype[1] = Zparticles;
MPI_Type_create_struct(2, zzblock, zzdisp, zztype, &Ztype);

MPI_Type_commit(&Ztype);
MPI_Send(MPI_BOTTOM, 1, Ztype, dest, tag, comm);


/* A probably more efficient way of defining Zparticles */

/* consecutive particles with index zero are handled as one block */
j=0;
for (i=0; i < 1000; i++)
   if (particle[i].type == 0)
      {
         for (k=i+1; (k < 1000)&&(particle[k].type == 0); k++);
         zdisp[j] = i;
         zblock[j] = k-i;
         j++;
         i = k;
      }
MPI_Type_indexed(j, zblock, zdisp, Particletype, &Zparticles);


/* 4.3: send the first two coordinates of all entries */

MPI_Datatype Allpairs;      /* datatype for all pairs of coordinates */

MPI_Type_get_extent(Particletype, &lb, &sizeofentry);

/* sizeofentry can also be computed by subtracting the address
   of particle[0] from the address of particle[1] */

MPI_Type_create_hvector(1000, 2, sizeofentry, MPI_DOUBLE, &Allpairs);
MPI_Type_commit(&Allpairs);
MPI_Send(particle[0].d, 1, Allpairs, dest, tag, comm);

/* an alternative solution to 4.3 */

MPI_Datatype Twodouble;

MPI_Type_contiguous(2, MPI_DOUBLE, &Twodouble);

MPI_Datatype Onepair;   /* datatype for one pair of coordinates, with
                          the extent of one particle entry */

MPI_Type_create_resized(Twodouble, 0, sizeofentry, &Onepair );
MPI_Type_commit(&Onepair);
MPI_Send(particle[0].d, 1000, Onepair, dest, tag, comm);
\end{lstlisting}
\end{example}

\begin{example}
\label{pt2pt-exHH}
\Exindex{C}{Using datatypes with array of structures with absolute addresses}{MPI\_Get\_address,MPI\_Type\_create\_struct,MPI\_Type\_indexed,MPI\_Send,MPI\_Type\_create\_hvector,MPI\_Type\_commit}%
\exindex{Datatype!absolute addresses}%
The same manipulations as in the previous example, but use absolute
addresses in datatypes.

%%HEADER
%%LANG: C
%%DECL: int dest, tag, i, j, k;
%%DECL: MPI_Comm comm;
%%SUBST:\/\* 5.2:}void testroutine2(void){/* 5.2
%%SUBST:MPI_Datatype Particletype;:MPI_Datatype Particletype;void testroutine(void){
%%TAIL: }
%%ENDHEADER
\begin{lstlisting}[language={[MPI]C}]
struct Partstruct
{
    int    type;
    double d[6];
    char   b[7];
};

struct Partstruct particle[1000];

/* build datatype describing first array entry */

MPI_Datatype Particletype;
MPI_Datatype type[3] = {MPI_INT, MPI_DOUBLE, MPI_CHAR};
int          block[3] = {1, 6, 7};
MPI_Aint     disp[3];

MPI_Get_address(particle, disp);
MPI_Get_address(particle[0].d, disp+1);
MPI_Get_address(particle[0].b, disp+2);
MPI_Type_create_struct(3, block, disp, type, &Particletype);

/* Particletype describes first array entry -- using absolute
   addresses */

/* 5.1: send the entire array */

MPI_Type_commit(&Particletype);
MPI_Send(MPI_BOTTOM, 1000, Particletype, dest, tag, comm);


/* 5.2: send the entries of type zero,
        preceded by the number of such entries */

MPI_Datatype Zparticles, Ztype;

int          zdisp[1000];
int          zblock[1000], i, j, k;
int          zzblock[2] = {1,1};
MPI_Datatype zztype[2];
MPI_Aint     zzdisp[2];

j=0;
for (i=0; i < 1000; i++)
    if (particle[i].type == 0)
        {
            for (k=i+1; (k < 1000)&&(particle[k].type == 0); k++);
            zdisp[j] = i;
            zblock[j] = k-i;
            j++;
            i = k;
        }
MPI_Type_indexed(j, zblock, zdisp, Particletype, &Zparticles);
/* Zparticles describe particles with type zero, using
   their absolute addresses*/

/* prepend particle count */
MPI_Get_address(&j, zzdisp);
zzdisp[1] = (MPI_Aint)0;
zztype[0] = MPI_INT;
zztype[1] = Zparticles;
MPI_Type_create_struct(2, zzblock, zzdisp, zztype, &Ztype);

MPI_Type_commit(&Ztype);
MPI_Send(MPI_BOTTOM, 1, Ztype, dest, tag, comm);
\end{lstlisting}
\end{example}

\begin{example}
\label{pt2pt-exII}
\Exindex{C}{Using datatypes with unions}{MPI\_Get\_address,MPI\_Type\_create\_resized,MPI\_Send,MPI\_Type\_commit}%
\exindex{Datatype!union}%
This example shows how datatypes can be used to handle unions.

%%HEADER
%%FRAGMENT
%%DECL: int dest, tag; MPI_Comm comm;
%%ENDHEADER
\begin{lstlisting}[language={[MPI]C}]
union {
   int     ival;
   float   fval;
      } u[1000];

int     i, utype;

/* All entries of u have identical type; variable
   utype keeps track of their current type */

MPI_Datatype   mpi_utype[2];
MPI_Aint       ubase, extent;

/* compute an MPI datatype for each possible union type;
   assume values are left-aligned in union storage. */

MPI_Get_address(u, &ubase);
MPI_Get_address(u+1, &extent);
extent = MPI_Aint_diff(extent, ubase);

MPI_Type_create_resized(MPI_INT, 0, extent, &mpi_utype[0]);

MPI_Type_create_resized(MPI_FLOAT, 0, extent, &mpi_utype[1]);

for(i=0; i<2; i++) MPI_Type_commit(&mpi_utype[i]);

/* actual communication */
MPI_Send(u, 1000, mpi_utype[utype], dest, tag, comm);
\end{lstlisting}
\end{example}

\begin{example}
\Exindex{C}{Decoding a datatype}{MPI\_Type\_get\_envelope,MPI\_Type\_get\_contents}%
This
example shows how a datatype can be decoded.  The routine
\gb\code{printdatatype} prints out the elements of the datatype.  Note the use
of \cfunc{MPI\_Type\_free} for datatypes that are not predefined.
%%WDG - MPI_Aint may not be a long, thus it is necessary to cast an
%% MPI_Aint value to (long) in order to use %ld to print the value.  This
%% is likely to be sufficient in most cases, though some systems may require
%% a different type for printing an MPI_Aint value.
%% The line that ends with ``array_of_dtypes);'' was indented slightly to
%% avoid an overfull (too long) line.  Do not ``fix'' the indentation; if
%% you must line up the argument list with the line above, you *must*
%% break that second line into two lines.
%%HEADER
%%LANG: C
%%SKIPELIPSIS
%%ENDHEADER
\begin{lstlisting}[language={[MPI]C}]
/*
  Example of decoding a datatype.

  Returns 0 if the datatype is predefined, 1 otherwise
 */
#include <stdio.h>
#include <stdlib.h>
#include "mpi.h"
int printdatatype(MPI_Datatype datatype)
{
    int *array_of_ints;
    MPI_Aint *array_of_adds;
    MPI_Datatype *array_of_dtypes;
    int num_ints, num_adds, num_dtypes, combiner;
    int i;

    MPI_Type_get_envelope(datatype,
                          &num_ints, &num_adds, &num_dtypes, &combiner);
    switch (combiner) {
    case MPI_COMBINER_NAMED:
        printf("Datatype is named:");
        /* To print the specific type, we can match against the
           predefined forms. We can NOT use a switch statement here
           We could also use MPI_TYPE_GET_NAME if we prefered to use
           names that the user may have changed.
         */
        if      (datatype == MPI_INT)    printf("MPI_INT\n");
        else if (datatype == MPI_DOUBLE) printf("MPI_DOUBLE\n");
        ... else test for other types ...
        return 0;
        break;
    case MPI_COMBINER_STRUCT:
        printf("Datatype is struct containing");
        array_of_ints   = (int *)malloc(num_ints * sizeof(int));
        array_of_adds   =
                   (MPI_Aint *) malloc(num_adds * sizeof(MPI_Aint));
        array_of_dtypes = (MPI_Datatype *)
            malloc(num_dtypes * sizeof(MPI_Datatype));
        MPI_Type_get_contents(datatype, num_ints, num_adds, num_dtypes,
                         array_of_ints, array_of_adds, array_of_dtypes);
        printf(" %d datatypes:\n", array_of_ints[0]);
        for (i=0; i<array_of_ints[0]; i++) {
            printf("blocklength %d, displacement %ld, type:\n",
                   array_of_ints[i+1], (long)array_of_adds[i]);
            if (printdatatype(array_of_dtypes[i])) {
                /* Note that we free the type ONLY if it
                   is not predefined */
                MPI_Type_free(&array_of_dtypes[i]);
            }
        }
        free(array_of_ints);
        free(array_of_adds);
        free(array_of_dtypes);
        break;
        ... other combiner values ...
    default:
        printf("Unrecognized combiner type\n");
    }
    return 1;
}
\end{lstlisting}
\end{example}



\section{Pack and Unpack}
\mpitermtitleindex{pack}
\mpitermtitleindex{unpack}
\label{sec:pt2pt-packing}

Some existing communication libraries provide pack/unpack procedures for sending
noncontiguous data. In these, the user explicitly packs
data into a contiguous buffer
before sending it, and unpacks it from a contiguous buffer after receiving it.
Derived datatypes, which are described in
Section~\ref{sec:pt2pt-datatype}, allow one, in most cases, to avoid
explicit packing and unpacking.  The user specifies the layout of
the data to be sent or received, and the communication library directly
accesses a noncontiguous buffer.  The pack/unpack routines are provided for
compatibility with previous libraries.   Also, they provide some functionality
that is not otherwise available in \MPI/.
For instance, a message can be received in several
parts, where the receive operation done on
a later part may depend on the content of a former part.
Another use is that outgoing messages
may be explicitly buffered in user supplied space, thus overriding the system
buffering policy.   Finally, the availability of pack and unpack operations
facilitates the development of additional communication libraries layered on top
of \MPI/.

\begin{funcdef}{MPI\_PACK}{(\mbox{inbuf}, \mbox{incount}, \mbox{datatype}, \mbox{outbuf}, \mbox{outsize}, \mbox{position}, \mbox{comm})}
\funcarg{\IN}{inbuf}{input buffer start (choice)}
\funcarg{\IN}{incount}{number of input data items (nonnegative integer)}
\funcarg{\IN}{datatype}{datatype of each input data item (handle)}
\funcarg{\OUT}{outbuf}{output buffer start (choice)}
\funcarg{\IN}{outsize}{output buffer size, in bytes (nonnegative integer)}
\funcarg{\INOUT}{position}{current position in buffer, in bytes (integer)}
\funcarg{\IN}{comm}{communicator for packed message (handle)}
\end{funcdef}
\mpibindmain{MPI\_Pack}{(\gb{}const~void~*inbuf, int~incount, MPI\_Datatype~datatype, void~*outbuf, int~outsize, int~*position, MPI\_Comm~comm)}
\MPIbindindex{MPI\_Pack\_c}
\mpibindmain{MPI\_Pack\_c}{(\gb{}const~void~*inbuf, MPI\_Count~incount, MPI\_Datatype~datatype, void~*outbuf, MPI\_Count~outsize, MPI\_Count~*position, MPI\_Comm~comm)}
\mpifnewbindmain{MPI\_Pack}{(inbuf, incount, datatype, outbuf, outsize, position, comm, ierror) \fargs TYPE(*), DIMENSION(..), INTENT(IN)\ ::\ \mbox{inbuf}\fnarg{}INTEGER, INTENT(IN)\ ::\ \mbox{incount}, \mbox{outsize}\fnarg{}TYPE(MPI\_Datatype), INTENT(IN)\ ::\ \mbox{datatype}\fnarg{}TYPE(*), DIMENSION(..)\ ::\ \mbox{outbuf}\fnarg{}INTEGER, INTENT(INOUT)\ ::\ \mbox{position}\fnarg{}TYPE(MPI\_Comm), INTENT(IN)\ ::\ \mbox{comm}\fnfinalarg{}INTEGER, OPTIONAL, INTENT(OUT)\ ::\ \mbox{ierror}}
\mpifnewbindmain{MPI\_Pack}{(inbuf, incount, datatype, outbuf, outsize, position, comm, ierror) !(\_c) \fargs TYPE(*), DIMENSION(..), INTENT(IN)\ ::\ \mbox{inbuf}\fnarg{}INTEGER(KIND=MPI\_COUNT\_KIND), INTENT(IN)\ ::\ \mbox{incount}, \mbox{outsize}\fnarg{}TYPE(MPI\_Datatype), INTENT(IN)\ ::\ \mbox{datatype}\fnarg{}TYPE(*), DIMENSION(..)\ ::\ \mbox{outbuf}\fnarg{}INTEGER(KIND=MPI\_COUNT\_KIND), INTENT(INOUT)\ ::\ \mbox{position}\fnarg{}TYPE(MPI\_Comm), INTENT(IN)\ ::\ \mbox{comm}\fnfinalarg{}INTEGER, OPTIONAL, INTENT(OUT)\ ::\ \mbox{ierror}}
\mpifbindmain{MPI\_PACK}{(INBUF, INCOUNT, DATATYPE, OUTBUF, OUTSIZE, POSITION, COMM, IERROR) \fargs <type> \mbox{INBUF(*)}, \mbox{OUTBUF(*)}\fnfinalarg{}INTEGER \mbox{INCOUNT}, \mbox{DATATYPE}, \mbox{OUTSIZE}, \mbox{POSITION}, \mbox{COMM}, \mbox{IERROR}}

Packs the message in the send buffer specified by \mpiarg{inbuf, incount,
datatype} into the buffer
space specified by \mpiarg{outbuf} and
\mpiarg{outsize}.  The
input buffer can
be any communication buffer allowed in \mpifunc{MPI\_SEND}.  The output buffer
is a contiguous storage area containing \mpiarg{outsize} bytes, starting at
the address \mpiarg{outbuf} (length is counted in \emph{bytes}, not elements,
as if it were a communication buffer for a message of type \type{MPI\_PACKED}).

The input value of \mpiarg{position} is the first
location in the output buffer to be used for packing.  \mpiarg{position} is
incremented by the size of the packed message, and the output value of
\mpiarg{position} is the first location in the output buffer
following the locations occupied by the packed message.
The \mpiarg{comm} argument
is the communicator that will be subsequently used for sending the packed
message.

\begin{funcdef}{MPI\_UNPACK}{(\mbox{inbuf}, \mbox{insize}, \mbox{position}, \mbox{outbuf}, \mbox{outcount}, \mbox{datatype}, \mbox{comm})}
\funcarg{\IN}{inbuf}{input buffer start (choice)}
\funcarg{\IN}{insize}{size of input buffer, in bytes (nonnegative integer)}
\funcarg{\INOUT}{position}{current position in bytes (integer)}
\funcarg{\OUT}{outbuf}{output buffer start (choice)}
\funcarg{\IN}{outcount}{number of items to be unpacked (integer)}
\funcarg{\IN}{datatype}{datatype of each output data item (handle)}
\funcarg{\IN}{comm}{communicator for packed message (handle)}
\end{funcdef}
\mpibindmain{MPI\_Unpack}{(\gb{}const~void~*inbuf, int~insize, int~*position, void~*outbuf, int~outcount, MPI\_Datatype~datatype, MPI\_Comm~comm)}
\MPIbindindex{MPI\_Unpack\_c}
\mpibindmain{MPI\_Unpack\_c}{(\gb{}const~void~*inbuf, MPI\_Count~insize, MPI\_Count~*position, void~*outbuf, MPI\_Count~outcount, MPI\_Datatype~datatype, MPI\_Comm~comm)}
\mpifnewbindmain{MPI\_Unpack}{(inbuf, insize, position, outbuf, outcount, datatype, comm, ierror) \fargs TYPE(*), DIMENSION(..), INTENT(IN)\ ::\ \mbox{inbuf}\fnarg{}INTEGER, INTENT(IN)\ ::\ \mbox{insize}, \mbox{outcount}\fnarg{}INTEGER, INTENT(INOUT)\ ::\ \mbox{position}\fnarg{}TYPE(*), DIMENSION(..)\ ::\ \mbox{outbuf}\fnarg{}TYPE(MPI\_Datatype), INTENT(IN)\ ::\ \mbox{datatype}\fnarg{}TYPE(MPI\_Comm), INTENT(IN)\ ::\ \mbox{comm}\fnfinalarg{}INTEGER, OPTIONAL, INTENT(OUT)\ ::\ \mbox{ierror}}
\mpifnewbindmain{MPI\_Unpack}{(inbuf, insize, position, outbuf, outcount, datatype, comm, ierror) !(\_c) \fargs TYPE(*), DIMENSION(..), INTENT(IN)\ ::\ \mbox{inbuf}\fnarg{}INTEGER(KIND=MPI\_COUNT\_KIND), INTENT(IN)\ ::\ \mbox{insize}, \mbox{outcount}\fnarg{}INTEGER(KIND=MPI\_COUNT\_KIND), INTENT(INOUT)\ ::\ \mbox{position}\fnarg{}TYPE(*), DIMENSION(..)\ ::\ \mbox{outbuf}\fnarg{}TYPE(MPI\_Datatype), INTENT(IN)\ ::\ \mbox{datatype}\fnarg{}TYPE(MPI\_Comm), INTENT(IN)\ ::\ \mbox{comm}\fnfinalarg{}INTEGER, OPTIONAL, INTENT(OUT)\ ::\ \mbox{ierror}}
\mpifbindmain{MPI\_UNPACK}{(INBUF, INSIZE, POSITION, OUTBUF, OUTCOUNT, DATATYPE, COMM, IERROR) \fargs <type> \mbox{INBUF(*)}, \mbox{OUTBUF(*)}\fnfinalarg{}INTEGER \mbox{INSIZE}, \mbox{POSITION}, \mbox{OUTCOUNT}, \mbox{DATATYPE}, \mbox{COMM}, \mbox{IERROR}}

Unpacks a message into the receive buffer specified by \mpiarg{outbuf,
outcount,
datatype} from the buffer
space specified by \mpiarg{inbuf} and \mpiarg{insize}.  The output buffer can
be any communication buffer allowed in \mpifunc{MPI\_RECV}.  The input
buffer is a contiguous storage area containing \mpiarg{insize} bytes,
starting at address \mpiarg{inbuf}.
The input value of \mpiarg{position} is the first location in
the input buffer occupied by the packed message.
\mpiarg{position} is incremented
by the size of the packed message, so that the output value of
\mpiarg{position} is the first location in the input buffer
after the locations occupied by the message that was unpacked.
\mpiarg{comm} is the communicator used to receive the packed message.

\begin{users}
Note the difference between \mpifunc{MPI\_RECV} and \mpifunc{MPI\_UNPACK}:  in
\mpifunc{MPI\_RECV}, the \mpiarg{count} argument specifies the maximum
number of
items that can be received.  The actual number of items received is determined
by the length of the incoming message.  In \mpifunc{MPI\_UNPACK}, the
\mpiarg{count} argument specifies the actual number
of items that are unpacked;
the ``size'' of the corresponding message is the increment in
\mpiarg{position}.
The reason for this change is that the ``incoming message size'' is not
predetermined since the user decides how much to unpack; nor is it easy to
determine
the ``message size'' from the number of items to be unpacked.  In fact, in a
heterogeneous system, this number may not be determined \emph{a priori}.
\end{users}

To understand the behavior of pack and unpack, it is convenient to think of the
data part of a message as being the sequence obtained by concatenating the
successive values sent in that message.   The pack operation stores this
sequence in the buffer space, as if sending the message to that buffer.  The
unpack operation retrieves this sequence from buffer space, as if receiving a
message from that buffer.  (It is helpful to think of internal Fortran files or
\code{sscanf} in C, for a similar function.)

Several messages can be successively packed into one \mpitermdef{packing unit}.
This
is effected by several successive \mpitermdef{related} calls to \mpifunc{MPI\_PACK},
where the first
call provides \mpiarg{position}\mpicode{ = 0}, and each successive call inputs the value
of \mpiarg{position} that was output by the previous call, and the same values
for \mpiarg{outbuf, outcount} and \mpiarg{comm}.   This packing unit
now contains
the equivalent information that would have been stored in a message by one send
call with a send buffer that is the ``concatenation'' of the individual send
buffers.

A  packing unit can
be sent using type \mpidtypemain{MPI\_PACKED}.  Any point-to-point
or collective communication operation can be used to move the sequence of bytes
that forms the packing unit from one process to another.  This packing unit
can now be
received using any receive operation, with any datatype:  the
type matching rules are relaxed for messages sent with type \type{MPI\_PACKED}.

A message sent with any type (including \type{MPI\_PACKED}) can be
received using the type \type{MPI\_PACKED}.  Such a message can then be
unpacked by calls to \mpifunc{MPI\_UNPACK}.

A packing unit (or a message created by a regular, ``typed'' send)
can be unpacked into several successive messages.  This is
effected by several successive related calls to \mpifunc{MPI\_UNPACK}, where
the first
call provides \mpiarg{position}\mpicode{ = 0}, and each successive call inputs the value
of \mpiarg{position} that was output by the previous call, and the same values
for \mpiarg{inbuf, insize} and \mpiarg{comm}.


The concatenation of two packing units is not necessarily a packing unit; nor is
a substring of a packing unit necessarily a packing unit.  Thus, one cannot
concatenate two packing units and then unpack the result as one packing
unit; nor can one unpack a substring of a packing unit as a separate
packing unit.  Each packing unit, that was created by a related
sequence of pack
calls, or by a regular send,
must be unpacked as a unit, by a sequence of related unpack calls.



\begin{rationale}
The restriction on ``atomic'' packing and unpacking of packing units allows
the implementation to add at the head of packing units additional
information, such as
a description of the sender architecture (to be used for type conversion, in a
heterogeneous environment)
\end{rationale}



The following call allows the user to find out how much space is
needed to pack a message and, thus, manage space allocation for
buffers.

\begin{funcdef}{MPI\_PACK\_SIZE}{(\mbox{incount}, \mbox{datatype}, \mbox{comm}, \mbox{size})}
\funcarg{\IN}{incount}{count argument to packing call (nonnegative integer)}
\funcarg{\IN}{datatype}{datatype argument to packing call (handle)}
\funcarg{\IN}{comm}{communicator argument to packing call (handle)}
\funcarg{\OUT}{size}{upper bound on size of packed message, in bytes (nonnegative integer)}
\end{funcdef}
\mpibindmain{MPI\_Pack\_size}{(\gb{}int~incount, MPI\_Datatype~datatype, MPI\_Comm~comm, int~*size)}
\MPIbindindex{MPI\_Pack\_size\_c}
\mpibindmain{MPI\_Pack\_size\_c}{(\gb{}MPI\_Count~incount, MPI\_Datatype~datatype, MPI\_Comm~comm, MPI\_Count~*size)}
\mpifnewbindmain{MPI\_Pack\_size}{(incount, datatype, comm, size, ierror) \fargs INTEGER, INTENT(IN)\ ::\ \mbox{incount}\fnarg{}TYPE(MPI\_Datatype), INTENT(IN)\ ::\ \mbox{datatype}\fnarg{}TYPE(MPI\_Comm), INTENT(IN)\ ::\ \mbox{comm}\fnarg{}INTEGER, INTENT(OUT)\ ::\ \mbox{size}\fnfinalarg{}INTEGER, OPTIONAL, INTENT(OUT)\ ::\ \mbox{ierror}}
\mpifnewbindmain{MPI\_Pack\_size}{(incount, datatype, comm, size, ierror) !(\_c) \fargs INTEGER(KIND=MPI\_COUNT\_KIND), INTENT(IN)\ ::\ \mbox{incount}\fnarg{}TYPE(MPI\_Datatype), INTENT(IN)\ ::\ \mbox{datatype}\fnarg{}TYPE(MPI\_Comm), INTENT(IN)\ ::\ \mbox{comm}\fnarg{}INTEGER(KIND=MPI\_COUNT\_KIND), INTENT(OUT)\ ::\ \mbox{size}\fnfinalarg{}INTEGER, OPTIONAL, INTENT(OUT)\ ::\ \mbox{ierror}}
\mpifbindmain{MPI\_PACK\_SIZE}{(INCOUNT, DATATYPE, COMM, SIZE, IERROR) \fargs \fnfinalarg{}INTEGER \mbox{INCOUNT}, \mbox{DATATYPE}, \mbox{COMM}, \mbox{SIZE}, \mbox{IERROR}}

A call to \mpifunc{MPI\_PACK\_SIZE}\mpicode{(incount, datatype, comm, size)}
returns in \mpiarg{size} an upper bound on the increment in \mpiarg{position}
that is effected by a call to \mpifunc{MPI\_PACK}\mpicode{(inbuf, incount, datatype,
outbuf, outcount, position, comm)}.
If the packed size of the datatype cannot be expressed by the \mpiarg{size}
parameter, then \mpifunc{MPI\_PACK\_SIZE} sets the value of \mpiarg{size} to
\mpiconst{MPI\_UNDEFINED}.

\begin{rationale}
The call returns an upper bound, rather than an exact bound, since the
exact amount of space needed to pack the message may depend on the
context (e.g., first message packed in a packing unit may take more
space).
\end{rationale}


\begin{example}
\label{pt2pt-exJJ}
\Exindex{C}{Using Pack}{MPI\_Pack}%
An example using \mpifunc{MPI\_PACK}.
%%HEADER
%%LANG: C
%%FRAGMENT
%%DECL: int myrank;
%%SKIPELIPSIS
%%ENDHEADER
\begin{lstlisting}[language={[MPI]C}]
int        position, i, j, a[2];
char       buff[1000];

MPI_Comm_rank(MPI_COMM_WORLD, &myrank);
if (myrank == 0)
{
    /* SENDER CODE */
    position = 0;
    MPI_Pack(&i, 1, MPI_INT, buff, 1000, &position, MPI_COMM_WORLD);
    MPI_Pack(&j, 1, MPI_INT, buff, 1000, &position, MPI_COMM_WORLD);
    MPI_Send(buff, position, MPI_PACKED, 1, 0, MPI_COMM_WORLD);
}
else  /* RECEIVER CODE */
    MPI_Recv(a, 2, MPI_INT, 0, 0, MPI_COMM_WORLD, MPI_STATUS_IGNORE);
\end{lstlisting}
\end{example}



\begin{example}
\label{pt2pt-exKK}
\Exindex{C}{Pack/Unpack with struct datatype}{MPI\_Pack,MPI\_Unpack,MPI\_Get\_address,MPI\_Type\_create\_struct,MPI\_Send,MPI\_Type\_commit}%
\exindex{Datatype!elaborate example}%
An elaborate example.
%%HEADER
%%LANG: C
%%FRAGMENT
%%DECL: int myrank;
%%SKIPELIPSIS
%%ENDHEADER
\begin{lstlisting}[language={[MPI]C}]
int   position, i = 200;
float a[200];
char  buff[1000]; /* larger than or equal to the size returned
                     from MPI_PACK_SIZE for 1,newtype */
MPI_Comm_rank(MPI_COMM_WORLD, &myrank);
if (myrank == 0)
{
    /* SENDER CODE */
    int len[2];
    MPI_Aint disp[2];
    MPI_Datatype type[2], newtype;

    /* build datatype for i followed by a[0]...a[i-1] */
    len[0] = 1;
    len[1] = i;
    MPI_Get_address(&i, disp);
    MPI_Get_address(a, disp+1);
    type[0] = MPI_INT;
    type[1] = MPI_FLOAT;
    MPI_Type_create_struct(2, len, disp, type, &newtype);
    MPI_Type_commit(&newtype);

    /* Pack i followed by a[0]...a[i-1]*/
    position = 0;
    MPI_Pack(MPI_BOTTOM, 1, newtype, buff, 1000, &position,
             MPI_COMM_WORLD);

    /* Send */
    MPI_Send(buff, position, MPI_PACKED, 1, 0,
             MPI_COMM_WORLD);

/* *****
   One can replace the last three lines with
   MPI_Send(MPI_BOTTOM, 1, newtype, 1, 0, MPI_COMM_WORLD);
   ***** */
}
else if (myrank == 1)
{
    /* RECEIVER CODE */
    MPI_Status status;

    /* Receive */
    MPI_Recv(buff, 1000, MPI_PACKED, 0, 0, MPI_COMM_WORLD, &status);

    /* Unpack i */
    position = 0;
    MPI_Unpack(buff, 1000, &position, &i, 1, MPI_INT, MPI_COMM_WORLD);

    /* Unpack a[0]...a[i-1] */
    MPI_Unpack(buff, 1000, &position, a, i, MPI_FLOAT, MPI_COMM_WORLD);
}
\end{lstlisting}
\end{example}

\begin{example}
\label{pt2pt-exLL}
\Exindex{C}{Pack and Pack\_size}{MPI\_Pack,MPI\_Unpack,MPI\_Pack\_size,MPI\_Gather,MPI\_Gatherv}%
\exindex{Datatype!elaborate example}%
Each process sends a count, followed by count characters to the root;
the root
concatenates
all characters into one string.
%%HEADER
%%LANG: C
%%FRAGMENT
%%DECL: MPI_Comm comm; int myrank, root, i;
%%SKIPELIPSIS
%%ENDHEADER
\begin{lstlisting}[language={[MPI]C}]
int  count, gsize, counts[64], totalcount, k1, k2, k,
     displs[64], position, concat_pos;
char chr[100], *lbuf, *rbuf, *cbuf;

MPI_Comm_size(comm, &gsize);
MPI_Comm_rank(comm, &myrank);

      /* allocate local pack buffer */
MPI_Pack_size(1, MPI_INT, comm, &k1);
MPI_Pack_size(count, MPI_CHAR, comm, &k2);
k = k1+k2;
lbuf = (char *)malloc(k);

      /* pack count, followed by count characters */
position = 0;
MPI_Pack(&count, 1, MPI_INT, lbuf, k, &position, comm);
MPI_Pack(chr, count, MPI_CHAR, lbuf, k, &position, comm);

if (myrank != root) {
    /* gather at root sizes of all packed messages */
    MPI_Gather(&position, 1, MPI_INT, NULL, 0,
               MPI_DATATYPE_NULL, root, comm);

    /* gather at root packed messages */
    MPI_Gatherv(lbuf, position, MPI_PACKED, NULL,
                NULL, NULL, MPI_DATATYPE_NULL, root, comm);

} else {   /* root code */
    /* gather sizes of all packed messages */
    MPI_Gather(&position, 1, MPI_INT, counts, 1,
               MPI_INT, root, comm);

    /* gather all packed messages */
    displs[0] = 0;
    for (i=1; i < gsize; i++)
        displs[i] = displs[i-1] + counts[i-1];
    totalcount = displs[gsize-1] + counts[gsize-1];
    rbuf = (char *)malloc(totalcount);
    cbuf = (char *)malloc(totalcount);
    MPI_Gatherv(lbuf, position, MPI_PACKED, rbuf,
                counts, displs, MPI_PACKED, root, comm);

    /* unpack all messages and concatenate strings */
    concat_pos = 0;
    for (i=0; i < gsize; i++) {
        position = 0;
        MPI_Unpack(rbuf+displs[i], totalcount-displs[i],
                   &position, &count, 1, MPI_INT, comm);
        MPI_Unpack(rbuf+displs[i], totalcount-displs[i],
                   &position, cbuf+concat_pos, count, MPI_CHAR, comm);
        concat_pos += count;
    }
    cbuf[concat_pos] = '\0';
}
\end{lstlisting}
\end{example}

\section{Canonical \texorpdfstring{\mpifuncskip{MPI\_PACK}}{MPI\_PACK} and \texorpdfstring{\mpifuncskip{MPI\_UNPACK}}{MPI\_UNPACK}}
\mpitermtitleindex{canonical pack and unpack}
\mpitermtitleindex{pack!canonical}
\mpitermtitleindex{unpack!canonical}
\label{canonical_pack}

These procedures read/write data to/from the buffer in the
\mpistring{external32} data
format specified in Section~\ref{subsec:ext32}, and calculate the size needed
for packing.  Their first arguments specify the data format, for future
extensibility, but
currently
the only valid value of the \mpiarg{datarep}
argument is \mpistring{external32}.
\begin{users}
  These procedures could be used, for example, to send typed data in a portable
  format from one \MPI/ implementation to another.
\end{users}

The buffer will contain exactly the packed data, without headers.
\type{MPI\_BYTE} should be used to send and receive data that is packed
using \mpifunc{MPI\_PACK\_EXTERNAL}.

\begin{rationale}
\mpifunc{MPI\_PACK\_EXTERNAL} specifies that there is no header on the message
and further
specifies the exact format of the data.  Since \mpifunc{MPI\_PACK} may (and is
allowed
to) use a header, the datatype \type{MPI\_PACKED} cannot be used for data
packed with
\mpifunc{MPI\_PACK\_EXTERNAL}.
\end{rationale}

\cdeclindex{MPI\_Aint}%
\begin{funcdef}{MPI\_PACK\_EXTERNAL}{(\mbox{datarep}, \mbox{inbuf}, \mbox{incount}, \mbox{datatype}, \mbox{outbuf}, \mbox{outsize}, \mbox{position})}
\funcarg{\IN}{datarep}{data representation (string)}
\funcarg{\IN}{inbuf}{input buffer start (choice)}
\funcarg{\IN}{incount}{number of input data items (integer)}
\funcarg{\IN}{datatype}{datatype of each input data item (handle)}
\funcarg{\OUT}{outbuf}{output buffer start (choice)}
\funcarg{\IN}{outsize}{output buffer size, in bytes (integer)}
\funcarg{\INOUT}{position}{current position in buffer, in bytes (integer)}
\end{funcdef}
\mpibindmain{MPI\_Pack\_external}{(\gb{}const~char~datarep[], const~void~*inbuf, int~incount, MPI\_Datatype~datatype, void~*outbuf, MPI\_Aint~outsize, MPI\_Aint~*position)}
\MPIbindindex{MPI\_Pack\_external\_c}
\mpibindmain{MPI\_Pack\_external\_c}{(\gb{}const~char~datarep[], const~void~*inbuf, MPI\_Count~incount, MPI\_Datatype~datatype, void~*outbuf, MPI\_Count~outsize, MPI\_Count~*position)}
\mpifnewbindmain{MPI\_Pack\_external}{(datarep, inbuf, incount, datatype, outbuf, outsize, position, ierror) \fargs CHARACTER(LEN=*), INTENT(IN)\ ::\ \mbox{datarep}\fnarg{}TYPE(*), DIMENSION(..), INTENT(IN)\ ::\ \mbox{inbuf}\fnarg{}INTEGER, INTENT(IN)\ ::\ \mbox{incount}\fnarg{}TYPE(MPI\_Datatype), INTENT(IN)\ ::\ \mbox{datatype}\fnarg{}TYPE(*), DIMENSION(..)\ ::\ \mbox{outbuf}\fnarg{}INTEGER(KIND=MPI\_ADDRESS\_KIND), INTENT(IN)\ ::\ \mbox{outsize}\fnarg{}INTEGER(KIND=MPI\_ADDRESS\_KIND), INTENT(INOUT)\ ::\ \mbox{position}\fnfinalarg{}INTEGER, OPTIONAL, INTENT(OUT)\ ::\ \mbox{ierror}}
\mpifnewbindmain{MPI\_Pack\_external}{(datarep, inbuf, incount, datatype, outbuf, outsize, position, ierror) !(\_c) \fargs CHARACTER(LEN=*), INTENT(IN)\ ::\ \mbox{datarep}\fnarg{}TYPE(*), DIMENSION(..), INTENT(IN)\ ::\ \mbox{inbuf}\fnarg{}INTEGER(KIND=MPI\_COUNT\_KIND), INTENT(IN)\ ::\ \mbox{incount}, \mbox{outsize}\fnarg{}TYPE(MPI\_Datatype), INTENT(IN)\ ::\ \mbox{datatype}\fnarg{}TYPE(*), DIMENSION(..)\ ::\ \mbox{outbuf}\fnarg{}INTEGER(KIND=MPI\_COUNT\_KIND), INTENT(INOUT)\ ::\ \mbox{position}\fnfinalarg{}INTEGER, OPTIONAL, INTENT(OUT)\ ::\ \mbox{ierror}}
\mpifbindmain{MPI\_PACK\_EXTERNAL}{(DATAREP, INBUF, INCOUNT, DATATYPE, OUTBUF, OUTSIZE, POSITION, IERROR) \fargs CHARACTER*(*) \mbox{DATAREP}\fnarg{}<type> \mbox{INBUF(*)}, \mbox{OUTBUF(*)}\fnarg{}INTEGER \mbox{INCOUNT}, \mbox{DATATYPE}, \mbox{IERROR}\fnfinalarg{}INTEGER(KIND=MPI\_ADDRESS\_KIND) \mbox{OUTSIZE}, \mbox{POSITION}}

\cdeclindex{MPI\_Aint}%
\begin{funcdef}{MPI\_UNPACK\_EXTERNAL}{(\mbox{datarep}, \mbox{inbuf}, \mbox{insize}, \mbox{position}, \mbox{outbuf}, \mbox{outcount}, \mbox{datatype})}
\funcarg{\IN}{datarep}{data representation (string)}
\funcarg{\IN}{inbuf}{input buffer start (choice)}
\funcarg{\IN}{insize}{input buffer size, in bytes (integer)}
\funcarg{\INOUT}{position}{current position in buffer, in bytes (integer)}
\funcarg{\OUT}{outbuf}{output buffer start (choice)}
\funcarg{\IN}{outcount}{number of output data items (integer)}
\funcarg{\IN}{datatype}{datatype of output data item (handle)}
\end{funcdef}
\mpibindmain{MPI\_Unpack\_external}{(\gb{}const~char~datarep[], const~void~*inbuf, MPI\_Aint~insize, MPI\_Aint~*position, void~*outbuf, int~outcount, MPI\_Datatype~datatype)}
\MPIbindindex{MPI\_Unpack\_external\_c}
\mpibindmain{MPI\_Unpack\_external\_c}{(\gb{}const~char~datarep[], const~void~*inbuf, MPI\_Count~insize, MPI\_Count~*position, void~*outbuf, MPI\_Count~outcount, MPI\_Datatype~datatype)}
\mpifnewbindmain{MPI\_Unpack\_external}{(datarep, inbuf, insize, position, outbuf, outcount, datatype, ierror) \fargs CHARACTER(LEN=*), INTENT(IN)\ ::\ \mbox{datarep}\fnarg{}TYPE(*), DIMENSION(..), INTENT(IN)\ ::\ \mbox{inbuf}\fnarg{}INTEGER(KIND=MPI\_ADDRESS\_KIND), INTENT(IN)\ ::\ \mbox{insize}\fnarg{}INTEGER(KIND=MPI\_ADDRESS\_KIND), INTENT(INOUT)\ ::\ \mbox{position}\fnarg{}TYPE(*), DIMENSION(..)\ ::\ \mbox{outbuf}\fnarg{}INTEGER, INTENT(IN)\ ::\ \mbox{outcount}\fnarg{}TYPE(MPI\_Datatype), INTENT(IN)\ ::\ \mbox{datatype}\fnfinalarg{}INTEGER, OPTIONAL, INTENT(OUT)\ ::\ \mbox{ierror}}
\mpifnewbindmain{MPI\_Unpack\_external}{(datarep, inbuf, insize, position, outbuf, outcount, datatype, ierror) !(\_c) \fargs CHARACTER(LEN=*), INTENT(IN)\ ::\ \mbox{datarep}\fnarg{}TYPE(*), DIMENSION(..), INTENT(IN)\ ::\ \mbox{inbuf}\fnarg{}INTEGER(KIND=MPI\_COUNT\_KIND), INTENT(IN)\ ::\ \mbox{insize}, \mbox{outcount}\fnarg{}INTEGER(KIND=MPI\_COUNT\_KIND), INTENT(INOUT)\ ::\ \mbox{position}\fnarg{}TYPE(*), DIMENSION(..)\ ::\ \mbox{outbuf}\fnarg{}TYPE(MPI\_Datatype), INTENT(IN)\ ::\ \mbox{datatype}\fnfinalarg{}INTEGER, OPTIONAL, INTENT(OUT)\ ::\ \mbox{ierror}}
\mpifbindmain{MPI\_UNPACK\_EXTERNAL}{(DATAREP, INBUF, INSIZE, POSITION, OUTBUF, OUTCOUNT, DATATYPE, IERROR) \fargs CHARACTER*(*) \mbox{DATAREP}\fnarg{}<type> \mbox{INBUF(*)}, \mbox{OUTBUF(*)}\fnarg{}INTEGER(KIND=MPI\_ADDRESS\_KIND) \mbox{INSIZE}, \mbox{POSITION}\fnfinalarg{}INTEGER \mbox{OUTCOUNT}, \mbox{DATATYPE}, \mbox{IERROR}}

\cdeclindex{MPI\_Aint}%
\begin{funcdef}{MPI\_PACK\_EXTERNAL\_SIZE}{(\mbox{datarep}, \mbox{incount}, \mbox{datatype}, \mbox{size})}
\funcarg{\IN}{datarep}{data representation (string)}
\funcarg{\IN}{incount}{number of input data items (integer)}
\funcarg{\IN}{datatype}{datatype of each input data item (handle)}
\funcarg{\OUT}{size}{output buffer size, in bytes (integer)}
\end{funcdef}
\mpibindmain{MPI\_Pack\_external\_size}{(\gb{}const~char~datarep[], int~incount, MPI\_Datatype~datatype, MPI\_Aint~*size)}
\MPIbindindex{MPI\_Pack\_external\_size\_c}
\mpibindmain{MPI\_Pack\_external\_size\_c}{(\gb{}const~char~datarep[], MPI\_Count~incount, MPI\_Datatype~datatype, MPI\_Count~*size)}
\mpifnewbindmain{MPI\_Pack\_external\_size}{(datarep, incount, datatype, size, ierror) \fargs CHARACTER(LEN=*), INTENT(IN)\ ::\ \mbox{datarep}\fnarg{}INTEGER, INTENT(IN)\ ::\ \mbox{incount}\fnarg{}TYPE(MPI\_Datatype), INTENT(IN)\ ::\ \mbox{datatype}\fnarg{}INTEGER(KIND=MPI\_ADDRESS\_KIND), INTENT(OUT)\ ::\ \mbox{size}\fnfinalarg{}INTEGER, OPTIONAL, INTENT(OUT)\ ::\ \mbox{ierror}}
\mpifnewbindmain{MPI\_Pack\_external\_size}{(datarep, incount, datatype, size, ierror) !(\_c) \fargs CHARACTER(LEN=*), INTENT(IN)\ ::\ \mbox{datarep}\fnarg{}INTEGER(KIND=MPI\_COUNT\_KIND), INTENT(IN)\ ::\ \mbox{incount}\fnarg{}TYPE(MPI\_Datatype), INTENT(IN)\ ::\ \mbox{datatype}\fnarg{}INTEGER(KIND=MPI\_COUNT\_KIND), INTENT(OUT)\ ::\ \mbox{size}\fnfinalarg{}INTEGER, OPTIONAL, INTENT(OUT)\ ::\ \mbox{ierror}}
\mpifbindmain{MPI\_PACK\_EXTERNAL\_SIZE}{(DATAREP, INCOUNT, DATATYPE, SIZE, IERROR) \fargs CHARACTER*(*) \mbox{DATAREP}\fnarg{}INTEGER \mbox{INCOUNT}, \mbox{DATATYPE}, \mbox{IERROR}\fnfinalarg{}INTEGER(KIND=MPI\_ADDRESS\_KIND) \mbox{SIZE}}

